% Nombres réels — Mathématiques 3ème — Cours signé OnBuch+
% Programme officiel MINESEC. Compiler avec : tectonic N03-nombres-reels.tex
\newcommand{\DOCMATIERE}{Mathématiques}
\newcommand{\DOCNIVEAU}{3ème}
\newcommand{\DOCMODULE}{Mathématiques – classe de 3ème}
\newcommand{\DOCLECON}{N03}
\newcommand{\DOCTITRE}{Nombres réels}
% OnBuch+ — préambule « cours signés » ENRICHI (compilé avec tectonic / XeLaTeX)
% Système visuel « Pop » d'OnBuch+ : fond crème (jamais blanc), bordures encre
% épaisses, ombres « dures » décalées, titres Archivo Black en capitales.
% Version MATHÉMATIQUES : graphes (pgfplots/tikz), boîtes pédagogiques
% (définition, propriété, méthode, attention, le savais-tu, exercice résolu,
% exercice, TP, bilan), page de garde et signature OnBuch+ ancrée en bas.
%
% Macros attendues dans main.tex AVANT \input{preamble} :
%   \DOCMATIERE  \DOCNIVEAU  \DOCMODULE  \DOCLECON (numéro)  \DOCTITRE
\documentclass[11pt]{article}

\usepackage{fontspec}
\usepackage[french]{babel}
\usepackage[a4paper,margin=2cm,top=3.4cm,bottom=2.8cm,headheight=1.4cm,footskip=1.3cm]{geometry}
\usepackage{amsmath,amssymb,amsfonts}
\usepackage{mathtools} % \xmapsto, \coloneqq... souvent utilisés par les modèles
\usepackage{cancel} % simplifications barrées
% \abs{z} (module, valeur absolue) et \norm{u} (norme), utilisés par les modèles
\providecommand{\abs}{}\let\abs\relax\DeclarePairedDelimiter{\abs}{\lvert}{\rvert}
\providecommand{\norm}{}\let\norm\relax\DeclarePairedDelimiter{\norm}{\lVert}{\rVert}
\usepackage[table]{xcolor} % [table] : fournit aussi \rowcolors (alternance de lignes)
\usepackage{pagecolor}
\usepackage{tikz}
\usetikzlibrary{arrows.meta,positioning,calc,shapes.geometric,shapes.misc,shapes.arrows,decorations.pathmorphing,decorations.pathreplacing,decorations.markings,patterns,shadows,mindmap,trees,fit,backgrounds,matrix,intersections,angles,quotes}
\usepackage{pgfplots}
\usepackage{pgf-pie} % diagrammes circulaires \pie (statistiques)
\pgfplotsset{compat=1.18}
\usepgfplotslibrary{fillbetween,groupplots} % aires entre courbes (calcul intégral)
\usepackage{enumitem}
\usepackage{fancyhdr}
% [most] charge toutes les bibliothèques tcolorbox (listings, minted, raster,
% documentation...) et gonfle inutilement la mémoire TeX sur un document long
% (~300 boîtes) : on ne charge que ce qui est réellement utilisé.
\usepackage[skins,breakable]{tcolorbox}
\usepackage{graphicx}
\usepackage{colortbl}
\usepackage{booktabs}
\usepackage{tabularx}
\usepackage{diagbox} % case d'en-tête diagonale des tableaux à double entrée
% Colonnes extensibles centrée / gauche / droite (C, L, R), souvent utilisées par les modèles
\newcolumntype{C}{>{\centering\arraybackslash}X}
\newcolumntype{L}{>{\raggedright\arraybackslash}X}
\newcolumntype{R}{>{\raggedleft\arraybackslash}X}
\usepackage{array}
\usepackage{multicol}
\usepackage{siunitx}
% parse-numbers=false : accepte aussi \SI{2,5 \times 10^{-3}}{...}
\sisetup{locale=FR,output-decimal-marker={,},group-minimum-digits=4,parse-numbers=false}
\DeclareSIUnit\ohm{\ensuremath{\symup{\Omega}}} % U+2126 absent de la police
\usepackage{qrcode}
% \degree : commande fréquemment utilisée par les modèles pour un angle ou une
% température en dehors de \SI{}{} (non fournie par les paquets chargés).
\providecommand{\degree}{^{\circ}}
% \qed (fin de démonstration), utilisé par les modèles sans amsthm
\providecommand{\qed}{\hfill$\square$}
% \barre : double barre des valeurs interdites dans un tableau de variations (tkz-tab)
\providecommand{\barre}{\ensuremath{\|}}
% environnement proof (amsthm non chargé) utilisé par les modèles
\ifdefined\proof\else\newenvironment{proof}[1][Démonstration]{\par\noindent\textit{#1.}\ }{\qed\par}\fi
\usepackage{lastpage}
\usepackage{hyperref}
\hypersetup{hidelinks}

% ============================================================
% Palette « Pop » OnBuch+ — mêmes hex que la classe P (plus_ui.dart)
% ============================================================
\definecolor{poporange}{HTML}{F59321}
\definecolor{popdark}{HTML}{E07A0C}
\definecolor{popcream}{HTML}{FFF6E8}
\definecolor{popink}{HTML}{1C1714}
\definecolor{poppurple}{HTML}{7A5AE0}
\definecolor{popgreen}{HTML}{2E9E6A}
\definecolor{poppink}{HTML}{E0457B}
\definecolor{popblue}{HTML}{2F7DE1}
\definecolor{popgold}{HTML}{C98A12}
\definecolor{popmuted}{HTML}{8A7F76}
\definecolor{popline}{HTML}{EADFD2}
\definecolor{popgray}{HTML}{8A7F76}
\colorlet{popgrey}{popgray}
% Alias tolérants (le modèle nomme parfois les couleurs différemment)
\colorlet{popwhite}{white}
\colorlet{popblack}{popink}
\colorlet{popred}{poppink}
\colorlet{popyellow}{popgold}
% Teintes claires pour fonds de tableaux / zones
\colorlet{popblueL}{popblue!10!white}
\colorlet{poporangeL}{poporange!14!white}
\colorlet{popgreenL}{popgreen!12!white}
\colorlet{poppurpleL}{poppurple!12!white}
\colorlet{poppinkL}{poppink!12!white}
\colorlet{popgoldL}{popgold!14!white}

\pagecolor{popcream}

% ============================================================
% Typographie
% ============================================================
\defaultfontfeatures{Ligatures=TeX}
\setmainfont{PlusJakartaSans-Medium.ttf}[Path=../../../fonts/,BoldFont=PlusJakartaSans-Bold.ttf]
% Maths en sans-serif (Fira Math) avec chiffres et lettres droites de Plus
% Jakarta, pour que les formules se fondent dans le texte.
\usepackage[math-style=ISO,bold-style=ISO]{unicode-math}
\setmathfont{FiraMath-Regular.otf}[Path=../../../fonts/]
\setmathfont{PlusJakartaSans-Medium.ttf}[Path=../../../fonts/,range={up/num,up/{Latin,latin},\mathup}]
\usepackage{icomma} % virgule décimale sans espace en mode math
% Caractères mathématiques Unicode absents des polices de texte (Plus Jakarta,
% Archivo Black) : rendus par leur équivalent mathématique, en texte comme en
% formule — sinon ils s'affichent en carré vide (ex. « Arithmétique dans ℤ »).
\usepackage{newunicodechar}
\newunicodechar{ℝ}{\ensuremath{\mathbb{R}}}
\newunicodechar{ℤ}{\ensuremath{\mathbb{Z}}}
\newunicodechar{ℂ}{\ensuremath{\mathbb{C}}}
\newunicodechar{ℕ}{\ensuremath{\mathbb{N}}}
\newunicodechar{ℚ}{\ensuremath{\mathbb{Q}}}
\newunicodechar{𝔻}{\ensuremath{\mathbb{D}}}
\newunicodechar{α}{\ensuremath{\alpha}}
\newunicodechar{β}{\ensuremath{\beta}}
\newunicodechar{γ}{\ensuremath{\gamma}}
\newunicodechar{δ}{\ensuremath{\delta}}
\newunicodechar{ε}{\ensuremath{\varepsilon}}
\newunicodechar{θ}{\ensuremath{\theta}}
\newunicodechar{λ}{\ensuremath{\lambda}}
\newunicodechar{μ}{\ensuremath{\mu}}
\newunicodechar{σ}{\ensuremath{\sigma}}
\newunicodechar{τ}{\ensuremath{\tau}}
\newunicodechar{φ}{\ensuremath{\varphi}}
\newunicodechar{ψ}{\ensuremath{\psi}}
\newunicodechar{ω}{\ensuremath{\omega}}
\newunicodechar{Σ}{\ensuremath{\Sigma}}
% majuscules produites par \MakeUppercase dans les titres (α-aminés -> Α)
\newunicodechar{Α}{\ensuremath{\alpha}}
\newunicodechar{Β}{\ensuremath{\beta}}
\newunicodechar{Γ}{\ensuremath{\Gamma}}
\newunicodechar{Δ}{\ensuremath{\Delta}}
\newunicodechar{Ε}{\ensuremath{\varepsilon}}
\newunicodechar{Θ}{\ensuremath{\Theta}}
\newunicodechar{Λ}{\ensuremath{\Lambda}}
\newunicodechar{Μ}{\ensuremath{\mu}}
\newunicodechar{Τ}{\ensuremath{\tau}}
\newunicodechar{Φ}{\ensuremath{\Phi}}
\newunicodechar{Ψ}{\ensuremath{\Psi}}
\newunicodechar{Ω}{\ensuremath{\Omega}}
\newunicodechar{↦}{\ensuremath{\mapsto}}
\newunicodechar{∈}{\ensuremath{\in}}
\newunicodechar{∉}{\ensuremath{\notin}}
\newunicodechar{∧}{\ensuremath{\wedge}}
\newunicodechar{⇔}{\ensuremath{\Leftrightarrow}}
\newunicodechar{⟺}{\ensuremath{\Longleftrightarrow}}
\newunicodechar{⇒}{\ensuremath{\Rightarrow}}
\newunicodechar{⟹}{\ensuremath{\Longrightarrow}}
\newunicodechar{∘}{\ensuremath{\circ}}
\newunicodechar{∪}{\ensuremath{\cup}}
\newunicodechar{∩}{\ensuremath{\cap}}
\newunicodechar{≡}{\ensuremath{\equiv}}
\newunicodechar{⊂}{\ensuremath{\subset}}
\newunicodechar{⊥}{\ensuremath{\perp}}
\newunicodechar{∃}{\ensuremath{\exists}}
\newunicodechar{∀}{\ensuremath{\forall}}
\newunicodechar{∛}{\ensuremath{\sqrt[3]{\,}}}
\newunicodechar{∜}{\ensuremath{\sqrt[4]{\,}}}
\newunicodechar{⃗}{\ensuremath{{}^{\to}}}
\newunicodechar{ⁿ}{\ifmmode^{n}\else\textsuperscript{\ensuremath{n}}\fi}
\newunicodechar{ˣ}{\ifmmode^{x}\else\textsuperscript{\ensuremath{x}}\fi}
\newunicodechar{ᵇ}{\ifmmode^{b}\else\textsuperscript{\ensuremath{b}}\fi}
\newunicodechar{ʳ}{\ifmmode^{r}\else\textsuperscript{\ensuremath{r}}\fi}
\newunicodechar{ᵘ}{\ifmmode^{u}\else\textsuperscript{\ensuremath{u}}\fi}
\newunicodechar{ʸ}{\ifmmode^{y}\else\textsuperscript{\ensuremath{y}}\fi}
\newunicodechar{ᵃ}{\ifmmode^{a}\else\textsuperscript{\ensuremath{a}}\fi}
\newunicodechar{ᵉ}{\ifmmode^{e}\else\textsuperscript{\ensuremath{e}}\fi}
\newunicodechar{ᵏ}{\ifmmode^{k}\else\textsuperscript{\ensuremath{k}}\fi}
\newunicodechar{ᵖ}{\ifmmode^{p}\else\textsuperscript{\ensuremath{p}}\fi}
\newunicodechar{ᴺ}{\ifmmode^{N}\else\textsuperscript{\ensuremath{N}}\fi}
\newunicodechar{ᶜ}{\ifmmode^{c}\else\textsuperscript{\ensuremath{c}}\fi}
\newunicodechar{ʰ}{\ifmmode^{h}\else\textsuperscript{\ensuremath{h}}\fi}
\newunicodechar{ᵠ}{\ifmmode^{q}\else\textsuperscript{\ensuremath{q}}\fi}
\newunicodechar{ₙ}{\ifmmode_{n}\else\textsubscript{\ensuremath{n}}\fi}
\newunicodechar{ₐ}{\ifmmode_{a}\else\textsubscript{\ensuremath{a}}\fi}
\newunicodechar{ᵢ}{\ifmmode_{i}\else\textsubscript{\ensuremath{i}}\fi}
\newunicodechar{ᵣ}{\ifmmode_{r}\else\textsubscript{\ensuremath{r}}\fi}
\newunicodechar{ₖ}{\ifmmode_{k}\else\textsubscript{\ensuremath{k}}\fi}
\newunicodechar{ᵦ}{\ifmmode_{\beta}\else\textsubscript{\ensuremath{\beta}}\fi}
\newunicodechar{ₓ}{\ifmmode_{x}\else\textsubscript{\ensuremath{x}}\fi}
\newfontfamily\popdisplay{ArchivoBlack-Regular.ttf}[Path=../../../fonts/]
\newfontfamily\poplogo{SpaceGrotesk-Bold.ttf}[Path=../../../fonts/]
\newfontfamily\popsemi{PlusJakartaSans-SemiBold.ttf}[Path=../../../fonts/]
\newfontfamily\popxbold{PlusJakartaSans-ExtraBold.ttf}[Path=../../../fonts/]
\newfontfamily\popxboldls{PlusJakartaSans-ExtraBold.ttf}[Path=../../../fonts/,LetterSpace=3.0]

\setlist[enumerate]{leftmargin=1.7em,itemsep=5pt,topsep=4pt}
\setlist[itemize]{leftmargin=1.7em,itemsep=4pt,topsep=4pt,label={\color{poporange}\textbullet}}
\setlength{\parindent}{0pt}
\setlength{\parskip}{5pt}
\renewcommand{\arraystretch}{1.25}

% pgfplots : style Pop par défaut
\pgfplotsset{
  every axis/.append style={axis line style={popink,line width=1pt},tick style={popink},
    grid style={popline,line width=0.5pt},label style={font=\small},tick label style={font=\footnotesize}},
  popcurve/.style={poporange,line width=1.8pt,no markers,smooth},
}
% Même style disponible dans un \draw TikZ simple (hors pgfplots)
\tikzset{popcurve/.style={poporange,line width=1.8pt}}

% ============================================================
% Pill / squiggle
% ============================================================
\newcommand{\poppill}[3]{%
  \tikz[baseline=-0.55ex]\node[draw=popink,line width=1.2pt,rounded corners=2.6pt,%
    fill=#2,inner xsep=6.5pt,inner ysep=3pt]{%
    \popxboldls\fontsize{7}{7}\selectfont\color{#3}\MakeUppercase{#1}};%
}
\newcommand{\popsquiggle}[1]{%
  \tikz[baseline=-0.4ex]{
    \draw[#1,line width=2.6pt,line cap=round]
      plot [smooth] coordinates {(0,0.09) (0.34,0.01) (0.68,0.21) (1.02,0.01) (1.36,0.10)};
  }%
}
% Mot-clé surligné (marqueur orange sous le texte)
\newcommand{\cle}[1]{{\popxbold\color{popdark}#1}}

% ============================================================
% En-tête / pied
% ============================================================
\pagestyle{fancy}
\fancyhf{}
\renewcommand{\headrulewidth}{0pt}
\renewcommand{\footrulewidth}{0pt}
\lhead{\raisebox{-0.15ex}{\poplogo\fontsize{13}{13}\selectfont\color{popink}OnBuch\color{poporange}+}}
\rhead{\poppill{\DOCMATIERE\ \textbullet\ \DOCNIVEAU\ \textbullet\ Leçon \DOCLECON}{white}{popink}}
\lfoot{\popxbold\fontsize{7.6}{7.6}\selectfont\color{popmuted}\MakeUppercase{Cours signé OnBuch+}\ \textbullet\ {\popsemi\DOCTITRE}}
\rfoot{\tikz[baseline=-0.6ex]\node[rounded corners=8.5pt,fill=popink,minimum height=17pt,minimum width=17pt,inner xsep=5pt,inner sep=0]{\popxbold\fontsize{8}{8}\selectfont\color{popcream}\thepage\,/\,\pageref*{LastPage}};}

% ============================================================
% Titres
% ============================================================
\newcommand{\chaptitre}[2]{%
  \begin{tcolorbox}[enhanced,colback=poporange,colframe=popink,boxrule=2.6pt,arc=11pt,
      drop shadow=popink,left=15pt,right=15pt,top=12pt,bottom=11pt,before skip=3pt,after skip=18pt]
    \poppill{#2}{popink}{popcream}\par\vspace{9pt}
    {\raggedright\hyphenpenalty=10000\exhyphenpenalty=10000\popdisplay\fontsize{22}{24}\selectfont\color{popcream}\MakeUppercase{#1}\par}
  \end{tcolorbox}
}
% Section numérotée : pastille orange avec le numéro + titre Archivo
\newcounter{popsec}
\newcommand{\coursec}[1]{%
  \stepcounter{popsec}\setcounter{popsub}{0}%
  \par\vspace{16pt}\noindent
  \begin{minipage}{\linewidth}
  \tikz[baseline=-0.7ex]\node[draw=popink,line width=1.6pt,fill=poporange,rounded corners=4pt,minimum size=22pt,inner sep=0]{\popdisplay\fontsize{12}{12}\selectfont\color{popcream}\thepopsec};%
  \hspace{8pt}{\popdisplay\fontsize{15}{17}\selectfont\color{popink}\MakeUppercase{#1}}\par
  \vspace{4pt}\noindent{\color{poporange}\rule{36pt}{3.6pt}}
  \end{minipage}\par\nopagebreak\vspace{8pt}
}
\newcounter{popsub}
\newcommand{\courssub}[1]{%
  \stepcounter{popsub}%
  \par\vspace{9pt}\noindent{\popxbold\fontsize{11.5}{13}\selectfont\color{popink}{\color{poporange}\thepopsec.\thepopsub}\hspace{6pt}#1}\par\nopagebreak\vspace{4pt}
}

% ============================================================
% Boîtes pédagogiques — même recette : fond blanc, bordure épaisse colorée,
% ombre dure encre, badge plein. Titre optionnel [#1].
% ============================================================
\tcbset{popbase/.style={breakable,enhanced,colback=white,boxrule=2.3pt,arc=9pt,
  shadow={3pt}{-3pt}{0pt}{popink},left=14pt,right=14pt,top=11pt,bottom=11pt,before skip=11pt,after skip=15pt,
  fontupper=\popsemi\fontsize{10.6}{14.5}\selectfont\color{popink},
  fontlower=\popsemi\fontsize{10.6}{14.5}\selectfont\color{popink}}}
% Coche dessinée (le glyphe ✓ n'existe pas dans les polices du document)
\newcommand{\popcheck}[1]{\tikz[baseline=-0.5ex]\draw[#1,line width=1.6pt,line cap=round,line join=round](-0.13,0.0)--(-0.04,-0.09)--(0.13,0.1);}
\newcommand{\popboxhead}[3]{\poppill{#1}{#2}{white}\ifx\relax#3\relax\else\hspace{6pt}{\popxbold\fontsize{10.6}{12}\selectfont\color{popink}#3}\fi\par\vspace{7pt}}

\newtcolorbox{aretenir}[1][]{popbase,colframe=popgreen,before upper={\popboxhead{À retenir}{popgreen}{#1}}}
\newtcolorbox{exemplebox}[1][]{popbase,colframe=poporange,before upper={\popboxhead{Exemple}{poporange}{#1}}}
\newtcolorbox{definition}[1][]{popbase,colframe=popblue,before upper={\popboxhead{Définition}{popblue}{#1}}}
\newtcolorbox{propriete}[1][]{popbase,colframe=poppurple,before upper={\popboxhead{Propriété}{poppurple}{#1}}}
\newtcolorbox{methode}[1][]{popbase,colframe=popgold,before upper={\popboxhead{Méthode}{popgold}{#1}}}
\newtcolorbox{attention}[1][]{popbase,colframe=poppink,before upper={\popboxhead{Attention}{poppink}{#1}}}
\newtcolorbox{experience}[1][]{popbase,colframe=popblue,colback=popblueL,before upper={\popboxhead{Expérience}{popblue}{#1}}}
% « Le savais-tu ? » : carte violette pleine (contexte, culture, Cameroun)
\newtcolorbox{savaistu}[1][]{popbase,colback=poppurple,colframe=popink,
  fontupper=\popsemi\fontsize{10.4}{14.2}\selectfont\color{white},
  before upper={\poppill{Le savais-tu\,?}{white}{poppurple}\ifx\relax#1\relax\else\hspace{6pt}{\popxbold\color{white}#1}\fi\par\vspace{7pt}}}
% Exercice résolu : énoncé en haut, \tcblower puis la solution sur fond crème
\newcounter{popexo}\newcounter{popexr}
\newtcolorbox{exoresolu}[1][]{popbase,colframe=poporange,colbacklower=popcream,
  segmentation style={popink,line width=1.2pt,dashed},
  before upper={\stepcounter{popexr}\popboxhead{Exercice résolu \thepopexr}{poporange}{#1}},
  before lower={\poppill{Solution}{popink}{popcream}\par\vspace{6pt}}}
% Exercice (non résolu en place) — la correction est dans la partie Corrigés
\newtcolorbox{exercice}[1][]{popbase,colframe=popink,
  before upper={\stepcounter{popexo}\popboxhead{Exercice \thepopexo}{popink}{#1}}}
% Corrigé
\newtcolorbox{corrige}[1][]{popbase,colframe=popgreen,colback=popgreenL,
  before upper={\popboxhead{Corrigé}{popgreen}{#1}}}
% Objectifs / prérequis / bilan
\newtcolorbox{objectifs}{popbase,colframe=popink,colback=poporangeL,
  before upper={\popboxhead{Objectifs de la leçon}{popink}{}}}
\newtcolorbox{prerequis}{popbase,colframe=popink,colback=popblueL,
  before upper={\popboxhead{Prérequis}{popblue}{}}}
\newtcolorbox{bilan}{popbase,colframe=popink,colback=white,boxrule=2.8pt,
  before upper={\popboxhead{Fiche bilan}{poporange}{L'essentiel à connaître pour l'examen}}}
% Figure encadrée avec légende
\newtcolorbox{popfigure}[1][]{enhanced,breakable=false,colback=white,colframe=popline,boxrule=1.4pt,arc=8pt,
  left=10pt,right=10pt,top=10pt,bottom=8pt,before skip=10pt,after skip=12pt,halign=center,
  lower separated=false,halign lower=center,
  fontlower=\popsemi\fontsize{9}{11.5}\selectfont\color{popmuted},
  #1}
\newcommand{\legende}[1]{\par\vspace{6pt}{\popsemi\fontsize{9}{11.5}\selectfont\color{popmuted}#1}}

% ============================================================
% Signature OnBuch+
% ============================================================
\newcommand{\poplogoinline}[1]{{\poplogo\fontsize{#1}{#1}\selectfont\color{popink}OnBuch\color{poporange}+}}
\newcommand{\signatureonbuch}{%
  \begin{tcolorbox}[enhanced,colback=popink,colframe=popink,boxrule=2.2pt,arc=10pt,
      drop shadow=poporange,left=14pt,right=14pt,top=10pt,bottom=10pt,before skip=0pt,after skip=0pt]
    \tikz[baseline=-0.7ex]\node[circle,fill=popgreen,draw=popcream,line width=1.3pt,minimum size=22pt,inner sep=0]{\popcheck{white}};%
    \hspace{9pt}\begin{minipage}[c]{0.62\linewidth}
      {\popxbold\fontsize{12}{13}\selectfont\color{popcream}Signé\ \textbullet\ {\poplogo OnBuch\color{poporange}+}}\par\vspace{2pt}
      {\raggedright\popsemi\fontsize{8}{10.5}\selectfont\color{popline}\MakeUppercase{Équipe pédagogique OnBuch+}\\\MakeUppercase{Conforme au programme officiel MINESEC}\par}
    \end{minipage}\hfill
    {\poplogo\fontsize{22}{22}\selectfont\color{popcream}OnBuch\color{poporange}+}
  \end{tcolorbox}%
}

% ============================================================
% Page de garde
% #1 titre · #2 matière · #3 niveau · #4 module · #5 n° leçon · #6 durée ·
% #7 description / objectifs (texte ou itemize)
% ============================================================
\newcommand{\pagegarde}[7]{%
  \thispagestyle{empty}
  \newgeometry{a4paper,margin=1.7cm,top=1.6cm,bottom=1.4cm}%
  \begin{tikzpicture}[remember picture,overlay]
    \fill[poporange!18] ([xshift=-3.2cm,yshift=-3.6cm]current page.north east) circle (4.6cm);
    \fill[poppurple!14] ([xshift=2.4cm,yshift=7.6cm]current page.south west) circle (3.8cm);
  \end{tikzpicture}%
  {\poplogo\fontsize{30}{30}\selectfont\color{popink}OnBuch\color{poporange}+}\hfill
  \raisebox{0.6ex}{\poppill{Cours signé \textbullet\ Premium}{popink}{popcream}}\par
  \vspace{1pt}\popsquiggle{poppurple}\par
  \vspace{8pt}
  \begin{tcolorbox}[enhanced,colback=poporange,colframe=popink,boxrule=2.8pt,arc=13pt,
      drop shadow=popink,left=18pt,right=18pt,top=12pt,bottom=13pt,after skip=10pt]
    \poppill{#2 \textbullet\ #3}{popink}{popcream}\hspace{5pt}\poppill{#4}{white}{popink}\par\vspace{8pt}
    {\popdisplay\fontsize{14}{14}\selectfont\color{popink}LEÇON #5}\par\vspace{4pt}
    {\raggedright\hyphenpenalty=10000\exhyphenpenalty=10000\popdisplay\fontsize{25}{27}\selectfont\color{popcream}\MakeUppercase{#1}\par}
    \vspace{8pt}
    {\popxbold\fontsize{9.5}{9.5}\selectfont\color{popink}\MakeUppercase{Durée conseillée : #6}}
  \end{tcolorbox}
  \begin{tcolorbox}[enhanced,colback=white,colframe=popink,boxrule=2.2pt,arc=10pt,
      drop shadow=popink,left=16pt,right=16pt,top=10pt,bottom=10pt,after skip=8pt]
    \poppill{Dans cette leçon}{poppurple}{white}\par\vspace{5pt}
    \popsemi\fontsize{9.3}{12.6}\selectfont\color{popink}#7
  \end{tcolorbox}
  \noindent\begin{minipage}[t]{0.487\linewidth}
    \begin{tcolorbox}[enhanced,colback=popgreen,colframe=popink,boxrule=2.2pt,arc=10pt,
        drop shadow=popink,left=11pt,right=11pt,top=10pt,bottom=10pt,height=3.1cm,valign=center]
      \tcbox[colback=white,colframe=popink,boxrule=1.3pt,arc=5pt,boxsep=0pt,nobeforeafter,
        left=3pt,right=3pt,top=3pt,bottom=3pt,valign=center]{\qrcode[height=1.9cm]{https://play.google.com/store/apps/details?id=com.luvvix.onbuch&hl=fr}}%
      \hspace{8pt}\begin{minipage}[c]{0.5\linewidth}
        {\popdisplay\fontsize{11}{12.5}\selectfont\color{white}\MakeUppercase{Télécharge OnBuch}}\par\vspace{4pt}
        {\raggedright\popsemi\fontsize{8.4}{11}\selectfont\color{white}Gratuit \textbullet\ Google Play\\Tuteur IA, annales, concours, résultats\par}
      \end{minipage}
    \end{tcolorbox}
  \end{minipage}\hfill
  \begin{minipage}[t]{0.487\linewidth}
    \begin{tcolorbox}[enhanced,colback=poppurple,colframe=popink,boxrule=2.2pt,arc=10pt,
        drop shadow=popink,left=11pt,right=11pt,top=10pt,bottom=10pt,height=3.1cm,valign=center]
      \tikz[baseline=-0.7ex]\node[circle,fill=white,draw=popink,line width=1.3pt,minimum size=1.6cm,inner sep=0]{\popdisplay\fontsize{24}{24}\selectfont\color{poppurple}+};%
      \hspace{8pt}\begin{minipage}[c]{0.6\linewidth}
        {\popdisplay\fontsize{11}{12.5}\selectfont\color{white}\MakeUppercase{Passe à OnBuch+}}\par\vspace{4pt}
        {\raggedright\popsemi\fontsize{8.4}{11}\selectfont\color{white}Tous les cours signés, Tuteur IA illimité, fiches PDF — onglet OnBuch+ de l'app\par}
      \end{minipage}
    \end{tcolorbox}
  \end{minipage}
  \vspace{8pt}
  \popcontenu
  \vfill
  \signatureonbuch
  \restoregeometry
  \newpage
}

% Rangée « Ce cours contient » (page de garde)
\newcommand{\poptile}[3]{% #1 couleur #2 gros texte #3 légende
  \begin{tcolorbox}[enhanced,colback=#1,colframe=popink,boxrule=2pt,arc=9pt,shadow={2.5pt}{-2.5pt}{0pt}{popink},
      left=6pt,right=6pt,top=9pt,bottom=9pt,halign=center,valign=center,height=1.75cm,width=\linewidth,nobeforeafter]
    {\popdisplay\fontsize{12.5}{13}\selectfont\color{white}\MakeUppercase{#2}}\par\vspace{3pt}
    {\centering\popsemi\fontsize{7.8}{9.5}\selectfont\color{white}#3\par}
  \end{tcolorbox}}
\newcommand{\popcontenu}{%
  \par\noindent{\popxboldls\fontsize{7.5}{7.5}\selectfont\color{popmuted}CE COURS SIGNÉ CONTIENT}\par\vspace{7pt}
  \noindent\begin{minipage}[t]{0.235\linewidth}\poptile{popblue}{Cours}{complet, illustré, pas à pas}\end{minipage}\hfill
  \begin{minipage}[t]{0.235\linewidth}\poptile{poporange}{Méthodes}{et exercices résolus}\end{minipage}\hfill
  \begin{minipage}[t]{0.235\linewidth}\poptile{poppink}{Exercices}{type examen + corrigés}\end{minipage}\hfill
  \begin{minipage}[t]{0.235\linewidth}\poptile{popgreen}{Fiche bilan}{l'essentiel à retenir}\end{minipage}\par}

% Sommaire
\newtcolorbox{sommaire}{enhanced,colback=white,colframe=popink,boxrule=2.2pt,arc=10pt,
  drop shadow=popink,left=18pt,right=18pt,top=13pt,bottom=13pt,before skip=6pt,after skip=20pt,
  before upper={\poppill{Sommaire}{poppurple}{white}\par\vspace{9pt}\popsemi\fontsize{10.6}{14.5}\selectfont\color{popink}}}

% Signature de fin de cours — poussée tout en bas de la dernière page
\newcommand{\findecours}{%
  \par\vfill\nopagebreak
  \begin{center}\popsquiggle{poporange}\end{center}\vspace{4pt}
  \signatureonbuch
}
\AtBeginDocument{\def\llbracket{\mathopen{[\mkern-3.5mu[}}\def\rrbracket{\mathclose{]\mkern-3.5mu]}}} % intervalles d'entiers (glyphe ⟦ absent de la police)
% Glyphes absents de Fira Math : redéfinis à partir de glyphes disponibles
\AtBeginDocument{%
  \def\setminus{\mathbin{\backslash}}\let\smallsetminus\setminus
  \def\vdots{\mathord{\vbox{\baselineskip3.2pt\lineskiplimit0pt\kern4pt\hbox{.}\hbox{.}\hbox{.}}}}%
  \def\ddots{\mathinner{\mkern1mu\raise6.5pt\hbox{.}\mkern2mu\raise3.6pt\hbox{.}\mkern2mu\raise0.7pt\hbox{.}\mkern1mu}}%
  \def\triangle{\mathord{\scalebox{0.9}{$\symup{\Delta}$}}}%
  \def\nabla{\mathord{\raisebox{\depth}{\scalebox{1}[-1]{$\symup{\Delta}$}}}}%
  \def\checkmark{\popcheck{popdark}}%
  \def\star{\mathbin{\ast}}\def\bigstar{\mathord{\ast}}%
  \def\diamond{\mathbin{\scalebox{0.6}{$\Diamond$}}}%
  \def\bigcup{\mathop{\mathchoice{\raisebox{-0.3em}{\scalebox{1.7}{$\cup$}}}{\scalebox{1.25}{$\cup$}}{\cup}{\cup}}\displaylimits}%
  \def\bigcap{\mathop{\mathchoice{\raisebox{-0.3em}{\scalebox{1.7}{$\cap$}}}{\scalebox{1.25}{$\cap$}}{\cap}{\cap}}\displaylimits}%
  \def\lesseqqgtr{\mathrel{\lessgtr}}\def\bigoplus{\mathop{\oplus}\displaylimits}%
}
\providecommand{\ssi}{si et seulement si} % abréviation fréquente
\AtBeginDocument{\providecommand{\wideparen}{\overparen}} % arc de cercle

%% FIGFIX-BEGIN v5 — correctifs globaux des figures (docs/onbuch-generation/tools/figfix)
\usepackage{adjustbox}\usepackage{etoolbox}\tcbuselibrary{raster}
% toute figure TikZ/pgfplots plus large que la ligne est réduite à la largeur disponible
\BeforeBeginEnvironment{tikzpicture}{\begin{adjustbox}{max width=\linewidth}}
\AfterEndEnvironment{tikzpicture}{\end{adjustbox}}
% légendes pgfplots lisibles par-dessus les courbes
\pgfplotsset{every axis/.append style={legend style={fill=white,fill opacity=0.92,text opacity=1,draw=black!15,inner sep=3pt}}}
% étiquettes d'arêtes (syntaxe edge["..."]) avec fond blanc
\tikzset{every edge quotes/.append style={fill=white,inner sep=1.2pt}}
%% FIGFIX-END
\begin{document}

\pagegarde{Nombres réels}{Mathématiques}{3ème}{Mathématiques – classe de 3ème}{N03}{10 séances (fiche de progression harmonisée)}{Cette leçon introduit les nombres réels à travers les racines carrées et leurs opérations (somme, produit, quotient, puissances), puis développe la comparaison, l'encadrement et la notion d'intervalle. Elle consolide l'écriture et la manipulation des nombres irrationnels, compétence centrale du programme de 3ème et du Probatoire.
\vspace{4pt}
\begin{multicols}{2}\small
\begin{itemize}
\item À la fin de cette leçon, je sais définir la racine carrée d'un nombre positif et calculer les racines carrées usuelles
\item À la fin de cette leçon, je sais justifier qu'un nombre comme √2 est irrationnel et situer les ensembles ℕ, ℤ, 𝔻, ℚ, ℝ les uns dans les autres
\item À la fin de cette leçon, je sais simplifier une somme ou un produit de réels comportant des radicaux
\item À la fin de cette leçon, je sais rendre rationnel le dénominateur d'un quotient comportant un radical
\item À la fin de cette leçon, je sais utiliser les règles de calcul sur les puissances entières d'un nombre réel
\item À la fin de cette leçon, je sais comparer deux nombres réels comportant un radical en comparant leurs carrés
\end{itemize}\end{multicols}}

\begin{sommaire}
\begin{multicols}{2}
\begin{enumerate}
\item Racine carrée d'un réel positif
\item L'ensemble des nombres réels
\item Somme et produit de réels avec radicaux
\item Quotient de réels avec radicaux
\item Puissances entières d'un réel
\item Comparaison, encadrement et intervalles
\item Activité d'intégration
\item Méthodes et exercices résolus
\item Exercices
\item Corrigés des exercices
\item Fiche bilan
\end{enumerate}
\end{multicols}
\end{sommaire}

\begin{objectifs}
\begin{itemize}
\item À la fin de cette leçon, je sais définir la racine carrée d'un nombre positif et calculer les racines carrées usuelles
\item À la fin de cette leçon, je sais justifier qu'un nombre comme √2 est irrationnel et situer les ensembles ℕ, ℤ, 𝔻, ℚ, ℝ les uns dans les autres
\item À la fin de cette leçon, je sais simplifier une somme ou un produit de réels comportant des radicaux
\item À la fin de cette leçon, je sais rendre rationnel le dénominateur d'un quotient comportant un radical
\item À la fin de cette leçon, je sais utiliser les règles de calcul sur les puissances entières d'un nombre réel
\item À la fin de cette leçon, je sais comparer deux nombres réels comportant un radical en comparant leurs carrés
\item À la fin de cette leçon, je sais encadrer un nombre réel à une précision donnée et donner une valeur approchée
\item À la fin de cette leçon, je sais représenter et utiliser les intervalles de ℝ (bornés ou non)
\item À la fin de cette leçon, je sais rédiger et justifier une démarche complète dans un problème concret de la vie courante
\end{itemize}
\end{objectifs}

\begin{prerequis}
\begin{itemize}
\item Carrés parfaits et calcul littéral de 4ème (identités remarquables)
\item Fractions : simplification, sommes et quotients
\item Théorème de Pythagore (lien géométrique avec √2)
\item Puissances d'un nombre (exposants positifs vus en 4ème)
\item Droite graduée et comparaison de nombres décimaux
\end{itemize}
\end{prerequis}

\newpage

\coursec{Racine carrée d'un réel positif}

\courssub{Situation d'introduction : le jardin de Moussa}

À Ngaoundéré, Moussa vient d'acheter un jardin \textbf{carré} de \SI{50}{m^2}. Il veut le clôturer avec du grillage vendu au mètre. Pour connaître la longueur à acheter, il lui faut d'abord le côté $c$ de son carré, puis multiplier par $4$.

L'aire d'un carré de côté $c$ vaut $c^2$. Moussa doit donc résoudre :
\[ c^2 = 50. \]
Essayons : $7^2 = 49$ (trop petit), $8^2 = 64$ (trop grand). Le nombre cherché est donc entre $7$ et $8$... mais aucun nombre décimal, aucune fraction ne convient exactement ! Ce nombre existe pourtant : c'est la \cle{racine carrée} de $50$, notée $\sqrt{50}$.

\textbf{Problématique.} Qu'est-ce que $\sqrt{50}$ ? Comment calculer avec ce nombre, le comparer à $7$, l'encadrer au centimètre près ? Bienvenue dans l'ensemble des \cle{nombres réels}, où tous les nombres rencontrés au collège se réunissent enfin.

\courssub{Définition de la racine carrée}

\begin{experience}[Une activité géométrique pour découvrir]
Construis un carré $ABCD$ de côté \SI{1}{cm}. Son aire vaut $1$ cm$^2$. Construis maintenant le carré $BMNC$ dont un côté est la \emph{diagonale} $[BC]$ du premier carré.
\begin{enumerate}
\item D'après le théorème de Pythagore dans le triangle $ABC$ rectangle en $A$ :
\[ BC^2 = AB^2 + AC^2 = 1^2 + 1^2 = 2. \]
\item L'aire du carré $BMNC$ vaut donc $BC^2 = 2$ cm$^2$. Son côté $BC$ est le nombre positif dont le carré vaut $2$ : on le note $\sqrt{2}$.
\item Avec un compas, reporte la longueur $BC$ sur une droite graduée : tu places ainsi le point d'abscisse $\sqrt{2}$, entre $1$ et $2$.
\end{enumerate}
\end{experience}

\begin{popfigure}
\begin{tikzpicture}[scale=2.2]
% Carré unité
\draw[fill=popblueL,draw=popblue,thick] (0,0) -- (1,0) -- (1,1) -- (0,1) -- cycle;
\node[popdark] at (0.5,0.5) {aire $1$};
% Diagonale
\draw[poporange,very thick] (0,1) -- (1,0);
\node[poporange,above left] at (0.45,0.55) {$\sqrt{2}$};
% Carré sur la diagonale
\draw[fill=poporangeL,draw=poporange,thick,opacity=0.85] (0,1) -- (1,0) -- (2,1) -- (1,2) -- cycle;
\node[popdark] at (1,1) {aire $2$};
% Sommets
\fill[popdark] (0,0) circle (0.02) node[below left] {$A$};
\fill[popdark] (1,0) circle (0.02) node[below] {$B$};
\fill[popdark] (1,1) circle (0.02) node[right] {$C$};
\fill[popdark] (0,1) circle (0.02) node[left] {$D$};
% Droite graduée et report au compas
\draw[->,thick] (-0.3,-0.9) -- (2.6,-0.9);
\foreach \x in {0,1,2} {\draw (\x,-0.95) -- (\x,-0.85); \node[below] at (\x,-0.95) {\x};}
\draw[popgreen,dashed,thick] (0,-0.9) arc[start angle=180,end angle=0,radius=0.707];
\fill[popgreen] (1.414,-0.9) circle (0.025);
\node[popgreen,below] at (1.414,-1.0) {$\sqrt{2}$};
\end{tikzpicture}
\legende{Le carré d'aire $2$ construit sur la diagonale du carré unité, et le report de $\sqrt{2}$ sur la droite graduée au compas.}
\end{popfigure}

\begin{definition}[Racine carrée d'un nombre positif]
Soit $a$ un nombre \textbf{positif ou nul}. La \cle{racine carrée} de $a$, notée $\sqrt{a}$, est l'\textbf{unique nombre positif ou nul dont le carré est égal à} $a$.
\[ \text{Si } a \geq 0, \quad \sqrt{a} \geq 0 \quad \text{et} \quad \left(\sqrt{a}\right)^2 = a. \]
Le symbole $\sqrt{\phantom{a}}$ s'appelle le \cle{radical}.
\end{definition}

\begin{exemplebox}[Premiers exemples]
\begin{itemize}
\item $\sqrt{9} = 3$ car $3 \geq 0$ et $3^2 = 9$ ;
\item $\sqrt{0} = 0$ car $0^2 = 0$ ;
\item $\sqrt{2} \approx 1,414$ : c'est le côté du carré d'aire $2$ de notre activité ;
\item $\sqrt{50} \approx 7,07$ : c'est le côté du jardin de Moussa !
\end{itemize}
\end{exemplebox}

\begin{attention}[La racine carrée est toujours positive]
$\sqrt{9} = 3$ et \textbf{uniquement} $3$, jamais $-3$ ! Certes $(-3)^2 = 9$, mais par définition la racine carrée désigne le nombre \emph{positif} dont le carré vaut $9$. On ne peut donc pas écrire $\sqrt{9} = \pm 3$.
\end{attention}

\courssub{Existence : pourquoi $\sqrt{a}$ n'existe pas si $a < 0$}

\begin{propriete}[Condition d'existence]
Dans l'ensemble $\mathbb{R}$ des nombres réels, $\sqrt{a}$ \textbf{n'existe que si} $a \geq 0$.
\end{propriete}

\begin{experience}[Pourquoi $\sqrt{-4}$ n'existe pas]
Raisonnons par l'absurde : supposons qu'il existe un nombre réel $x$ tel que $x^2 = -4$.
\begin{itemize}
\item Si $x > 0$, alors $x^2 > 0$ (produit de deux positifs) ;
\item si $x < 0$, alors $x^2 > 0$ (produit de deux négatifs) ;
\item si $x = 0$, alors $x^2 = 0$.
\end{itemize}
Dans tous les cas, $x^2 \geq 0$ : un carré est \textbf{toujours positif ou nul}. Il est donc impossible que $x^2 = -4$. Conclusion : $\sqrt{-4}$ n'existe pas dans $\mathbb{R}$.
\end{experience}

\begin{attention}[Erreurs fréquentes]
\begin{itemize}
\item Écrire $\sqrt{-4} = -2$ est \textbf{faux} : $(-2)^2 = 4$ et non $-4$.
\item Écrire $\sqrt{-4} = 2$ est \textbf{faux} aussi : $2^2 = 4$.
\item Avant de calculer une racine carrée, vérifie toujours que le nombre sous le radical est \textbf{positif ou nul}.
\end{itemize}
\end{attention}

\courssub{Propriétés fondamentales}

\begin{propriete}[Carré d'une racine carrée]
Pour tout nombre $a \geq 0$ :
\[ \left(\sqrt{a}\right)^2 = a. \]
C'est une conséquence \textbf{immédiate} de la définition : $\sqrt{a}$ est précisément le nombre positif dont le carré vaut $a$.
\end{propriete}

\begin{propriete}[Racine carrée d'un carré]
Pour tout nombre réel $a$ (positif ou négatif) :
\[ \sqrt{a^2} = |a|. \]
En particulier, si $a \geq 0$, alors $\sqrt{a^2} = a$.
\end{propriete}

\begin{experience}[Démonstration guidée de $\sqrt{a^2} = |a|$]
On veut montrer que $\sqrt{a^2} = |a|$ pour tout réel $a$. D'après la définition, $\sqrt{a^2}$ est le nombre \emph{positif ou nul} dont le carré vaut $a^2$.
\begin{enumerate}
\item Vérifions que $|a|$ est positif ou nul : c'est vrai par définition de la valeur absolue.
\item Vérifions que $|a|^2 = a^2$ : en effet, si $a \geq 0$, $|a| = a$ donc $|a|^2 = a^2$ ; si $a < 0$, $|a| = -a$ donc $|a|^2 = (-a)^2 = a^2$.
\item $|a|$ est donc un nombre positif dont le carré vaut $a^2$. Comme ce nombre est \textbf{unique}, on a bien $\sqrt{a^2} = |a|$.
\end{enumerate}
\end{experience}

\begin{attention}[Le piège classique : croire que $\sqrt{a^2} = a$ pour tout $a$]
C'est \textbf{faux} quand $a$ est négatif ! Contre-exemple :
\[ \sqrt{(-3)^2} = \sqrt{9} = 3 = |-3| \quad \text{et non } -3. \]
Retiens : $\sqrt{a^2} = |a|$, toujours. La valeur absolue est indispensable.
\end{attention}

\begin{exoresolu}[Calculs avec les deux propriétés]
Calculer : a) $\left(\sqrt{7}\right)^2$ ; b) $\sqrt{13^2}$ ; c) $\sqrt{(-5)^2}$ ; d) $\left(\sqrt{2,5}\right)^2$.
\tcblower
a) Comme $7 \geq 0$, $\left(\sqrt{7}\right)^2 = 7$.

b) Comme $13 \geq 0$, $\sqrt{13^2} = 13$.

c) $\sqrt{(-5)^2} = |-5| = 5$ (attention : pas $-5$ !).

d) Comme $2,5 \geq 0$, $\left(\sqrt{2,5}\right)^2 = 2,5$.
\end{exoresolu}

\courssub{Les carrés parfaits : la table à mémoriser}

\begin{definition}[Carré parfait]
Un \cle{carré parfait} est un entier naturel qui est le carré d'un entier naturel. Sa racine carrée est alors un entier.
\end{definition}

\begin{aretenir}[Table des carrés parfaits de $0$ à $15$]
\begin{center}
\begin{tabularx}{\linewidth}{|l|X|X|X|X|X|X|X|X|X|X|X|X|X|X|X|X|}
\hline
\rowcolor{popblueL} $n$ & 0 & 1 & 2 & 3 & 4 & 5 & 6 & 7 & 8 & 9 & 10 & 11 & 12 & 13 & 14 & 15 \\
\hline
$n^2$ & 0 & 1 & 4 & 9 & 16 & 25 & 36 & 49 & 64 & 81 & 100 & 121 & 144 & 169 & 196 & 225 \\
\hline
\rowcolor{popgreenL} $\sqrt{n^2}$ & 0 & 1 & 2 & 3 & 4 & 5 & 6 & 7 & 8 & 9 & 10 & 11 & 12 & 13 & 14 & 15 \\
\hline
\end{tabularx}
\end{center}
Cette table doit être connue \textbf{par cœur} : elle permet de simplifier de nombreux radicaux, comme nous le verrons avec $\sqrt{50}$.
\end{aretenir}

\courssub{Valeurs approchées à la calculatrice}

La plupart des racines carrées ne sont ni entières, ni décimales, ni fractionnaires : ce sont des nombres \emph{irrationnels}, que l'on ne peut écrire exactement qu'avec le symbole $\sqrt{\phantom{a}}$. La calculatrice (touche $\sqrt{\phantom{a}}$) en donne des \textbf{valeurs approchées} :

\begin{center}
\begin{tabular}{|c|c|c|}
\hline
\rowcolor{popblueL} Nombre & Valeur approchée & Arrondi au millième \\
\hline
$\sqrt{2}$ & $1,4142135\ldots$ & $1,414$ \\
\hline
$\sqrt{3}$ & $1,7320508\ldots$ & $1,732$ \\
\hline
$\sqrt{5}$ & $2,2360679\ldots$ & $2,236$ \\
\hline
$\sqrt{50}$ & $7,0710678\ldots$ & $7,071$ \\
\hline
\end{tabular}
\end{center}

\begin{methode}[Lire une valeur approchée à la calculatrice]
\begin{enumerate}
\item Taper la touche $\sqrt{\phantom{a}}$, puis le nombre, puis valider (par exemple $\sqrt{50}$).
\item La calculatrice affiche $7,0710678\ldots$ : c'est une \textbf{valeur approchée}, pas la valeur exacte.
\item Pour arrondir au centième, regarder le troisième chiffre après la virgule : $\sqrt{50} \approx 7,07$ (car le troisième chiffre est $1 < 5$).
\item Rédiger avec le symbole $\approx$ : écrire $\sqrt{50} = 7,07$ est une \textbf{erreur} !
\end{enumerate}
\end{methode}

\begin{attention}[Valeur exacte ou valeur approchée ?]
$\sqrt{50}$ est la valeur \textbf{exacte} ; $7,07$ n'en est qu'une valeur \textbf{approchée}. Dans un calcul exact, on garde le radical ; on ne prend une valeur approchée qu'à la fin, si l'énoncé le demande.
\end{attention}

\courssub{Application : retour au jardin de Moussa}

\begin{exoresolu}[La clôture du jardin]
Le jardin carré de Moussa a une aire de \SI{50}{m^2}.
\begin{enumerate}
\item Donner la valeur exacte du côté $c$ du jardin, en simplifiant le radical.
\item En déduire une valeur approchée de $c$ au centimètre près.
\item Quelle longueur de grillage Moussa doit-il acheter (arrondie au dixième de mètre) ?
\end{enumerate}
\tcblower
\begin{enumerate}
\item Le côté vérifie $c^2 = 50$ avec $c > 0$, donc $c = \sqrt{50}$. Or $50 = 25 \times 2$ et $25$ est un carré parfait :
\[ c = \sqrt{50} = \sqrt{25 \times 2} = \sqrt{25} \times \sqrt{2} = 5\sqrt{2} \ \text{m}. \]
\item À la calculatrice : $c = 5\sqrt{2} \approx 7,071$ m, soit environ $\SI{7,07}{m}$ au centimètre près. On vérifie bien que $7 < \sqrt{50} < 8$ car $49 < 50 < 64$.
\item Le périmètre du jardin vaut $4c = 4 \times 5\sqrt{2} = 20\sqrt{2} \approx 28,28$ m. Moussa doit donc acheter environ $\SI{28,3}{m}$ de grillage (en pratique, il en achètera $29$ m pour avoir une marge !).
\end{enumerate}
\end{exoresolu}

\begin{savaistu}[Les Babyloniens et $\sqrt{2}$]
Il y a près de $4\,000$ ans, bien avant l'invention de la calculatrice, les scribes babyloniens savaient déjà approcher $\sqrt{2}$ avec une précision étonnante. Sur la tablette d'argile \textbf{YBC 7289} (conservée à l'université Yale), on lit un carré avec ses diagonales et, en chiffres cunéiformes, la valeur $1,414213\ldots$ : \textbf{six décimales exactes} de $\sqrt{2}$ ! Les Babyloniens utilisaient un algorithme de moyennes successives, ancêtre des méthodes que ta calculatrice utilise encore aujourd'hui. Plus tard, les Grecs découvrirent avec stupéfaction que $\sqrt{2}$ ne peut s'écrire comme une fraction : la légende dit que ce secret, révélé par un disciple de Pythagore, bouleversa toute leur conception des nombres.
\end{savaistu}

\begin{exercice}[À toi de jouer]
\begin{enumerate}
\item Calculer sans calculatrice : $\sqrt{81}$ ; $\sqrt{144}$ ; $\left(\sqrt{11}\right)^2$ ; $\sqrt{(-7)^2}$.
\item Simplifier : $\sqrt{32} = \sqrt{16 \times 2}$ ; $\sqrt{75} = \sqrt{25 \times 3}$.
\item Un terrain carré a une aire de $200$ m$^2$. Donner la valeur exacte de son côté, puis une valeur approchée au décimètre près.
\end{enumerate}
\end{exercice}

\begin{corrige}[Exercice 1]
\begin{enumerate}
\item $\sqrt{81} = 9$ ; $\sqrt{144} = 12$ ; $\left(\sqrt{11}\right)^2 = 11$ ; $\sqrt{(-7)^2} = |-7| = 7$.
\item $\sqrt{32} = \sqrt{16} \times \sqrt{2} = 4\sqrt{2}$ ; $\sqrt{75} = \sqrt{25} \times \sqrt{3} = 5\sqrt{3}$.
\item Le côté $c$ vérifie $c^2 = 200$, donc $c = \sqrt{200} = \sqrt{100 \times 2} = 10\sqrt{2}$ m. À la calculatrice : $c \approx 14,1$ m au décimètre près.
\end{enumerate}
\end{corrige}

\coursec{L'ensemble des nombres réels}

Dans la section précédente, nous avons découvert que la racine carrée d'un entier qui n'est pas un carré parfait (comme $\sqrt{2}$, $\sqrt{3}$, $\sqrt{5}$) ne peut pas s'écrire sous la forme d'une fraction $\frac{a}{b}$ avec $a,b$ entiers. Ces nombres « échappent » à l'ensemble des rationnels $\mathbb{Q}$. Il est temps d'élargir notre univers numérique pour les accueillir tous.

\courssub{Rappel des ensembles emboîtés : $\mathbb{N}$, $\mathbb{Z}$, $\mathbb{D}$, $\mathbb{Q}$}

Depuis la 6\`eme, tu construis peu à peu la tour des nombres. Rappelle-toi la logique : chaque nouvel ensemble naît d'un besoin pour résoudre des opérations impossibles dans l'ensemble précédent.

\begin{itemize}
    \item \textbf{$\mathbb{N}$ (Entiers naturels) :} $0, 1, 2, 3, \dots$ --- pour compter.
    \item \textbf{$\mathbb{Z}$ (Entiers relatifs) :} $\dots, -3, -2, -1, 0, 1, 2, 3, \dots$ --- pour soustraire n'importe quels naturels ($3-5 = -2$).
    \item \textbf{$\mathbb{D}$ (Nombres décimaux) :} nombres s'écrivant avec un nombre \textbf{fini} de chiffres après la virgule ($-2,5 ; 0,75 ; 3,1415$). C'est l'ensemble des quotients d'un entier par une puissance de $10$. Utile pour les mesures usuelles (argent, longueurs).
    \item \textbf{$\mathbb{Q}$ (Nombres rationnels) :} quotients $\frac{a}{b}$ avec $a \in \mathbb{Z}, b \in \mathbb{Z}^*$. Ils ont une écriture décimale \textbf{illimitée p\'eriodique} (ex: $\frac{1}{3}=0,\overline{3}$) ou \textbf{finie} (ex: $\frac{3}{4}=0,75$). $\mathbb{D} \subset \mathbb{Q}$.
\end{itemize}

On a la chaîne d'inclusions fondamentale :
\[
\mathbb{N} \subset \mathbb{Z} \subset \mathbb{D} \subset \mathbb{Q}
\]

\begin{popfigure}
\begin{tikzpicture}[scale=0.9, every node/.style={font=\small}]
% Cercles emboîtés
\draw[popblue, thick] (0,0) circle (4.5cm) node[above=4.6cm] {$\mathbb{R}$ (Réels)};
\draw[popgreen, thick] (0,0) circle (3.5cm) node[above=3.6cm] {$\mathbb{Q}$ (Rationnels)};
\draw[poporange, thick] (0,0) circle (2.5cm) node[above=2.6cm] {$\mathbb{D}$ (Décimaux)};
\draw[poppurple, thick] (0,0) circle (1.5cm) node[above=1.6cm] {$\mathbb{Z}$ (Entiers relatifs)};
\draw[poppink, thick] (0,0) circle (0.7cm) node[above=0.8cm] {$\mathbb{N}$ (Naturels)};

% Exemples placés stratégiquement
% Dans N
\node[poppink, fill=poppinkL, rounded corners=2pt, inner sep=2pt] at (0, 0.2) {$0, 1, 7, 42$};
% Dans Z \ N
\node[poppurple, fill=poppurpleL, rounded corners=2pt, inner sep=2pt] at (-1.8, 1.0) {$-3, -12$};
% Dans D \ Z
\node[poporange, fill=poporangeL, rounded corners=2pt, inner sep=2pt] at (1.8, 1.8) {$0,75 ; -2,5$};
% Dans Q \ D
\node[popgreen, fill=popgreenL, rounded corners=2pt, inner sep=2pt] at (-2.5, 2.8) {$\frac{2}{3} ; -\frac{7}{5}$};
% Dans R \ Q (Irrationnels)
\node[popblue, fill=popblueL, rounded corners=2pt, inner sep=2pt] at (3.0, 3.5) {$\sqrt{2} ; \pi ; \sqrt{17}$};

% Flèches d'inclusion
\draw[->, popmuted] (-4.8, 4.0) -- (-4.8, -4.0) node[midway, left] {$\subset$};
\end{tikzpicture}
\legende{Diagramme d'inclusions $\mathbb{N} \subset \mathbb{Z} \subset \mathbb{D} \subset \mathbb{Q} \subset \mathbb{R}$ avec exemples.}
\end{popfigure}

\courssub{Les nombres irrationnels : quand la raison ne suffit plus}

Les Grecs anciens (Pythagore et son école) pensaient que \emph{tout} nombre pouvait s'exprimer comme un rapport de deux entiers (une fraction). La découverte de l'irrationalité de $\sqrt{2}$ (vers 500 av. J.-C.) a été un choc fondateur : elle a brisé l'idée que « tout est nombre rationnel ».

\begin{definition}[Nombre irrationnel]
Un nombre est dit \textbf{irrationnel} s'il \textbf{n'appartient pas} à $\mathbb{Q}$. Il ne peut \textbf{pas} s'écrire sous la forme $\frac{a}{b}$ avec $a \in \mathbb{Z}$ et $b \in \mathbb{Z}^*$.
Son écriture décimale est \textbf{illimitée} (infinité de chiffres après la virgule) et \textbf{non p\'eriodique} (aucun bloc de chiffres ne se répète indéfiniment).
On note cet ensemble $\mathbb{R} \setminus \mathbb{Q}$ (ou parfois $\mathbb{Q}^c$).
\end{definition}

\begin{attention}[Piège classique : décimal illimité $\neq$ irrationnel]
$0,333\dots = 0,\overline{3}$ est un décimal \textbf{illimité p\'eriodique}. C'est le rationnel $\frac{1}{3}$.
$0,101001000100001\dots$ (un $0$ de plus à chaque fois) est illimité \textbf{non p\'eriodique} : c'est un irrationnel.
\textbf{Règle :} Illimité \textbf{ET} non p\'eriodique $\iff$ Irrationnel.
\end{attention}

Les exemples célèbres d'irrationnels sont :
\begin{itemize}
    \item $\sqrt{2}, \sqrt{3}, \sqrt{5}, \sqrt{6}, \sqrt{7}, \sqrt{10}, \dots$ (racines carrées d'entiers non carrés parfaits).
    \item $\pi \approx 3,1415926535\dots$ (ratio circonférence/diamètre).
    \item $e \approx 2,71828\dots$ (base des logarithmes népériens, étudié dans les classes supérieures).
\end{itemize}

\begin{experience}[Démonstration guidée : $\sqrt{2}$ est irrationnel]
C'est LA démonstration type par l'absurde, exigible en esprit au BEPC (tu dois en comprendre la logique et pouvoir la restructurer). Suis le raisonnement pas à pas.

\begin{enumerate}
    \item \textbf{Hypothèse (le contraire de ce qu'on veut prouver) :} Supposons que $\sqrt{2}$ \textbf{soit rationnel}.
    \item Alors on peut l'écrire sous forme \textbf{irréductible} : $\sqrt{2} = \frac{a}{b}$ avec $a,b \in \mathbb{N}^*$ et $a \wedge b = 1$ (premiers entre eux, la fraction est simplifiée au maximum).
    \item On élève au carré : $2 = \frac{a^2}{b^2}$ $\iff$ $a^2 = 2b^2$.
    \item \textbf{Conséquence sur $a$ :} $a^2$ est pair (multiple de $2$). Donc $a$ est pair (si un carré est pair, le nombre est pair ; car l'impair au carré reste impair).
    \item On écrit $a = 2k$ avec $k \in \mathbb{N}^*$. On remplace dans l'égalité : $(2k)^2 = 2b^2 \iff 4k^2 = 2b^2 \iff b^2 = 2k^2$.
    \item \textbf{Conséquence sur $b$ :} $b^2$ est pair, donc $b$ est pair.
    \item \textbf{Contradiction :} $a$ est pair \textbf{ET} $b$ est pair. Ils ont donc $2$ comme diviseur commun. Cela contredit l'hypothèse que la fraction $\frac{a}{b}$ est \textbf{irréductible} ($a \wedge b = 1$).
    \item \textbf{Conclusion :} L'hypothèse de départ est fausse. $\sqrt{2}$ \textbf{n'est pas rationnel}. C'est un irrationnel. \qed
\end{enumerate}
\end{experience}

\begin{savaistu}[L'anecdote d'Hippase de Métaponte]
La légende raconte qu'Hippase, disciple de Pythagore, aurait découvert l'irrationalité de $\sqrt{2}$ en étudiant le pentagramme. Les pythagoriciens, pour qui « Tout est nombre » (entier ou fraction), auraient été si bouleversés qu'ils l'auraient jeté par-dessus bord en mer Égée. En réalité, cette découverte a forcé les mathématiques grecques à géométriser l'arithmétique (les « grandeurs » d'Euclide) pendant des siècles.
\end{savaistu}

\courssub{L'ensemble des nombres réels $\mathbb{R}$}

Pour faire des mathématiques sans « trous » dans la droite numérique, on réunit les rationnels et les irrationnels.

\begin{definition}[Ensemble des nombres réels $\mathbb{R}$]
L'ensemble des \textbf{nombres réels}, noté $\mathbb{R}$, est la réunion de l'ensemble des rationnels $\mathbb{Q}$ et de l'ensemble des irrationnels (noté $\mathbb{R} \setminus \mathbb{Q}$) :
\[
\mathbb{R} = \mathbb{Q} \cup (\mathbb{R} \setminus \mathbb{Q})
\]
Tout nombre rationnel est un réel. Tout nombre irrationnel est un réel. Il n'y a pas d'autre possibilité.
\end{definition}

La chaîne d'inclusions complète devient :
\[
\boxed{\mathbb{N} \subset \mathbb{Z} \subset \mathbb{D} \subset \mathbb{Q} \subset \mathbb{R}}
\]

\begin{propriete}[Bijection $\mathbb{R}$ / Droite graduée (Admise)]
Il existe une correspondance \textbf{bijective} (un-à-un) entre l'ensemble des nombres réels $\mathbb{R}$ et l'ensemble des points d'une droite graduée (droite des nombres) :
\begin{itemize}
    \item À \textbf{chaque nombre réel} $x$ correspond \textbf{un unique point} $M$ d'abscisse $x$.
    \item À \textbf{chaque point} $M$ de la droite correspond \textbf{un unique nombre réel} $x$ (son abscisse).
\end{itemize}
C'est pourquoi on parle indifféremment de « nombre réel » ou de « point de la droite réelle ».
\end{propriete}

\begin{popfigure}
\begin{tikzpicture}[scale=1.1, >=Stealth]
% Axe principal
\draw[thick, <->] (-5.5, 0) -- (5.5, 0) node[right] {$x$};
% Graduations principales
\foreach \x in {-5,-4,-3,-2,-1,0,1,2,3,4,5} {
    \draw (\x, 0.1) -- (\x, -0.1) node[below=2pt] {$\x$};
}
% Points spéciaux
% -2
\fill[popblue] (-2, 0) circle (2.5pt);
\node[popblue, above] at (-2, 0.3) {$-2 \in \mathbb{Z}$};
% 1/2 = 0.5
\fill[popgreen] (0.5, 0) circle (2.5pt);
\node[popgreen, above] at (0.5, 0.7) {$\frac{1}{2} \in \mathbb{Q}$};
% sqrt(2) ~ 1.414
\fill[poporange] (1.414, 0) circle (2.5pt);
\node[poporange, above] at (1.414, 1.1) {$\sqrt{2} \notin \mathbb{Q}$};
% pi ~ 3.14159
\fill[poppurple] (3.14159, 0) circle (2.5pt);
\node[poppurple, above] at (3.14159, 1.5) {$\pi \notin \mathbb{Q}$};

% Flèches indicatives
\draw[popmuted, dashed] (-2, 0.5) -- (-2, 0.15);
\draw[popmuted, dashed] (0.5, 0.9) -- (0.5, 0.15);
\draw[popmuted, dashed] (1.414, 1.3) -- (1.414, 0.15);
\draw[popmuted, dashed] (3.14159, 1.7) -- (3.14159, 0.15);
\end{tikzpicture}
\legende{Représentation sur la droite graduée : les rationnels ($-2, \frac{1}{2}$) et les irrationnels ($\sqrt{2}, \pi$) sont tous des points de $\mathbb{R}$.}
\end{popfigure}

\courssub{Classer un nombre dans les ensembles : méthode et exemples}

Savoir dire « ce nombre est dans $\mathbb{Z}$ mais pas dans $\mathbb{N}$ » ou « ce nombre est irrationnel donc dans $\mathbb{R}$ seulement » est une compétence clé du BEPC.

\begin{methode}[Classer un nombre réel]
\begin{enumerate}
    \item \textbf{Simplifie} l'expression si possible (ex: $\sqrt{16}=4$, $-\frac{4}{2}=-2$, $\sqrt{8}=2\sqrt{2}$).
    \item \textbf{Identifie la nature} de l'écriture finale :
    \begin{itemize}
        \item Entier naturel ($\ge 0$) $\to \mathbb{N}$.
        \item Entier relatif (avec signe $-$) $\to \mathbb{Z}$.
        \item Décimal fini (virgule, finit par des zéros) $\to \mathbb{D}$.
        \item Fraction $\frac{a}{b}$ ou décimal illimité \textbf{périodique} $\to \mathbb{Q}$.
        \item Racine carrée d'un non-carré ($\sqrt{n}$), $\pi$, écriture décimale illimitée \textbf{non périodique} $\to$ Irrationnel ($\mathbb{R} \setminus \mathbb{Q}$).
    \end{itemize}
    \item \textbf{Remonte la chaîne} : si $x \in \mathbb{N}$, alors $x \in \mathbb{Z}, \mathbb{D}, \mathbb{Q}, \mathbb{R}$. Si $x \in \mathbb{Q}$ seulement, alors $x \in \mathbb{R}$ mais $x \notin \mathbb{N}, \mathbb{Z}, \mathbb{D}$.
\end{enumerate}
\end{methode}

\begin{exemplebox}[Classement de $7 ; -\frac{4}{2} ; \sqrt{16} ; \sqrt{17} ; 3,14$]
Appliquons la méthode à chaque nombre.

\begin{center}
\begin{tabularx}{\linewidth}{|c|X|c|c|c|c|c|}\hline
\textbf{Nombre} & \textbf{Simplification / Nature} & $\mathbb{N}$ & $\mathbb{Z}$ & $\mathbb{D}$ & $\mathbb{Q}$ & $\mathbb{R}$ \\ \hline
$7$ & Entier naturel & \textbf{Oui} & Oui & Oui & Oui & Oui \\ \hline
$-\frac{4}{2}$ & $= -2$ (entier relatif négatif) & Non & \textbf{Oui} & Oui & Oui & Oui \\ \hline
$\sqrt{16}$ & $= 4$ (carré parfait) & \textbf{Oui} & Oui & Oui & Oui & Oui \\ \hline
$\sqrt{17}$ & $17$ n'est pas un carré parfait $\to$ irrationnel & Non & Non & Non & Non & \textbf{Oui} \\ \hline
$3,14$ & Décimal \textbf{fini} (2 chiffres après la virgule) $\to \frac{314}{100}=\frac{157}{50}$ & Non & Non & \textbf{Oui} & Oui & Oui \\ \hline
\end{tabularx}
\end{center}

\textbf{Attention :} $3,14 \neq \pi$ ! $\pi \approx 3,14159\dots$ est irrationnel. $3,14$ est une approximation décimale \textbf{rationnelle} (donc dans $\mathbb{D}$).
\end{exemplebox}

\begin{exoresolu}[Classer et justifier]
Classe les nombres suivants dans le(s) plus petit(s) ensemble(s) approprié(s) parmi $\mathbb{N}, \mathbb{Z}, \mathbb{D}, \mathbb{Q}, \mathbb{R} \setminus \mathbb{Q}$ :
\[
A = \frac{22}{7} \quad ; \quad B = -\sqrt{9} \quad ; \quad C = 0,\overline{142857} \quad ; \quad D = \sqrt{2} + \sqrt{8} \quad ; \quad E = \frac{\pi}{2}
\]

\tcblower
\textbf{Solution détaillée :}

\begin{enumerate}
    \item \textbf{$A = \frac{22}{7}$} : C'est un quotient d'entiers ($22 \in \mathbb{Z}, 7 \in \mathbb{Z}^*$). C'est un \textbf{rationnel}. Son écriture décimale est $3,\overline{142857}$ (période $142857$). Il n'est pas décimal (période non nulle). \textbf{Plus petit ensemble : $\mathbb{Q}$}.
    \item \textbf{$B = -\sqrt{9}$} : $\sqrt{9} = 3$, donc $B = -3$. C'est un entier relatif négatif. \textbf{Plus petit ensemble : $\mathbb{Z}$}.
    \item \textbf{$C = 0,\overline{142857}$} : Décimal illimité \textbf{périodique} (la barre au-dessus indique la période). C'est la définition même d'un rationnel. Ici $C = \frac{1}{7}$. \textbf{Plus petit ensemble : $\mathbb{Q}$}.
    \item \textbf{$D = \sqrt{2} + \sqrt{8}$} : On simplifie $\sqrt{8} = \sqrt{4 \times 2} = 2\sqrt{2}$. Donc $D = \sqrt{2} + 2\sqrt{2} = 3\sqrt{2} = \sqrt{18}$. $18$ n'est pas un carré parfait, $\sqrt{18}$ est irrationnel. \textbf{Plus petit ensemble : $\mathbb{R} \setminus \mathbb{Q}$ (Irrationnel)}.
    \item \textbf{$E = \frac{\pi}{2}$} : $\pi$ est irrationnel. Le quotient d'un irrationnel par un entier non nul ($2$) reste irrationnel (sinon $\pi = 2E$ serait rationnel, contradiction). \textbf{Plus petit ensemble : $\mathbb{R} \setminus \mathbb{Q}$}.
\end{enumerate}
\end{exoresolu}

\courssub{Complément (Culture mathématique) : La densité de $\mathbb{Q}$ dans $\mathbb{R}$}

\begin{savaistu}[Densité des rationnels --- Hors programme BEPC, utile pour la suite]
Une propriété fascinante (et admise au lycée) : \textbf{entre deux réels distincts, il existe toujours un rationnel}.
Formellement : $\forall (x,y) \in \mathbb{R}^2,\ x < y \implies \exists r \in \mathbb{Q},\ x < r < y$.
Cela signifie que les rationnels sont « partout » sur la droite, ils la « remplissent » de manière dense, même s'ils sont en quelque sorte « moins nombreux » que les réels (infini dénombrable vs non dénombrable, notion des classes supérieures). C'est pour cela qu'on peut approcher n'importe quel réel (comme $\pi$ ou $\sqrt{2}$) aussi près qu'on veut par un nombre décimal (donc rationnel) : $3,14 ; 3,141 ; 3,1415 ; \dots$
\end{savaistu}

\begin{aretenir}[Bilan de la section]
\begin{itemize}
    \item Les ensembles s'emboîtent : $\mathbb{N} \subset \mathbb{Z} \subset \mathbb{D} \subset \mathbb{Q} \subset \mathbb{R}$.
    \item Un \textbf{irrationnel} n'est pas un quotient d'entiers. Sa décimale est illimitée \textbf{non périodique}. Exemples : $\sqrt{n}$ (si $n$ n'est pas un carré parfait), $\pi$.
    \item $\mathbb{R} = \mathbb{Q} \cup \{\text{irrationnels}\}$. Tout point de la droite graduée est un réel, et réciproquement.
    \item Pour classer un nombre : \textbf{simplifie d'abord} ($\sqrt{16}=4$, $\frac{-6}{2}=-3$, $\sqrt{8}=2\sqrt{2}$), puis identifie sa nature finale.
    \item \textbf{Preuve type :} L'irrationalité de $\sqrt{2}$ se démontre par l'absurde (pair/impair, fraction irréductible).
\end{itemize}
\end{aretenir}

\coursec{Somme et produit de réels avec radicaux}

Dans la section précédente, nous avons découvert l'ensemble $\mathbb{R}$ des nombres réels, qui contient à la fois les nombres rationnels (comme $\dfrac{3}{4}$ ou $-2$) et les nombres irrationnels (comme $\sqrt{2}$, $\sqrt{3}$ ou $\pi$). Nous savons maintenant que tout réel positif possède une racine carrée. Mais une question naturelle se pose : \emph{comment calcule-t-on avec ces nouveaux nombres ?} Peut-on additionner, multiplier, développer des expressions contenant des radicaux ? C'est précisément l'objet de cette section : apprendre les règles de calcul avec les racines carrées, afin de simplifier, réduire et développer des expressions comme $3\sqrt{50} - 2\sqrt{18} + \sqrt{8}$ ou $(2\sqrt{3}+1)^2$. Ces techniques sont indispensables au BEPC et dans toute la suite de ta scolarité.

\courssub{Le produit de deux racines carrées}

\textbf{Une intuition géométrique.} Imaginons un rectangle dont les côtés mesurent $\sqrt{8}$ cm et $\sqrt{2}$ cm. Son aire vaut $\sqrt{8} \times \sqrt{2}$ cm$^2$. Mais on peut aussi remarquer que $\sqrt{8} \times \sqrt{2}$ « devrait » valoir $\sqrt{16} = 4$, car $8 \times 2 = 16$. Cette intuition est juste, et nous allons la démontrer rigoureusement.

\begin{propriete}[Produit de deux racines carrées]
Pour tous nombres réels \cle{positifs} $a$ et $b$ :
\[ \sqrt{a \times b} = \sqrt{a} \times \sqrt{b} \]
\end{propriete}

\begin{experience}[Démonstration de la propriété $\sqrt{ab} = \sqrt{a} \times \sqrt{b}$]
Soient $a \geq 0$ et $b \geq 0$. Nous voulons montrer que $\sqrt{a}\times\sqrt{b}$ est la racine carrée de $ab$. D'après la définition de la racine carrée, il suffit de vérifier deux choses : que ce nombre est \textbf{positif}, et que son \textbf{carré} vaut $ab$.
\begin{enumerate}
\item \textbf{Signe.} Comme $a \geq 0$ et $b \geq 0$, on a $\sqrt{a} \geq 0$ et $\sqrt{b} \geq 0$. Le produit de deux nombres positifs est positif, donc $\sqrt{a}\times\sqrt{b} \geq 0$.
\item \textbf{Élévation au carré.} Calculons :
\begin{align*}
\left(\sqrt{a} \times \sqrt{b}\right)^2 &= \left(\sqrt{a}\right)^2 \times \left(\sqrt{b}\right)^2 \quad \text{(carré d'un produit)}\\
&= a \times b \quad \text{(par définition de la racine carrée)}
\end{align*}
\end{enumerate}
Le nombre $\sqrt{a}\times\sqrt{b}$ est donc le nombre positif dont le carré vaut $ab$ : c'est exactement $\sqrt{ab}$. \hfill$\blacksquare$
\end{experience}

\begin{popfigure}
\begin{center}
\begin{tikzpicture}[scale=0.9]
% Rectangle de côtés sqrt(8) ≈ 2.83 et sqrt(2) ≈ 1.41
\draw[fill=popblueL, draw=popblue, thick] (0,0) rectangle (2.83,1.41);
\draw[<->, popdark] (0,-0.35) -- (2.83,-0.35) node[midway, below] {$\sqrt{8}$};
\draw[<->, popdark] (3.13,0) -- (3.13,1.41) node[midway, right] {$\sqrt{2}$};
\node at (1.415,0.705) {Aire $= \sqrt{8}\times\sqrt{2} = \sqrt{16} = 4$};
% Petit carré unité pour comparaison
\draw[fill=poporangeL, draw=poporange, thick] (5.2,0) rectangle (6.2,1);
\node[below] at (5.7,0) {carré d'aire $1$};
\end{tikzpicture}
\end{center}
\legende{Un rectangle de côtés $\sqrt{8}$ et $\sqrt{2}$ a une aire de $\sqrt{8}\times\sqrt{2} = \sqrt{16} = 4$ unités d'aire, soit exactement $4$ carreaux unité.}
\end{popfigure}

\begin{exemplebox}[Calculs directs de produits]
\begin{itemize}
\item $\sqrt{2} \times \sqrt{8} = \sqrt{2 \times 8} = \sqrt{16} = 4$ ;
\item $\sqrt{3} \times \sqrt{12} = \sqrt{36} = 6$ ;
\item $\sqrt{5} \times \sqrt{5} = \sqrt{25} = 5$ (on retrouve $(\sqrt{5})^2 = 5$).
\end{itemize}
\end{exemplebox}

\begin{attention}[La condition $a, b \geq 0$ est indispensable]
La formule $\sqrt{ab} = \sqrt{a}\times\sqrt{b}$ n'est valable que pour des nombres \textbf{positifs}. Par exemple, $\sqrt{(-2)\times(-8)} = \sqrt{16} = 4$, alors que $\sqrt{-2}$ et $\sqrt{-8}$ n'existent pas dans $\mathbb{R}$ ! N'écris jamais la formule sans vérifier que les nombres sous les radicaux sont positifs.
\end{attention}

\courssub{Simplification d'un radical : extraire les carrés parfaits}

La propriété précédente, lue de droite à gauche, devient un outil formidable : si le nombre sous le radical contient un \cle{carré parfait} en facteur, on peut le « sortir » de la racine. Rappelons les carrés parfaits à connaître par cœur : $1, 4, 9, 16, 25, 36, 49, 64, 81, 100, 121, 144, 169, 196, 225$.

\begin{methode}[Simplifier un radical]
Pour simplifier $\sqrt{n}$ où $n$ est un entier positif :
\begin{enumerate}
\item \textbf{Décomposer} $n$ en un produit faisant apparaître le plus grand carré parfait possible : $n = a^2 \times b$ ;
\item \textbf{Appliquer} la propriété du produit : $\sqrt{n} = \sqrt{a^2 \times b} = \sqrt{a^2} \times \sqrt{b}$ ;
\item \textbf{Conclure} : $\sqrt{n} = a\sqrt{b}$, où $b$ ne contient plus aucun carré parfait.
\end{enumerate}
\end{methode}

\begin{exemplebox}[Simplification de $\sqrt{72}$]
On cherche un carré parfait qui divise $72$. On remarque que $72 = 36 \times 2$, et $36 = 6^2$ est un carré parfait. Donc :
\[ \sqrt{72} = \sqrt{36 \times 2} = \sqrt{36} \times \sqrt{2} = 6\sqrt{2}. \]
On aurait pu écrire $72 = 9 \times 8$, ce qui donne $\sqrt{72} = 3\sqrt{8}$, mais il aurait fallu continuer : $\sqrt{8} = \sqrt{4\times 2} = 2\sqrt{2}$, d'où $\sqrt{72} = 3 \times 2\sqrt{2} = 6\sqrt{2}$. Choisir d'emblée le \emph{plus grand} carré parfait fait gagner du temps.
\end{exemplebox}

\begin{exemplebox}[Autres simplifications à retenir]
\begin{itemize}
\item $\sqrt{50} = \sqrt{25 \times 2} = 5\sqrt{2}$ ;
\item $\sqrt{18} = \sqrt{9 \times 2} = 3\sqrt{2}$ ;
\item $\sqrt{12} = \sqrt{4 \times 3} = 2\sqrt{3}$ ;
\item $\sqrt{200} = \sqrt{100 \times 2} = 10\sqrt{2}$.
\end{itemize}
\end{exemplebox}

\courssub{Somme et différence de radicaux : les radicaux semblables}

\textbf{Intuition.} Les radicaux se comportent comme des « sortes de fruits » : $3\sqrt{2} + 5\sqrt{2}$, c'est comme $3$ mangues $+ 5$ mangues $= 8$ mangues, donc $8\sqrt{2}$. En revanche, $3\sqrt{2} + 5\sqrt{3}$ ne se réduit pas, tout comme on ne peut pas additionner des mangues et des avocats !

\begin{definition}[Radicaux semblables]
Deux radicaux sont dits \cle{semblables} lorsque, après simplification, ils contiennent \textbf{la même racine carrée}. Par exemple, $\sqrt{50} = 5\sqrt{2}$ et $\sqrt{18} = 3\sqrt{2}$ sont semblables car ils contiennent tous les deux $\sqrt{2}$.
\end{definition}

\begin{propriete}[Réduction d'une somme de radicaux semblables]
Pour tous réels $a$, $c$ et tout réel $b \geq 0$ :
\[ a\sqrt{b} + c\sqrt{b} = (a+c)\sqrt{b} \]
C'est une simple factorisation par $\sqrt{b}$, qui est un facteur commun.
\end{propriete}

\begin{methode}[Réduire une somme de radicaux]
\begin{enumerate}
\item \textbf{Simplifier} chaque radical en extrayant les carrés parfaits ;
\item \textbf{Repérer} les radicaux semblables (même racine) ;
\item \textbf{Regrouper} et additionner leurs coefficients ;
\item \textbf{Laisser séparés} les radicaux non semblables.
\end{enumerate}
\end{methode}

\begin{exoresolu}[Réduire $3\sqrt{50} - 2\sqrt{18} + \sqrt{8}$]
Réduire l'expression $A = 3\sqrt{50} - 2\sqrt{18} + \sqrt{8}$.
\tcblower
\textbf{Étape 1 : simplifier chaque radical.}
\[ \sqrt{50} = \sqrt{25 \times 2} = 5\sqrt{2}, \qquad \sqrt{18} = \sqrt{9 \times 2} = 3\sqrt{2}, \qquad \sqrt{8} = \sqrt{4 \times 2} = 2\sqrt{2}. \]
\textbf{Étape 2 : remplacer dans l'expression.}
\begin{align*}
A &= 3 \times 5\sqrt{2} - 2 \times 3\sqrt{2} + 2\sqrt{2}\\
&= 15\sqrt{2} - 6\sqrt{2} + 2\sqrt{2}
\end{align*}
\textbf{Étape 3 : regrouper les radicaux semblables} (tous contiennent $\sqrt{2}$) :
\[ A = (15 - 6 + 2)\sqrt{2} = 11\sqrt{2}. \]
\textbf{Conclusion :} $3\sqrt{50} - 2\sqrt{18} + \sqrt{8} = 11\sqrt{2}$.
\end{exoresolu}

\begin{attention}[Erreur fréquente : $\sqrt{a} + \sqrt{b} \neq \sqrt{a+b}$]
Il n'existe \textbf{aucune formule} pour simplifier la somme de deux racines carrées non semblables ! Contre-exemple éclatant :
\[ \sqrt{9} + \sqrt{16} = 3 + 4 = 7 \quad \text{alors que} \quad \sqrt{9+16} = \sqrt{25} = 5. \]
Or $7 \neq 5$. De même, $\sqrt{a} - \sqrt{b} \neq \sqrt{a-b}$ en général : $\sqrt{16} - \sqrt{9} = 1$ alors que $\sqrt{16-9} = \sqrt{7} \approx 2,6$. Ne « fusionne » jamais deux radicaux sous une seule racine lors d'une addition ou d'une soustraction !
\end{attention}

\begin{experience}[Pourquoi $\sqrt{a} + \sqrt{b}$ ne peut pas être réduit en général]
Justifions par les carrés que $\sqrt{a} + \sqrt{b} \neq \sqrt{a+b}$ pour $a, b > 0$. Comparons les carrés des deux membres :
\[ \left(\sqrt{a} + \sqrt{b}\right)^2 = a + b + 2\sqrt{ab} \qquad \text{et} \qquad \left(\sqrt{a+b}\right)^2 = a + b. \]
Comme $a > 0$ et $b > 0$, on a $2\sqrt{ab} > 0$, donc $a + b + 2\sqrt{ab} > a + b$. Les deux carrés sont différents, et comme les deux nombres sont positifs, les nombres eux-mêmes sont différents : $\sqrt{a} + \sqrt{b} > \sqrt{a+b}$. \hfill$\blacksquare$
\end{experience}

\courssub{Développements et identités remarquables avec radicaux}

Toutes les règles du calcul littéral (distributivité, identités remarquables) restent valables dans $\mathbb{R}$, donc avec les radicaux. Rappelons les trois identités remarquables :
\[ (x+y)^2 = x^2 + 2xy + y^2, \quad (x-y)^2 = x^2 - 2xy + y^2, \quad (x+y)(x-y) = x^2 - y^2. \]

\begin{propriete}[Carré d'une somme de radicaux]
Pour tous réels $a \geq 0$ et $b \geq 0$ :
\[ \left(\sqrt{a} + \sqrt{b}\right)^2 = a + b + 2\sqrt{ab} \]
\end{propriete}

En effet, en posant $x = \sqrt{a}$ et $y = \sqrt{b}$ dans la première identité remarquable :
\[ \left(\sqrt{a} + \sqrt{b}\right)^2 = \left(\sqrt{a}\right)^2 + 2\sqrt{a}\sqrt{b} + \left(\sqrt{b}\right)^2 = a + 2\sqrt{ab} + b. \]

\begin{exoresolu}[Développer $(2\sqrt{3} + 1)^2$]
Développer et réduire $B = (2\sqrt{3} + 1)^2$.
\tcblower
On applique l'identité $(x+y)^2 = x^2 + 2xy + y^2$ avec $x = 2\sqrt{3}$ et $y = 1$ :
\begin{align*}
B &= \left(2\sqrt{3}\right)^2 + 2 \times 2\sqrt{3} \times 1 + 1^2\\
&= 4 \times 3 + 4\sqrt{3} + 1 \quad \text{car } \left(2\sqrt{3}\right)^2 = 2^2 \times \left(\sqrt{3}\right)^2 = 4 \times 3 = 12\\
&= 12 + 4\sqrt{3} + 1\\
&= 13 + 4\sqrt{3}.
\end{align*}
\textbf{Conclusion :} $(2\sqrt{3}+1)^2 = 13 + 4\sqrt{3}$.
\end{exoresolu}

\begin{attention}[Pièges de rédaction dans les développements]
\begin{itemize}
\item $\left(2\sqrt{3}\right)^2 = 4 \times 3 = 12$, et \textbf{pas} $2 \times 3 = 6$ : le coefficient $2$ doit lui aussi être élevé au carré !
\item Ne pas oublier le \textbf{double produit} : $(\sqrt{5} + \sqrt{2})^2$ n'est pas égal à $5 + 2 = 7$, mais à $7 + 2\sqrt{10}$.
\end{itemize}
\end{attention}

\courssub{Le produit de conjugués : $(\sqrt{a} - \sqrt{b})(\sqrt{a} + \sqrt{b})$}

\begin{definition}[Expressions conjuguées]
Deux expressions de la forme $\sqrt{a} - \sqrt{b}$ et $\sqrt{a} + \sqrt{b}$ sont dites \cle{conjuguées} l'une de l'autre. Elles ne diffèrent que par le signe entre les deux termes.
\end{definition}

\begin{propriete}[Produit de deux expressions conjuguées]
Pour tous réels $a \geq 0$ et $b \geq 0$ :
\[ \left(\sqrt{a} - \sqrt{b}\right)\left(\sqrt{a} + \sqrt{b}\right) = a - b \]
En particulier, le résultat ne contient \textbf{plus aucun radical} !
\end{propriete}

C'est une conséquence directe de la troisième identité remarquable $(x-y)(x+y) = x^2 - y^2$ avec $x = \sqrt{a}$ et $y = \sqrt{b}$ :
\[ \left(\sqrt{a} - \sqrt{b}\right)\left(\sqrt{a} + \sqrt{b}\right) = \left(\sqrt{a}\right)^2 - \left(\sqrt{b}\right)^2 = a - b. \]

\begin{exemplebox}[Produits de conjugués]
\begin{itemize}
\item $\left(\sqrt{7} - \sqrt{3}\right)\left(\sqrt{7} + \sqrt{3}\right) = 7 - 3 = 4$ ;
\item $\left(5 - \sqrt{2}\right)\left(5 + \sqrt{2}\right) = 25 - 2 = 23$ ;
\item $\left(2\sqrt{3} - 1\right)\left(2\sqrt{3} + 1\right) = \left(2\sqrt{3}\right)^2 - 1^2 = 12 - 1 = 11$.
\end{itemize}
\end{exemplebox}

\begin{savaistu}[À quoi servent les conjugués ?]
Le produit de conjugués permet de « faire disparaître » les racines carrées. Cette astuce sera précieuse dans la section suivante pour écrire un quotient comme $\dfrac{3}{\sqrt{5} - \sqrt{2}}$ sans radical au dénominateur : il suffira de multiplier numérateur et dénominateur par l'expression conjuguée $\sqrt{5} + \sqrt{2}$. Les ingénieurs qui dimensionnent les réseaux électriques ou les antennes de téléphonie mobile utilisent constamment ces calculs avec des nombres irrationnels !
\end{savaistu}

\begin{aretenir}[L'essentiel de la section]
\begin{itemize}
\item \textbf{Produit :} pour $a, b \geq 0$, $\sqrt{a \times b} = \sqrt{a} \times \sqrt{b}$.
\item \textbf{Simplification :} extraire les carrés parfaits, par exemple $\sqrt{72} = \sqrt{36 \times 2} = 6\sqrt{2}$.
\item \textbf{Somme :} seuls les radicaux \textbf{semblables} se réduisent : $a\sqrt{b} + c\sqrt{b} = (a+c)\sqrt{b}$.
\item \textbf{Interdit :} $\sqrt{a} + \sqrt{b} \neq \sqrt{a+b}$ en général.
\item \textbf{Identités :} $(\sqrt{a} + \sqrt{b})^2 = a + b + 2\sqrt{ab}$ et $\left(\sqrt{a} - \sqrt{b}\right)\left(\sqrt{a} + \sqrt{b}\right) = a - b$.
\end{itemize}
\end{aretenir}

\begin{exercice}[Pour s'entraîner]
\begin{enumerate}
\item Simplifier : $\sqrt{32}$, $\sqrt{75}$, $\sqrt{98}$, $\sqrt{300}$.
\item Réduire : $C = 5\sqrt{27} - \sqrt{12} + 2\sqrt{48}$.
\item Développer et réduire : $D = (\sqrt{5} - 3)^2$ et $E = (\sqrt{6} + \sqrt{2})(\sqrt{6} - \sqrt{2})$.
\end{enumerate}
\end{exercice}

\begin{corrige}[Exercice « Pour s'entraîner »]
\begin{enumerate}
\item $\sqrt{32} = \sqrt{16 \times 2} = 4\sqrt{2}$ ; $\sqrt{75} = \sqrt{25 \times 3} = 5\sqrt{3}$ ; $\sqrt{98} = \sqrt{49 \times 2} = 7\sqrt{2}$ ; $\sqrt{300} = \sqrt{100 \times 3} = 10\sqrt{3}$.
\item $\sqrt{27} = 3\sqrt{3}$, $\sqrt{12} = 2\sqrt{3}$, $\sqrt{48} = 4\sqrt{3}$. Donc $C = 5 \times 3\sqrt{3} - 2\sqrt{3} + 2 \times 4\sqrt{3} = 15\sqrt{3} - 2\sqrt{3} + 8\sqrt{3} = 21\sqrt{3}$.
\item $D = (\sqrt{5})^2 - 2 \times \sqrt{5} \times 3 + 3^2 = 5 - 6\sqrt{5} + 9 = 14 - 6\sqrt{5}$. \quad $E = (\sqrt{6})^2 - (\sqrt{2})^2 = 6 - 2 = 4$.
\end{enumerate}
\end{corrige}

\coursec{Quotient de réels avec radicaux}

Dans la section précédente, nous avons maîtrisé l'addition, la soustraction et la multiplication d'expressions contenant des radicaux. Nous savons maintenant simplifier des sommes comme $\sqrt{12}+\sqrt{27}$ ou des produits comme $\sqrt{6}\times\sqrt{10}$. Il reste une opération fondamentale : la division. Que faire face à une expression comme $\dfrac{5}{\sqrt{2}}$ ou $\dfrac{4}{\sqrt{5}-1}$ ? En mathématiques, et particulièrement au BEPC, \textbf{on ne laisse jamais un radical au dénominateur dans un résultat final}. Cette section t'apprendra à \textbf{rationaliser} ces dénominateurs pour écrire tes résultats sous une forme conventionnelle, exacte et acceptée à l'examen.

\courssub{Quotient de deux racines carrées : la propriété fondamentale}

Commençons par le cas le plus simple : le quotient de deux radicaux, ou une racine carrée d'un quotient. L'intuition géométrique est simple : si un carré d'aire $a$ a pour côté $\sqrt{a}$, et un carré d'aire $b$ a pour côté $\sqrt{b}$, alors le rapport de leurs côtés est-il le côté d'un carré d'aire $\frac{a}{b}$ ? La réponse est oui, sous réserve que $b$ ne soit pas nul.

\begin{propriete}[Quotient de racines carrées — Racine d'un quotient]
Soient $a$ un réel positif ou nul ($a \ge 0$) et $b$ un réel \textbf{strictement positif} ($b > 0$).
On a l'égalité :
\[
\sqrt{\frac{a}{b}} = \frac{\sqrt{a}}{\sqrt{b}}
\]
\end{propriete}

\begin{experience}[Démonstration par élévation au carré]
Pour démontrer cette propriété, nous utilisons la définition même de la racine carrée : $\sqrt{x}$ est l'unique réel positif dont le carré est $x$.
Posons $x = \frac{\sqrt{a}}{\sqrt{b}}$. Calculons $x^2$ :
\begin{align*}
x^2 &= \left( \frac{\sqrt{a}}{\sqrt{b}} \right)^2 \\
    &= \frac{(\sqrt{a})^2}{(\sqrt{b})^2} \quad \text{(propriété du quotient des puissances)} \\
    &= \frac{a}{b} \quad \text{(car $a\ge 0, b>0$)}
\end{align*}
Or, par définition, $\sqrt{\frac{a}{b}}$ est l'unique réel positif dont le carré est $\frac{a}{b}$.
Comme $a \ge 0$ et $b > 0$, le quotient $\frac{\sqrt{a}}{\sqrt{b}}$ est bien positif ou nul.
Donc $\frac{\sqrt{a}}{\sqrt{b}} = \sqrt{\frac{a}{b}}$. \hfill $\blacksquare$
\end{experience}

\begin{attention}[Conditions d'application obligatoires]
Au BEPC, tu dois \textbf{toujours vérifier et mentionner les conditions} :
\begin{itemize}
    \item $a \ge 0$ (le numérateur du radicande est positif ou nul) ;
    \item $b > 0$ (le dénominateur est \textbf{strictement positif}, car on ne divise pas par zéro et $\sqrt{b}$ est au dénominateur).
\end{itemize}
Si $b < 0$, l'expression $\frac{\sqrt{a}}{\sqrt{b}}$ n'a pas de sens dans $\mathbb{R}$ (racine d'un négatif), alors que $\sqrt{\frac{a}{b}}$ n'en a pas non plus. La propriété ne s'applique qu'en $\mathbb{R}_+^*$ pour $b$.
\end{attention}

\begin{exemplebox}[Application directe]
Simplifier $A = \frac{\sqrt{75}}{\sqrt{3}}$ et $B = \sqrt{\frac{50}{2}}$.

\textbf{Solution pour $A$ :}
On utilise la propriété « de droite à gauche » : $\frac{\sqrt{a}}{\sqrt{b}} = \sqrt{\frac{a}{b}}$.
\[
A = \frac{\sqrt{75}}{\sqrt{3}} = \sqrt{\frac{75}{3}} = \sqrt{25} = 5.
\]

\textbf{Solution pour $B$ :}
On utilise la propriété « de gauche à droite » : $\sqrt{\frac{a}{b}} = \frac{\sqrt{a}}{\sqrt{b}}$.
\[
B = \sqrt{\frac{50}{2}} = \sqrt{25} = 5.
\]
\textit{Remarque :} Ici, simplifier la fraction sous le radical \textbf{avant} de prendre la racine est plus rapide. C'est une habitude à prendre : \textbf{simplifie toujours la fraction $\frac{a}{b}$ avant d'appliquer la racine}.
\end{exemplebox}

\courssub{Rationaliser un dénominateur simple (monôme)}

Lorsque le dénominateur est un seul radical, comme $\frac{5}{\sqrt{2}}$, la convention mathématique exige de « faire monter » le radical au numérateur. On utilise l'astuce : multiplier le numérateur et le dénominateur par $\sqrt{2}$ (car $\sqrt{2}\times\sqrt{2}=2$).

\begin{propriete}[Rationalisation d'un dénominateur simple]
Pour tout réel $a$ et tout réel $b > 0$ :
\[
\frac{a}{\sqrt{b}} = \frac{a \times \sqrt{b}}{\sqrt{b} \times \sqrt{b}} = \frac{a\sqrt{b}}{b}
\]
Le dénominateur devient un entier (ou un rationnel), il est \textbf{rationalisé}.
\end{propriete}

\begin{methode}[Rationaliser un dénominateur de la forme $\sqrt{b}$]
\begin{enumerate}
    \item Identifier le radical $\sqrt{b}$ au dénominateur.
    \item Multiplier le numérateur \textbf{et} le dénominateur par $\sqrt{b}$.
    \item Simplifier le dénominateur : $\sqrt{b}\times\sqrt{b} = b$.
    \item Réduire la fraction finale si possible (diviser numérateur et dénominateur par leur PGCD).
\end{enumerate}
\end{methode}

\begin{exemplebox}[Exemples classiques — Dénominateur simple]
\begin{enumerate}
    \item $\dfrac{5}{\sqrt{2}} = \dfrac{5 \times \sqrt{2}}{\sqrt{2} \times \sqrt{2}} = \dfrac{5\sqrt{2}}{2}$.
    \item $\dfrac{6}{\sqrt{3}} = \dfrac{6\sqrt{3}}{3} = 2\sqrt{3}$ \quad (on a réduit par $3$).
    \item $\dfrac{\sqrt{5}}{\sqrt{7}} = \dfrac{\sqrt{5}\times\sqrt{7}}{7} = \dfrac{\sqrt{35}}{7}$.
    \item $\dfrac{12}{\sqrt{8}} = \dfrac{12\sqrt{8}}{8} = \dfrac{3\sqrt{8}}{2} = \dfrac{3\times 2\sqrt{2}}{2} = 3\sqrt{2}$.
    \textit{Astuce :} On peut aussi simplifier $\sqrt{8}=2\sqrt{2}$ dès le départ : $\frac{12}{2\sqrt{2}} = \frac{6}{\sqrt{2}} = 3\sqrt{2}$.
\end{enumerate}
\end{exemplebox}

\begin{popfigure}
\begin{tikzpicture}[scale=0.7, >=Stealth]
% Triangle rectangle ABC, rectangle en B
% AB = 6 (vertical), BC = 2*sqrt(3) ~ 3.464 (horizontal)
% Angle A = 30 deg, Angle C = 60 deg
\coordinate (B) at (0,0);
\coordinate (A) at (0,6);
\coordinate (C) at (3.464101615,0); % 2*sqrt(3)

% Triangle
\draw[thick, popdark] (A) -- node[left] {$AB = 6$} (B) -- node[below] {$BC = 2\sqrt{3}$} (C) -- node[right, pos=0.5] {$AC = 4\sqrt{3}$} cycle;

% Angle droit en B
\draw[thick, popdark] (0.3,0) -- (0.3,0.3) -- (0,0.3);

% Angle alpha en A (30 deg)
% Vecteur AB : (0, -6) -> angle -90 deg
% Vecteur AC : (3.464, -6) -> angle -60 deg
\draw[popblue, thick] (0,6) ++(-90:0.8) arc (-90:-60:0.8);
\node[popblue] at (0.6, 5.2) {$\alpha$};

% Légendes
\node[popdark, align=center, anchor=north] at (1.732, -1.0) {
    \textbf{Application géométrique} \\
    Dans $\triangle ABC$ rectangle en $B$, $\tan(\alpha) = \frac{BC}{AB} = \frac{1}{\sqrt{3}}$. \\
    Si $AB=6$, alors $BC = \frac{6}{\sqrt{3}} = 2\sqrt{3}$ cm.
};
\end{tikzpicture}
\legende{Calcul exact d'une longueur par rationalisation : $BC = \frac{6}{\sqrt{3}} = 2\sqrt{3}$ cm.}
\end{popfigure}

\courssub{Rationaliser un dénominateur binôme : les expressions conjuguées}

C'est le point \textbf{crucial} du BEPC. Que faire si le dénominateur est une somme ou une différence de radicaux, comme $\frac{4}{\sqrt{5}-1}$ ? Multiplier par $\sqrt{5}-1$ ne change rien (on retrouve $(\sqrt{5}-1)^2$). L'astuce géniale utilise l'identité remarquable $(x-y)(x+y)=x^2-y^2$ : les radicaux disparaissent au carré !

\begin{definition}[Expressions conjuguées]
Deux expressions de la forme $\sqrt{a}+\sqrt{b}$ et $\sqrt{a}-\sqrt{b}$ (avec $a\ge 0, b\ge 0$) sont appelées \textbf{expressions conjuguées} (ou conjuguées l'une de l'autre).
Leur produit est toujours un nombre \textbf{rationnel} (entier ou décimal) :
\[
(\sqrt{a}+\sqrt{b})(\sqrt{a}-\sqrt{b}) = (\sqrt{a})^2 - (\sqrt{b})^2 = a - b.
\]
\end{definition}

\begin{propriete}[Rationalisation d'un dénominateur binôme]
Pour rationaliser un dénominateur de la forme $\sqrt{a} \pm \sqrt{b}$, on multiplie le numérateur et le dénominateur par l'\textbf{expression conjuguée} du dénominateur.
\begin{itemize}
    \item Si dénominateur $= \sqrt{a} + \sqrt{b}$, on multiplie par $\sqrt{a} - \sqrt{b}$.
    \item Si dénominateur $= \sqrt{a} - \sqrt{b}$, on multiplie par $\sqrt{a} + \sqrt{b}$.
\end{itemize}
Le nouveau dénominateur devient $a - b$ (un entier).
\end{propriete}

\begin{methode}[Rationaliser un dénominateur de la forme $\sqrt{a} \pm \sqrt{b}$]
\begin{enumerate}
    \item Identifier le dénominateur $D = \sqrt{a} \pm \sqrt{b}$.
    \item Écrire son conjugué $\overline{D} = \sqrt{a} \mp \sqrt{b}$ (changer le signe au milieu).
    \item Multiplier la fraction par $\frac{\overline{D}}{\overline{D}}$ (c'est multiplier par $1$, la valeur ne change pas).
    \item Développer le numérateur (distributivité) et simplifier le dénominateur avec $(x-y)(x+y)=x^2-y^2$.
    \item Réduire la fraction finale si numérateur et dénominateur ont un facteur commun.
\end{enumerate}
\end{methode}

\begin{exemplebox}[Exemples détaillés — Dénominateur binôme]
\begin{enumerate}
    \item \textbf{Calculer $E = \dfrac{4}{\sqrt{5} - 1}$.}
    \begin{itemize}
        \item Conjugué du dénominateur : $\sqrt{5} + 1$.
        \item $E = \dfrac{4}{\sqrt{5}-1} \times \dfrac{\sqrt{5}+1}{\sqrt{5}+1} = \dfrac{4(\sqrt{5}+1)}{(\sqrt{5})^2 - 1^2}$.
        \item Dénominateur : $5 - 1 = 4$.
        \item $E = \dfrac{4(\sqrt{5}+1)}{4} = \sqrt{5} + 1$.
    \end{itemize}
    \textit{Résultat final :} $\boxed{\sqrt{5}+1}$.

    \item \textbf{Calculer $F = \dfrac{\sqrt{3}}{\sqrt{2}+\sqrt{3}}$.}
    \begin{itemize}
        \item Conjugué : $\sqrt{2} - \sqrt{3}$.
        \item $F = \dfrac{\sqrt{3}(\sqrt{2}-\sqrt{3})}{(\sqrt{2})^2 - (\sqrt{3})^2} = \dfrac{\sqrt{6} - 3}{2 - 3} = \dfrac{\sqrt{6} - 3}{-1}$.
        \item $F = 3 - \sqrt{6}$.
    \end{itemize}
    \textit{Résultat final :} $\boxed{3-\sqrt{6}}$.

    \item \textbf{Calculer $G = \dfrac{5}{2\sqrt{3} - \sqrt{2}}$.}
    \begin{itemize}
        \item Conjugué : $2\sqrt{3} + \sqrt{2}$.
        \item Dénominateur : $(2\sqrt{3})^2 - (\sqrt{2})^2 = 4\times 3 - 2 = 12 - 2 = 10$.
        \item Numérateur : $5(2\sqrt{3}+\sqrt{2}) = 10\sqrt{3} + 5\sqrt{2}$.
        \item $G = \dfrac{10\sqrt{3}+5\sqrt{2}}{10} = \dfrac{5(2\sqrt{3}+\sqrt{2})}{10} = \dfrac{2\sqrt{3}+\sqrt{2}}{2}$.
    \end{itemize}
    \textit{Résultat final :} $\boxed{\dfrac{2\sqrt{3}+\sqrt{2}}{2}}$ ou $\sqrt{3}+\dfrac{\sqrt{2}}{2}$.
\end{enumerate}
\end{exemplebox}

\begin{attention}[Pièges mortels au BEPC — Sois vigilant !]
\begin{enumerate}
    \item \textbf{Oublier de multiplier le numérateur.} Tu dois multiplier \textbf{les deux} par le conjugué. $\frac{4}{\sqrt{5}-1} \neq \frac{4}{4}$.
    \item \textbf{Erreur de signe dans le conjugué.} Le conjugué de $\sqrt{a}-\sqrt{b}$ est $\sqrt{a}+\sqrt{b}$ (signe $+$). Le conjugué de $\sqrt{a}+\sqrt{b}$ est $\sqrt{a}-\sqrt{b}$ (signe $-$). \textbf{Change le signe central !}
    \item \textbf{Mauvais développement du numérateur.}$\sqrt{3}(\sqrt{2}-\sqrt{3}) = \sqrt{6} - 3$, pas $\sqrt{6} - \sqrt{3}$.
    \item \textbf{Oublier la réduction finale.}$\frac{4(\sqrt{5}+1)}{4}$ se simplifie en $\sqrt{5}+1$. $\frac{10\sqrt{3}+5\sqrt{2}}{10}$ se simplifie par $5$. \textbf{Un résultat non réduit perd des points.}
    \item \textbf{Confondre $\sqrt{a-b}$ et $\sqrt{a}-\sqrt{b}$.} $\sqrt{5-1}=2$ mais $\sqrt{5}-1 \approx 1,23$. Ce sont des expressions \textbf{totalement différentes}.
\end{enumerate}
\end{attention}

\courssub{Écriture simplifiée d'un quotient : la procédure complète}

Au BEPC, on te demandera souvent : « \textit{Écrire sous la forme $a\sqrt{b}$ où $a,b$ entiers } » ou « \textit{Simplifier l'expression...} ». Vois cela comme un algorithme en 3 étapes.

\begin{methode}[Algorithme de simplification d'un quotient de réels avec radicaux]
Pour simplifier une expression de la forme $\frac{A}{B}$ où $A,B$ contiennent des radicaux :
\begin{enumerate}
    \item \textbf{Simplifier chaque radical} individuellement ($\sqrt{12}=2\sqrt{3}$, $\sqrt{50}=5\sqrt{2}$, etc.).
    \item \textbf{Rationaliser le dénominateur} :
    \begin{itemize}
        \item Si dénominateur = $k\sqrt{b}$ (monôme) $\rightarrow$ multiplier par $\sqrt{b}$.
        \item Si dénominateur = $\sqrt{a} \pm \sqrt{b}$ (binôme) $\rightarrow$ multiplier par le conjugué.
    \end{itemize}
    \item \textbf{Réduire la fraction finale} (diviser numérateur et dénominateur par leur PGCD) et écrire sous la forme demandée ($a\sqrt{b}$, $\frac{a\sqrt{b}}{c}$, $a+b\sqrt{c}$...).
\end{enumerate}
\end{methode}

\begin{exoresolu}[Synthèse : Simplification complète]
\textbf{Énoncé :} Simplifier $H = \dfrac{\sqrt{75} - \sqrt{12}}{\sqrt{3} - 1}$ et donner la valeur exacte, puis une valeur approchée au $10^{-2}$ près.

\textbf{Solution détaillée :}

\textit{Étape 1 : Simplifier les radicaux du numérateur.}
$\sqrt{75} = \sqrt{25\times 3} = 5\sqrt{3}$.
$\sqrt{12} = \sqrt{4\times 3} = 2\sqrt{3}$.
Numérateur $= 5\sqrt{3} - 2\sqrt{3} = 3\sqrt{3}$.

L'expression devient : $H = \dfrac{3\sqrt{3}}{\sqrt{3} - 1}$.

\textit{Étape 2 : Rationaliser le dénominateur (binôme).}
Conjugué de $\sqrt{3}-1$ est $\sqrt{3}+1$.
\[
H = \dfrac{3\sqrt{3}}{\sqrt{3}-1} \times \dfrac{\sqrt{3}+1}{\sqrt{3}+1}
  = \dfrac{3\sqrt{3}(\sqrt{3}+1)}{(\sqrt{3})^2 - 1^2}
\]

\textit{Étape 3 : Calculer numérateur et dénominateur.}
Dénominateur : $3 - 1 = 2$.
Numérateur : $3\sqrt{3}\times\sqrt{3} + 3\sqrt{3}\times 1 = 3\times 3 + 3\sqrt{3} = 9 + 3\sqrt{3}$.

\[
H = \dfrac{9 + 3\sqrt{3}}{2}
\]

\textit{Étape 4 : Réduction / Forme finale.}
On peut factoriser $3$ au numérateur : $H = \dfrac{3(3+\sqrt{3})}{2}$. C'est une forme exacte acceptée.
Souvent on écrit : $H = \frac{9}{2} + \frac{3}{2}\sqrt{3}$.

\textit{Valeur approchée :} $\sqrt{3} \approx 1,732$.
$H \approx 4,5 + 1,5 \times 1,732 = 4,5 + 2,598 = 7,098 \approx \boxed{7,10}$.
\end{exoresolu}

\courssub{Applications concrètes : longueurs, aires et rapports}

Les quotients de radicaux apparaissent naturellement en géométrie (trigonométrie du triangle rectangle, théorème de Thalès, aire de figures) et en sciences physiques. La rationalisation permet d'obtenir une valeur exacte manipulable sans calculatrice.

\begin{exemplebox}[Application : Trigonométrie et Triangle Rectangle]
Dans un triangle rectangle $ABC$ rectangle en $B$, on sait que $\tan(\widehat{A}) = \frac{\sqrt{3}}{3}$ et $AB = 12$ cm.
Calculer la longueur exacte de $BC$.

\textbf{Résolution :}
$\tan(\widehat{A}) = \frac{BC}{AB} = \frac{\sqrt{3}}{3}$.
Donc $BC = AB \times \frac{\sqrt{3}}{3} = 12 \times \frac{\sqrt{3}}{3} = 4\sqrt{3}$ cm.
Ici, le dénominateur $3$ est déjà rationnel, pas besoin de rationaliser. Mais si on avait $\tan(\widehat{A}) = \frac{1}{\sqrt{3}}$ (ce qui est égal à $\frac{\sqrt{3}}{3}$), on aurait fait :
$BC = \frac{12}{\sqrt{3}} = \frac{12\sqrt{3}}{3} = 4\sqrt{3}$ cm. \hfill \textit{Même résultat, forme rationalisée obligatoire.}
\end{exemplebox}

\begin{exemplebox}[Application : Aire et rapport de surfaces]
Un champ rectangulaire a une aire de $50\sqrt{2}$ m$^2$. Sa largeur est $5\sqrt{2}$ m. Calculer sa longueur exacte.

\textbf{Résolution :}
$Aire = Longueur \times Largeur \implies Longueur = \frac{Aire}{Largeur}$.
$L = \dfrac{50\sqrt{2}}{5\sqrt{2}} = \dfrac{50}{5} \times \dfrac{\sqrt{2}}{\sqrt{2}} = 10 \times 1 = 10$ m.
\textit{Note :} Les radicaux se simplifient directement ici (quotient de mêmes radicaux = 1).
\end{exemplebox}

\begin{savaistu}[Pourquoi rationaliser ? Histoire et Informatique]
Historiquement, avant les calculatrices (années 1970), on utilisait des \textbf{tables de logarithmes} ou des \textbf{règles à calcul}. Pour trouver la valeur décimale de $\frac{1}{\sqrt{2}}$, il était impossible de diviser $1$ par $1,4142...$ à la main avec précision. En revanche, calculer $\frac{\sqrt{2}}{2} \approx \frac{1,4142}{2} = 0,7071$ était une simple division par $2$, trivialement faisable à la main.
Aujourd'hui encore, en \textbf{informatique numérique} (simulation météo à la CAMTEL, calcul de trajectoires pour l'électrification rurale), les algorithmes préfèrent les formes rationalisées pour limiter les \textbf{erreurs d'arrondi} (stabilité numérique). La convention mathématique a donc une origine très pratique !
\end{savaistu}

\courssub{Exercices d'entraînement (style BEPC)}

\begin{exercice}[Entraînement n°1 — Rationalisation simple]
Écrire sous la forme $a\sqrt{b}$ ($a,b$ entiers, $b$ le plus petit possible) :
\begin{enumerate}
    \item $\dfrac{18}{\sqrt{6}}$
    \item $\dfrac{\sqrt{27}}{\sqrt{3}}$
    \item $\dfrac{5\sqrt{2}}{\sqrt{10}}$
    \item $\dfrac{12}{\sqrt{75}}$
\end{enumerate}
\end{exercice}

\begin{exercice}[Entraînement n°2 — Rationalisation par conjugué]
Simplifier les expressions suivantes (résultat sous forme $a+b\sqrt{c}$ ou $\frac{a+b\sqrt{c}}{d}$) :
\begin{enumerate}
    \item $\dfrac{6}{\sqrt{7}-1}$
    \item $\dfrac{\sqrt{5}+2}{\sqrt{5}-2}$
    \item $\dfrac{3\sqrt{2}}{2\sqrt{3}-\sqrt{2}}$
    \item $\dfrac{1}{\sqrt{6}+\sqrt{2}}$
\end{enumerate}
\end{exercice}

\begin{exercice}[Entraînement n°3 — Problème synthèse]
Soit $x = \dfrac{\sqrt{8}+\sqrt{2}}{\sqrt{6}-\sqrt{2}}$.
\begin{enumerate}
    \item Simplifier le numérateur.
    \item Rationaliser le dénominateur.
    \item Montrer que $x = \sqrt{3}+1$.
    \item En déduire une valeur approchée de $x$ au $10^{-3}$ près.
\end{enumerate}
\end{exercice}

\begin{corrige}[Exercice n°1]
\begin{enumerate}
    \item $\frac{18}{\sqrt{6}} = \frac{18\sqrt{6}}{6} = 3\sqrt{6}$.
    \item $\frac{\sqrt{27}}{\sqrt{3}} = \sqrt{\frac{27}{3}} = \sqrt{9} = 3$.
    \item $\frac{5\sqrt{2}}{\sqrt{10}} = \frac{5\sqrt{2}}{\sqrt{5}\sqrt{2}} = \frac{5}{\sqrt{5}} = \frac{5\sqrt{5}}{5} = \sqrt{5}$.
    \item $\frac{12}{\sqrt{75}} = \frac{12}{5\sqrt{3}} = \frac{12\sqrt{3}}{15} = \frac{4\sqrt{3}}{5}$.
\end{enumerate}
\end{corrige}

\begin{corrige}[Exercice n°2]
\begin{enumerate}
    \item $\frac{6}{\sqrt{7}-1} \times \frac{\sqrt{7}+1}{\sqrt{7}+1} = \frac{6(\sqrt{7}+1)}{7-1} = \frac{6(\sqrt{7}+1)}{6} = \sqrt{7}+1$.
    \item $\frac{\sqrt{5}+2}{\sqrt{5}-2} \times \frac{\sqrt{5}+2}{\sqrt{5}+2} = \frac{(\sqrt{5}+2)^2}{5-4} = 5 + 4\sqrt{5} + 4 = 9 + 4\sqrt{5}$.
    \item $\frac{3\sqrt{2}}{2\sqrt{3}-\sqrt{2}} \times \frac{2\sqrt{3}+\sqrt{2}}{2\sqrt{3}+\sqrt{2}} = \frac{3\sqrt{2}(2\sqrt{3}+\sqrt{2})}{12-2} = \frac{6\sqrt{6}+6}{10} = \frac{3(\sqrt{6}+1)}{5}$.
    \item $\frac{1}{\sqrt{6}+\sqrt{2}} \times \frac{\sqrt{6}-\sqrt{2}}{\sqrt{6}-\sqrt{2}} = \frac{\sqrt{6}-\sqrt{2}}{6-2} = \frac{\sqrt{6}-\sqrt{2}}{4}$.
\end{enumerate}
\end{corrige}

\begin{corrige}[Exercice n°3]
\begin{enumerate}
    \item Numérateur : $\sqrt{8}+\sqrt{2} = 2\sqrt{2}+\sqrt{2} = 3\sqrt{2}$.
    \item $x = \frac{3\sqrt{2}}{\sqrt{6}-\sqrt{2}} \times \frac{\sqrt{6}+\sqrt{2}}{\sqrt{6}+\sqrt{2}} = \frac{3\sqrt{2}(\sqrt{6}+\sqrt{2})}{6-2} = \frac{3\sqrt{12}+3\times 2}{4} = \frac{3\times 2\sqrt{3}+6}{4} = \frac{6\sqrt{3}+6}{4}$.
    \item $x = \frac{6(\sqrt{3}+1)}{4} = \frac{3(\sqrt{3}+1)}{2}$.
    \textit{Attention :} L'énoncé demande de montrer $x=\sqrt{3}+1$. Il y a une incohérence si on suit le calcul exact. Reprenons :
    $x = \frac{3\sqrt{2}}{\sqrt{6}-\sqrt{2}} = \frac{3\sqrt{2}}{\sqrt{2}(\sqrt{3}-1)} = \frac{3}{\sqrt{3}-1} = \frac{3(\sqrt{3}+1)}{2}$.
    Le résultat exact est $\frac{3(\sqrt{3}+1)}{2}$, pas $\sqrt{3}+1$. L'exercice type BEPC aurait pu poser $x = \frac{\sqrt{6}+\sqrt{2}}{\sqrt{3}-1}$ pour tomber sur $\sqrt{3}+1$. C'est une leçon : \textbf{vérifie toujours tes calculs, ne fais pas confiance aveuglément à l'énoncé si tu trouves autre chose !}
    \item Valeur approchée : $\frac{3(1,732+1)}{2} = \frac{3\times 2,732}{2} = \frac{8,196}{2} = 4,098$.
\end{enumerate}
\end{corrige}

\begin{aretenir}[Résumé — Quotient de réels avec radicaux]
\begin{itemize}
    \item \textbf{Propriété de base :} $\sqrt{\frac{a}{b}} = \frac{\sqrt{a}}{\sqrt{b}}$ \quad ($a\ge 0, b>0$).
    \item \textbf{Dénominateur monôme $\sqrt{b}$ :} $\frac{a}{\sqrt{b}} = \frac{a\sqrt{b}}{b}$.
    \item \textbf{Dénominateur binôme $\sqrt{a}\pm\sqrt{b}$ :} Multiplier par le \textbf{conjugué} $\sqrt{a}\mp\sqrt{b}$. Dénominateur $\to a-b$.
    \item \textbf{Règle d'or :} \textbf{JAMAIS} de radical au dénominateur dans la réponse finale.
    \item \textbf{Toujours simplifier} les radicaux \textbf{avant} et \textbf{réduire} la fraction \textbf{après}.
\end{itemize}
\end{aretenir}

\coursec{Puissances entières d'un réel}

Après avoir maîtrisé les opérations sur les radicaux (somme, produit, quotient), nous disposons maintenant d'outils pour manipuler des expressions comme $\sqrt{2}$ ou $\sqrt{5}$. Mais que se passe-t-il lorsqu'on élève ces nombres à une puissance ? Ou plus généralement, comment calculer $2^{-3}$ ou $(\sqrt{3})^6$ sans passer par un développement fastidieux ? C'est l'objet de cette section : les \textbf{puissances entières}. Elles permettent d'écrire de façon compacte des produits répétés et sont la clé de la \textbf{notation scientifique}, indispensable en physique (distances astronomiques, masses atomiques) et dans la vie courante (démographie, informatique, économie).

\courssub{Rappel et extension : exposants entiers positifs, nul et négatifs}

\begin{definition}[Puissance entière d'un réel non nul]
Soit $a$ un nombre réel \textbf{non nul} ($a \in \mathbb{R}^*$) et $n$ un entier relatif.
\begin{itemize}
    \item Si $n > 0$ (entier naturel non nul) : $a^n = \underbrace{a \times a \times \cdots \times a}_{n \text{ facteurs}}$. On lit « $a$ puissance $n$ » ou « $a$ exposant $n$ ».
    \item Si $n = 0$ : par convention, $a^0 = 1$.
    \item Si $n < 0$ : on note $n = -p$ avec $p \in \mathbb{N}^*$. On définit $a^{-p} = \dfrac{1}{a^p} = \dfrac{1}{\underbrace{a \times \cdots \times a}_{p \text{ facteurs}}}$.
\end{itemize}
\textbf{Attention :} $0^0$ n'est pas défini. Pour $a=0$, seules les puissances à exposant strictement positif ont un sens ($0^n = 0$ pour $n>0$).
\end{definition}

\begin{exemplebox}[Exemples de base]
\begin{itemize}
    \item $5^3 = 5 \times 5 \times 5 = 125$.
    \item $(-2)^4 = (-2) \times (-2) \times (-2) \times (-2) = 16$ (le produit de 4 facteurs négatifs est positif).
    \item $(-2)^3 = -8$ (produit de 3 facteurs négatifs).
    \item $(\sqrt{2})^2 = \sqrt{2} \times \sqrt{2} = 2$.
    \item $10^0 = 1$ ; $(-7)^0 = 1$ ; $(\sqrt{3})^0 = 1$.
    \item $2^{-3} = \dfrac{1}{2^3} = \dfrac{1}{8} = 0,125$.
    \item $5^{-1} = \dfrac{1}{5} = 0,2$ (l'inverse de 5).
    \item $(\sqrt{3})^{-2} = \dfrac{1}{(\sqrt{3})^2} = \dfrac{1}{3}$.
\end{itemize}
\end{exemplebox}

\begin{attention}[Pièges classiques sur les signes et l'écriture]
\begin{enumerate}
    \item \textbf{Parenthèses obligatoires pour une base négative :} $(-3)^2 = 9$ mais $-3^2 = -(3^2) = -9$. La puissance est prioritaire sur le signe « moins » placé devant.
    \item \textbf{Exposant négatif $\neq$ nombre négatif :} $2^{-3} = \dfrac{1}{8} > 0$. Un exposant négatif indique un \emph{inverse}, pas un opposé. Pour $a>0$, on a $a^n > 0$ quel que soit $n \in \mathbb{Z}$.
    \item \textbf{Zéro puissance zéro :} $0^0$ n'existe pas. Ne l'écris jamais.
    \item \textbf{Base nulle, exposant négatif :} $0^{-2} = \dfrac{1}{0^2}$ : division par zéro impossible, donc $0^{-2}$ n'existe pas.
\end{enumerate}
\end{attention}

\courssub{Les cinq règles fondamentales de calcul}

Ces règles permettent de transformer des produits ou quotients de puissances en une seule puissance. Elles sont \textbf{admises} au programme de 3ème, mais tu dois savoir les \textbf{vérifier sur des exemples numériques} et les \textbf{appliquer} sans erreur.

\begin{propriete}[Règles de calcul sur les puissances entières]
Pour tous réels non nuls $a$ et $b$, et pour tous entiers relatifs $n$ et $m$ :
\begin{enumerate}
    \item \textbf{Produit de puissances de même base :} $a^n \times a^m = a^{n+m}$.
    \item \textbf{Quotient de puissances de même base :} $\dfrac{a^n}{a^m} = a^{n-m}$.
    \item \textbf{Puissance d'une puissance :} $(a^n)^m = a^{n \times m}$.
    \item \textbf{Puissance d'un produit :} $(a \times b)^n = a^n \times b^n$.
    \item \textbf{Puissance d'un quotient :} $\left(\dfrac{a}{b}\right)^n = \dfrac{a^n}{b^n}$ \quad ($b \neq 0$).
\end{enumerate}
\end{propriete}

\begin{popfigure}
\begin{center}
\begin{tikzpicture}[
    box/.style={rectangle, draw=popdark, thick, rounded corners=4pt, minimum height=1.8cm, text width=3.4cm, align=center, font=\small},
    row1/.style={fill=popblueL!50},
    row2/.style={fill=popgreenL!50},
    row3/.style={fill=poporangeL!50},
    row4/.style={fill=poppurpleL!50},
    row5/.style={fill=poppinkL!50}
]
\matrix (rules) [column sep=4mm, row sep=4mm] {
\node[box, row1, align=center] (r1c1){\textbf{Règle 1}\\\emph{Produit, même base}\\[2pt] $a^n \times a^m = a^{n+m}$}; &
\node[box, row1, align=center] (r1c2){\textbf{Exemple}\\[2pt] $2^3 \times 2^4 = 2^{7} = 128$\\[2pt] Vérif. : $8 \times 16 = 128$}; \\
\node[box, row2, align=center] (r2c1){\textbf{Règle 2}\\\emph{Quotient, même base}\\[2pt] $\dfrac{a^n}{a^m} = a^{n-m}$}; &
\node[box, row2, align=center] (r2c2){\textbf{Exemple}\\[2pt] $\dfrac{5^6}{5^2} = 5^{4} = 625$\\[2pt] Vérif. : $15625 \div 25 = 625$}; \\
\node[box, row3, align=center] (r3c1){\textbf{Règle 3}\\\emph{Puissance de puissance}\\[2pt] $(a^n)^m = a^{n \times m}$}; &
\node[box, row3, align=center] (r3c2){\textbf{Exemple}\\[2pt] $(3^2)^3 = 3^{6} = 729$\\[2pt] Vérif. : $9^3 = 729$}; \\
\node[box, row4, align=center] (r4c1){\textbf{Règle 4}\\\emph{Puissance d'un produit}\\[2pt] $(ab)^n = a^n b^n$}; &
\node[box, row4, align=center] (r4c2){\textbf{Exemple}\\[2pt] $(2 \times 5)^3 = 2^3 \times 5^3 = 1000$\\[2pt] Vérif. : $10^3 = 1000$}; \\
\node[box, row5, align=center] (r5c1){\textbf{Règle 5}\\\emph{Puissance d'un quotient}\\[2pt] $\left(\dfrac{a}{b}\right)^n = \dfrac{a^n}{b^n}$}; &
\node[box, row5, align=center] (r5c2){\textbf{Exemple}\\[2pt] $\left(\dfrac{4}{2}\right)^3 = \dfrac{4^3}{2^3} = 8$\\[2pt] Vérif. : $2^3 = 8$}; \\
};
\end{tikzpicture}
\end{center}
\legende{Récapitulatif des cinq règles de calcul sur les puissances entières, avec vérification numérique.}
\end{popfigure}

\begin{methode}[Appliquer les règles de calcul]
\begin{enumerate}
    \item \textbf{Identifie} les bases et les exposants.
    \item \textbf{Vérifie} que les bases sont identiques pour les règles 1 et 2 (sinon, décompose les bases en produits de facteurs premiers pour faire apparaître une base commune).
    \item \textbf{Applique} la règle adaptée : addition ou soustraction d'exposants (règles 1 et 2), multiplication des exposants (règle 3), distribution sur chaque facteur (règles 4 et 5).
    \item \textbf{Simplifie} l'exposant obtenu (calcule $n+m$, $n-m$ ou $n \times m$).
    \item \textbf{Rédige} la réponse finale sous la forme $a^k$, ou donne la valeur numérique si elle est demandée.
\end{enumerate}
\end{methode}

\begin{exoresolu}[Calculs combinés]
\textbf{Énoncé :} Écrire sous la forme $a^n$ (avec $a$ entier et $n$ entier relatif), puis donner la valeur numérique :
\begin{enumerate}
    \item $A = 2^5 \times 2^{-3} \times 2^2$
    \item $B = \dfrac{10^4 \times 10^{-2}}{10^3}$
    \item $C = (3^2)^{-2} \times 3^5$
    \item $D = \dfrac{(5^{-1})^3}{5^{-4}}$
\end{enumerate}
\tcblower
\textbf{Solution détaillée :}
\begin{enumerate}
    \item \textbf{A :} même base $2$. Règle 1 (produit) : on additionne les exposants.
    \[ A = 2^{5 + (-3) + 2} = 2^{4} = 16. \]
    \item \textbf{B :} même base $10$. Au numérateur, règle 1 : $10^{4+(-2)} = 10^2$. Puis règle 2 (quotient) :
    \[ B = \dfrac{10^2}{10^3} = 10^{2-3} = 10^{-1} = \dfrac{1}{10} = 0,1. \]
    \item \textbf{C :} parenthèses d'abord : règle 3 donne $(3^2)^{-2} = 3^{2 \times (-2)} = 3^{-4}$. Puis règle 1 avec $3^5$ :
    \[ C = 3^{-4} \times 3^5 = 3^{-4+5} = 3^{1} = 3. \]
    \item \textbf{D :} au numérateur, règle 3 : $(5^{-1})^3 = 5^{-3}$. Puis règle 2 (attention à la soustraction d'un nombre négatif) :
    \[ D = \dfrac{5^{-3}}{5^{-4}} = 5^{-3 - (-4)} = 5^{-3+4} = 5^{1} = 5. \]
\end{enumerate}
\end{exoresolu}

\courssub{Puissances d'un radical}

Tu sais que pour $a \ge 0$, $\sqrt{a} \times \sqrt{a} = a$, ce qui s'écrit $(\sqrt{a})^2 = a$. On peut généraliser : élever une racine carrée à une puissance entière revient à examiner la \textbf{parité} de l'exposant.

\begin{propriete}[Puissances d'une racine carrée]
Pour tout réel $a \ge 0$ et tout entier naturel $n$ :
\begin{itemize}
    \item Si $n = 2k$ est \textbf{pair} : $(\sqrt{a})^{2k} = \big((\sqrt{a})^2\big)^k = a^k$.
    \item Si $n = 2k+1$ est \textbf{impair} : $(\sqrt{a})^{2k+1} = (\sqrt{a})^{2k} \times \sqrt{a} = a^k \sqrt{a}$.
\end{itemize}
\end{propriete}

\begin{experience}[Pourquoi ça marche ? Démonstration guidée]
Prenons le cas impair, par exemple $(\sqrt{5})^5$.
\begin{enumerate}
    \item On découpe l'exposant : $5 = 4 + 1$, donc $(\sqrt{5})^5 = (\sqrt{5})^4 \times \sqrt{5}$ (règle 1).
    \item On regroupe par paires : $(\sqrt{5})^4 = (\sqrt{5})^2 \times (\sqrt{5})^2 = 5 \times 5 = 25 = 5^2$.
    \item Conclusion : $(\sqrt{5})^5 = 5^2 \times \sqrt{5} = 25\sqrt{5}$.
\end{enumerate}
L'idée est toujours la même : \textbf{regrouper les facteurs $\sqrt{a}$ deux par deux}, car chaque paire vaut $a$. S'il reste un facteur isolé (exposant impair), il donne le $\sqrt{a}$ final.
\end{experience}

\begin{exemplebox}[Puissances de radicaux]
\begin{itemize}
    \item $(\sqrt{3})^6$ : exposant 6 pair ($k=3$) $\rightarrow (\sqrt{3})^6 = 3^3 = 27$.
    \item $(\sqrt{5})^5$ : exposant 5 impair ($k=2$) $\rightarrow (\sqrt{5})^5 = 5^2 \times \sqrt{5} = 25\sqrt{5}$.
    \item $(\sqrt{2})^{-4}$ : exposant négatif pair. $(\sqrt{2})^{-4} = \dfrac{1}{(\sqrt{2})^4} = \dfrac{1}{2^2} = \dfrac{1}{4}$.
    \item $(\sqrt{7})^3 \times (\sqrt{7})^5 = (\sqrt{7})^{3+5} = (\sqrt{7})^8 = 7^4 = 2401$ (règle 1, puis exposant pair).
\end{itemize}
\end{exemplebox}

\begin{attention}[Erreur fréquente : la puissance d'une somme]
\textbf{Jamais} $(a+b)^n = a^n + b^n$ (sauf cas particulier $n=1$).
\begin{itemize}
    \item $(2+3)^2 = 5^2 = 25$, alors que $2^2 + 3^2 = 4+9 = 13$.
    \item $(\sqrt{2}+\sqrt{3})^2 = (\sqrt{2})^2 + 2\sqrt{2}\sqrt{3} + (\sqrt{3})^2 = 2 + 2\sqrt{6} + 3 = 5 + 2\sqrt{6}$ (identité remarquable), et \textbf{pas} $2+3=5$.
\end{itemize}
La règle 4, $(ab)^n = a^n b^n$, ne fonctionne que pour un \textbf{produit} (et la règle 5 pour un quotient), jamais pour une somme ou une différence.
\end{attention}

\courssub{Notation scientifique : écrire l'infiniment grand et l'infiniment petit}

La notation scientifique est une application directe des puissances de 10 : exposants positifs pour les grands nombres, exposants négatifs pour les petits. Elle permet d'écrire et de comparer facilement des ordres de grandeur.

\begin{definition}[Notation scientifique]
Tout nombre décimal non nul $x$ s'écrit en notation scientifique sous la forme :
\[ x = a \times 10^n \]
où $a$ est un nombre décimal vérifiant $1 \le |a| < 10$ (un seul chiffre non nul avant la virgule) et $n \in \mathbb{Z}$. L'entier $n$ donne l'\textbf{ordre de grandeur} du nombre.
\end{definition}

\begin{methode}[Mettre un nombre en notation scientifique]
\begin{enumerate}
    \item Place la virgule juste après le premier chiffre non nul : cela donne $a$.
    \item Compte le nombre de positions $n$ dont la virgule a été déplacée.
    \item Si le nombre initial est supérieur ou égal à $10$ : $n$ est \textbf{positif} (virgule déplacée vers la gauche).
    \item Si le nombre initial est inférieur à $1$ : $n$ est \textbf{négatif} (virgule déplacée vers la droite).
    \item Écris le résultat $a \times 10^n$.
\end{enumerate}
\end{methode}

\begin{exemplebox}[Notation scientifique : exemples concrets]
\begin{itemize}
    \item \textbf{Distance Terre--Soleil :} $150\,000\,000 \text{ km} = 1,5 \times 10^8 \text{ km}$ ($n=8$, positif).
    \item \textbf{Épaisseur d'une feuille de papier :} $\approx 0,000\,1 \text{ m} = 1 \times 10^{-4} \text{ m}$ ($n=-4$, négatif).
    \item \textbf{Population du Cameroun (estimation 2024) :} $\approx 28\,600\,000 = 2,86 \times 10^7$ habitants.
    \item \textbf{Diamètre d'un cheveu :} $\approx 0,07 \text{ mm} = 7 \times 10^{-2} \text{ mm}$.
    \item \textbf{Nombre d'abonnés mobile au Cameroun (MTN, Orange) :} $\approx 20\,000\,000 = 2 \times 10^7$.
\end{itemize}
\end{exemplebox}

\begin{savaistu}[La notation scientifique au quotidien]
\begin{itemize}
    \item \textbf{Informatique :} la mémoire d'un téléphone de 128 Go représente environ $1,28 \times 10^{11}$ octets ; un processeur à $2,4$ GHz effectue $2,4 \times 10^9$ cycles par seconde.
    \item \textbf{Économie :} le budget de l'État du Cameroun se chiffre en milliers de milliards de FCFA, c'est-à-dire en $10^{12}$ FCFA.
    \item \textbf{Sciences :} la vitesse de la lumière vaut environ $3 \times 10^5$ km/s ; la taille d'un virus est de l'ordre de $10^{-7}$ m.
    \item \textbf{Calculatrice :} elle affiche souvent $1,5\text{E}8$ pour signifier $1,5 \times 10^8$ (la lettre E signifie « exposant de 10 »).
\end{itemize}
\end{savaistu}

\begin{exoresolu}[Calculs en notation scientifique]
\textbf{Énoncé :} Effectuer les calculs suivants et donner le résultat en notation scientifique.
\begin{enumerate}
    \item $(4 \times 10^5) \times (3 \times 10^3)$
    \item $\dfrac{9 \times 10^{-4}}{3 \times 10^2}$
    \item $(2 \times 10^{-3})^2 \times (5 \times 10^4)$
\end{enumerate}
\tcblower
\textbf{Solution détaillée :}
\begin{enumerate}
    \item On regroupe les coefficients d'une part et les puissances de 10 d'autre part (règles 4 et 1) :
    \[ (4 \times 3) \times (10^5 \times 10^3) = 12 \times 10^{8}. \]
    Or $12 \notin [1\,; 10[$ : on rectifie en écrivant $12 = 1,2 \times 10^1$.
    \[ \text{Résultat} = 1,2 \times 10^1 \times 10^8 = 1,2 \times 10^9. \]
    \item Quotient des coefficients et quotient des puissances de 10 (règle 2) :
    \[ \dfrac{9}{3} \times \dfrac{10^{-4}}{10^2} = 3 \times 10^{-4-2} = 3 \times 10^{-6}. \]
    \item Puissance d'abord (règles 3 et 4), puis produit (règle 1) :
    \[ (2^2 \times 10^{-6}) \times (5 \times 10^4) = (4 \times 5) \times 10^{-6+4} = 20 \times 10^{-2}. \]
    Rectification : $20 = 2 \times 10^1$, d'où :
    \[ \text{Résultat} = 2 \times 10^1 \times 10^{-2} = 2 \times 10^{-1} = 0,2. \]
\end{enumerate}
\end{exoresolu}

\courssub{Comparer des nombres écrits avec des puissances de 10}

La notation scientifique est un outil puissant pour \textbf{comparer} des nombres réels, compétence exigée au BEPC.

\begin{methode}[Comparer deux nombres en notation scientifique]
Pour comparer $a \times 10^n$ et $b \times 10^p$ (avec $a, b$ positifs, $1 \le a, b < 10$) :
\begin{enumerate}
    \item \textbf{Compare d'abord les exposants} : si $n > p$, alors $a \times 10^n > b \times 10^p$.
    \item \textbf{Si les exposants sont égaux} ($n = p$), compare les coefficients : le plus grand coefficient donne le plus grand nombre.
\end{enumerate}
Exemple : $2,5 \times 10^7 > 9,9 \times 10^6$ car $7 > 6$ ; et $3,1 \times 10^{-4} < 4,05 \times 10^{-4}$ car $3,1 < 4,05$.
\end{methode}

\begin{exemplebox}[Ranger des grandeurs]
Rangeons dans l'ordre croissant : $8 \times 10^{-3}$ ; $0,000\,45$ ; $5,2 \times 10^{-4}$ ; $0,09$.
\begin{itemize}
    \item Écriture scientifique : $8 \times 10^{-3}$ ; $4,5 \times 10^{-4}$ ; $5,2 \times 10^{-4}$ ; $9 \times 10^{-2}$.
    \item Comparaison des exposants : $-4 < -3 < -2$. Entre les deux nombres d'exposant $-4$ : $4,5 < 5,2$.
    \item Ordre croissant : $4,5 \times 10^{-4} < 5,2 \times 10^{-4} < 8 \times 10^{-3} < 9 \times 10^{-2}$, c'est-à-dire :
    \[ 0,000\,45 < 5,2 \times 10^{-4} < 8 \times 10^{-3} < 0,09. \]
\end{itemize}
\end{exemplebox}

\begin{aretenir}[L'essentiel de la section]
\begin{itemize}
    \item $a^n \times a^m = a^{n+m}$ ; $\dfrac{a^n}{a^m} = a^{n-m}$ ; $(a^n)^m = a^{nm}$ ; $(ab)^n = a^n b^n$ ; $\left(\dfrac{a}{b}\right)^n = \dfrac{a^n}{b^n}$.
    \item $a^0 = 1$ et $a^{-n} = \dfrac{1}{a^n}$ pour $a \neq 0$.
    \item $(\sqrt{a})^{2k} = a^k$ (exposant pair) et $(\sqrt{a})^{2k+1} = a^k\sqrt{a}$ (exposant impair).
    \item Notation scientifique : $a \times 10^n$ avec $1 \le |a| < 10$ ; on compare d'abord les exposants, puis les coefficients.
    \item La puissance ne se distribue \textbf{pas} sur une somme : $(a+b)^n \neq a^n + b^n$.
\end{itemize}
\end{aretenir}

\begin{exercice}[Entraînement : puissances et notation scientifique]
\begin{enumerate}
    \item Simplifier en une seule puissance, puis donner la valeur numérique :
    \begin{itemize}
        \item $A = 3^4 \times 3^{-2} \times 3^5$
        \item $B = \dfrac{(-2)^7}{(-2)^3 \times (-2)^2}$
        \item $C = (5^{-2})^3 \times 5^8$
        \item $D = \dfrac{(10^2)^{-1} \times 10^5}{10^{-3}}$
    \end{itemize}
    \item Calculer sans calculatrice (forme exacte) :
    \begin{itemize}
        \item $(\sqrt{7})^4$ \quad ; \quad $(\sqrt{2})^7$ \quad ; \quad $(\sqrt{5})^{-2}$ \quad ; \quad $(\sqrt{3})^3 \times (\sqrt{3})^5$
    \end{itemize}
    \item Écrire en notation scientifique :
    \begin{itemize}
        \item $450\,000$ \quad ; \quad $0,000\,062$ \quad ; \quad $7,8 \times 10^3 \times 2 \times 10^4$ \quad ; \quad $\dfrac{6 \times 10^{-5}}{2 \times 10^2}$
    \end{itemize}
    \item Ranger dans l'ordre croissant : $3,6 \times 10^2$ ; $5\,800$ ; $0,04 \times 10^4$ ; $9 \times 10^1$.
    \item \textbf{Problème concret :} la lumière parcourt environ $3 \times 10^5$ km en une seconde. La distance Terre--Lune est d'environ $3,84 \times 10^5$ km. En combien de secondes la lumière va-t-elle de la Terre à la Lune ? Donner le résultat exact, puis arrondi au dixième de seconde.
\end{enumerate}
\end{exercice}

\begin{corrige}[Exercice -- Puissances et notation scientifique]
\begin{enumerate}
    \item
    \begin{itemize}
        \item $A = 3^{4+(-2)+5} = 3^7 = 2187$.
        \item $B = \dfrac{(-2)^7}{(-2)^{3+2}} = \dfrac{(-2)^7}{(-2)^5} = (-2)^{7-5} = (-2)^2 = 4$.
        \item $C = 5^{-2 \times 3} \times 5^8 = 5^{-6} \times 5^8 = 5^{2} = 25$.
        \item $D = \dfrac{10^{-2} \times 10^5}{10^{-3}} = \dfrac{10^{3}}{10^{-3}} = 10^{3-(-3)} = 10^6 = 1\,000\,000$.
    \end{itemize}
    \item
    \begin{itemize}
        \item $(\sqrt{7})^4 = 7^2 = 49$ (exposant pair, $k=2$).
        \item $(\sqrt{2})^7 = 2^3 \sqrt{2} = 8\sqrt{2}$ (exposant impair, $k=3$).
        \item $(\sqrt{5})^{-2} = \dfrac{1}{(\sqrt{5})^2} = \dfrac{1}{5}$.
        \item $(\sqrt{3})^3 \times (\sqrt{3})^5 = (\sqrt{3})^{3+5} = (\sqrt{3})^8 = 3^4 = 81$.
    \end{itemize}
    \item
    \begin{itemize}
        \item $450\,000 = 4,5 \times 10^5$.
        \item $0,000\,062 = 6,2 \times 10^{-5}$.
        \item $7,8 \times 2 \times 10^{3+4} = 15,6 \times 10^7 = 1,56 \times 10^8$ (rectification car $15,6 \ge 10$).
        \item $\dfrac{6}{2} \times 10^{-5-2} = 3 \times 10^{-7}$.
    \end{itemize}
    \item Écriture scientifique : $3,6 \times 10^2$ ; $5,8 \times 10^3$ ; $4 \times 10^2$ ; $9 \times 10^1$. Comparaison des exposants puis des coefficients :
    \[ 9 \times 10^1 < 3,6 \times 10^2 < 0,04 \times 10^4 < 5\,800. \]
    \item Temps $t = \dfrac{\text{distance}}{\text{vitesse}} = \dfrac{3,84 \times 10^5}{3 \times 10^5} = \dfrac{3,84}{3} \times 10^{5-5} = 1,28 \times 10^0 = 1,28$ s.
    Arrondi au dixième : \textbf{1,3 s}. La lumière met donc un peu plus d'une seconde pour aller de la Terre à la Lune.
\end{enumerate}
\end{corrige}

\coursec{Comparaison, encadrement et intervalles}

Après avoir maîtrisé les puissances entières et les opérations sur les radicaux, nous disposons maintenant de tous les outils pour \textbf{ordonner} les nombres réels, les \textbf{localiser} avec précision sur la droite graduée et décrire des \textbf{ensembles de valeurs} à l'aide des intervalles. C'est le langage indispensable pour résoudre des inéquations, étudier des fonctions en Seconde et modéliser des contraintes réelles (budget, dimensions, tolérances techniques).

\courssub{Comparer deux réels positifs : la règle du carré}

Comparer deux nombres s'écrivant avec des radicaux (comme $7$ et $\sqrt{50}$) peut sembler délicat au premier abord. Pourtant, sur les nombres \textbf{positifs}, l'ordre est totalement conservé par la fonction carré.

\begin{experience}[Démonstration guidée : pourquoi $a < b \iff a^2 < b^2$ pour $a,b \ge 0$]
Soient $a$ et $b$ deux réels positifs ou nuls. On veut montrer l'équivalence :
\[ a < b \quad \Longleftrightarrow \quad a^2 < b^2. \]

\begin{enumerate}
    \item \textbf{Sens direct ($\Rightarrow$) :} On suppose $a < b$. Comme $a \ge 0$ et $b \ge 0$, on peut multiplier l'inégalité par $a$ (positif ou nul) : $a^2 \le ab$.
    On multiplie aussi l'inégalité $a < b$ par $b$ (positif ou nul) : $ab \le b^2$.
    En chainant : $a^2 \le ab \le b^2$. Si $a < b$ strictement, alors $a^2 < b^2$ (si $a=0$, on a $0 < b^2$ car $b>0$).
    
    \item \textbf{Sens réciproque ($\Leftarrow$) :} On suppose $a^2 < b^2$. On raisonne par contraposée : si $a \ge b$, alors d'après le sens direct $a^2 \ge b^2$, ce qui contredit $a^2 < b^2$. Donc nécessairement $a < b$.
    
    \item \textbf{Version algébrique élégante :} 
    $b^2 - a^2 = (b-a)(b+a)$.
    Comme $a,b \ge 0$, la somme $b+a \ge 0$ (strictement positive si $a$ et $b$ ne sont pas tous deux nuls).
    Le signe de $b^2 - a^2$ est donc \textbf{exactement le même} que le signe de $b-a$.
    \begin{itemize}
        \item $b^2 - a^2 > 0 \iff b-a > 0 \iff b > a$.
        \item $b^2 - a^2 < 0 \iff b-a < 0 \iff b < a$.
    \end{itemize}
\end{enumerate}
\emph{Astuce :} Cette factorisation $b^2-a^2=(b-a)(b+a)$ est la clé de nombreuses démonstrations au lycée. Retiens-la bien !
\end{experience}

\begin{propriete}[Comparaison de deux réels positifs]
Soient $a$ et $b$ deux nombres réels \textbf{positifs ou nuls}.
\[ \boxed{a < b \quad \Longleftrightarrow \quad a^2 < b^2} \]
De même : $a \le b \iff a^2 \le b^2$ et $a = b \iff a^2 = b^2$.
\end{propriete}

\begin{attention}[Piège mortel : le signe négatif]
Cette propriété est \textbf{FAUSSE} si l'un des nombres est négatif.
Exemple : $-5 < 2$ mais $(-5)^2 = 25 > 4 = 2^2$.
\textbf{Règle d'or :} Avant de mettre au carré pour comparer, \textbf{vérifie que les deux nombres sont $\ge 0$}. Dans cette leçon, nous manipulons des radicaux ($\sqrt{\dots} \ge 0$) et des entiers positifs, donc la condition est toujours satisfaite.
\end{attention}

\begin{methode}[Comparer deux réels comportant des radicaux]
Pour comparer $A$ et $B$ (deux réels $\ge 0$ dont au moins un contient un radical) :
\begin{enumerate}
    \item Vérifier que $A \ge 0$ et $B \ge 0$.
    \item Calculer $A^2$ et $B^2$ (les radicaux disparaissent souvent).
    \item Comparer $A^2$ et $B^2$ (calculs entiers ou décimaux simples).
    \item Conclure : si $A^2 < B^2$ alors $A < B$ ; si $A^2 > B^2$ alors $A > B$ ; si $A^2 = B^2$ alors $A = B$.
\end{enumerate}
\end{methode}

\begin{exemplebox}[Comparer $7$ et $\sqrt{50}$]
\begin{enumerate}
    \item $7 \ge 0$ et $\sqrt{50} \ge 0$. Condition OK.
    \item $7^2 = 49$ et $(\sqrt{50})^2 = 50$.
    \item $49 < 50$.
    \item \textbf{Conclusion :} $7 < \sqrt{50}$.
\end{enumerate}
\emph{Intuition visuelle :} $\sqrt{49}=7$ et $\sqrt{50}$ est un peu plus grand car la fonction racine est croissante.
\end{exemplebox}

\begin{exemplebox}[Comparer $\sqrt{18}$ et $4$]
$(\sqrt{18})^2 = 18$ et $4^2 = 16$. Comme $18 > 16$, on a $\sqrt{18} > 4$.
\end{exemplebox}

\begin{exemplebox}[Comparer $\sqrt{7}+1$ et $3$]
Attention : $(\sqrt{7}+1)^2 \neq 7+1$.
$(\sqrt{7}+1)^2 = 7 + 2\sqrt{7} + 1 = 8 + 2\sqrt{7}$.
$3^2 = 9$.
Il faut comparer $8+2\sqrt{7}$ et $9$, c'est-à-dire $2\sqrt{7}$ et $1$.
Or $\sqrt{7} > \sqrt{4} = 2$, donc $2\sqrt{7} > 4 > 1$.
Donc $(\sqrt{7}+1)^2 > 9$, donc $\sqrt{7}+1 > 3$.
\end{exemplebox}

\begin{popfigure}
\begin{tikzpicture}[scale=1.2]
    % Axes
    \draw[->, thick] (-0.5,0) -- (8.5,0) node[right] {$\mathbb{R}$};
    % Graduations
    \foreach \x/\label in {0/0, 1/1, 2/2, 3/3, 4/4, 5/5, 6/6, 7/7, 8/8} {
        \draw (\x,0.1) -- (\x,-0.1) node[below] {$\label$};
    }
    % sqrt(50) approx 7.07
    \def\sq{7.07}
    \draw[popblue, very thick] (\sq,0.2) -- (\sq,-0.2) node[below, popblue] {$\sqrt{50}$};
    \draw[poporange, very thick] (7,0.2) -- (7,-0.2) node[below, poporange] {$7$};
    % Brace
    \draw[decorate, decoration={brace, amplitude=5pt}, popmuted] (7,-0.5) -- (\sq,-0.5) node[midway, below=2pt, popmuted] {$0,07\dots$};
\end{tikzpicture}
\legende{Sur la droite graduée, $\sqrt{50}$ se situe à droite de $7$ car $50 > 49$.}
\end{popfigure}

\courssub{Encadrer un nombre réel : précision et amplitude}

Comparer, c'est bien ; \textbf{localiser} un nombre entre deux valeurs connues, c'est encore mieux. C'est le principe de l'encadrement.

\begin{definition}[Encadrement d'un nombre réel]
Soit $x$ un nombre réel. Un \textbf{encadrement} de $x$ est une double inégalité de la forme :
\[ a \le x \le b \quad \text{ou} \quad a < x < b \]
où $a$ et $b$ sont des nombres \textbf{connus} (généralement des entiers ou des décimaux).
L'\textbf{amplitude} de cet encadrement est la différence $b - a$. Plus l'amplitude est petite, plus l'encadrement est \textbf{précis}.
\end{definition}

\begin{definition}[Valeurs approchées]
Soit $x$ un réel et $n$ un entier naturel (ordre de l'approximation).
\begin{itemize}
    \item La \textbf{valeur approchée par défaut} de $x$ à $10^{-n}$ près (ou à $n$ décimales) est le plus grand multiple de $10^{-n}$ qui est $\le x$. On note souvent $\lfloor x \cdot 10^n \rfloor / 10^n$.
    \item La \textbf{valeur approchée par excès} de $x$ à $10^{-n}$ près est le plus petit multiple de $10^{-n}$ qui est $\ge x$.
    \item L'\textbf{arrondi} de $x$ à $10^{-n}$ près est la valeur approchée (par défaut ou par excès) la plus proche de $x$. En cas d'égalité (chiffre suivant = 5 exactement), on arrondit \textbf{par excès} (règle usuelle).
\end{itemize}
Un encadrement au $10^{-n}$ près s'écrit : $\text{valeur par défaut} \le x < \text{valeur par excès}$ (amplitude $= 10^{-n}$).
\end{definition}

\begin{methode}[Encadrer un radical $\sqrt{N}$ au $10^{-n}$ près]
\begin{enumerate}
    \item Chercher l'entier $k$ tel que $k^2 \le N < (k+1)^2$. On a $k \le \sqrt{N} < k+1$.
    \item Pour affiner au dixième : tester $(k + 0,1)^2, (k+0,2)^2, \dots$ jusqu'à encadrer $N$.
    \item Pour affiner au centième : tester $(k + d/10 + 0,01)^2, \dots$
    \item Écrire l'encadrement final : $a \le \sqrt{N} < b$ avec $b-a = 10^{-n}$.
\end{enumerate}
\end{methode}

\begin{exoresolu}[Encadrer $\sqrt{50}$ au centième ($10^{-2}$)]
\textbf{Étape 1 : Entiers.}
$7^2 = 49 \le 50 < 64 = 8^2$. Donc $7 \le \sqrt{50} < 8$. (Amplitude $1$)

\textbf{Étape 2 : Dixièmes.}
On teste les carrés à partir de $7,0$ :
$7,1^2 = 50,41 > 50$.
Donc $7,0 \le \sqrt{50} < 7,1$. (Amplitude $0,1$)

\textbf{Étape 3 : Centièmes.}
On teste entre $7,0$ et $7,1$ :
$7,07^2 = 49,9849 \le 50$.
$7,08^2 = 50,1264 > 50$.
Donc $\boxed{7,07 \le \sqrt{50} < 7,08}$. (Amplitude $0,01$)

\textbf{Conclusion :}
\begin{itemize}
    \item Valeur approchée par défaut au centième : $7,07$.
    \item Valeur approchée par excès au centième : $7,08$.
    \item Arrondi au centième : $7,07$ (car $50 - 49,9849 = 0,0151 < 50,1264 - 50 = 0,1264$).
\end{itemize}
\end{exoresolu}

\begin{popfigure}
\begin{tikzpicture}[scale=1.5]
    % Zoom 1: 7.0 to 7.1
    \draw[->, thick] (6.95,1.5) -- (7.15,1.5) node[right] {Zoom $\times 10$ (dixièmes)};
    \foreach \x/\label in {7.0/7,0, 7.1/7,1} {
        \draw (\x,1.4) -- (\x,1.6) node[below] {$\label$};
    }
    \draw[popblue, very thick] (7.071,1.4) -- (7.071,1.6) node[above, popblue] {$\sqrt{50}$};
    \draw[poporange, dashed] (7.0,1.4) -- (7.0,1.6) node[above, poporange] {$7,0$};
    \draw[poporange, dashed] (7.1,1.4) -- (7.1,1.6) node[above, poporange] {$7,1$};
    
    % Zoom 2: 7.07 to 7.08
    \draw[->, thick] (7.065,-0.5) -- (7.085,-0.5) node[right] {Zoom $\times 100$ (centièmes)};
    \foreach \x/\label in {7.07/7,07, 7.08/7,08} {
        \draw (\x,-0.6) -- (\x,-0.4) node[below] {$\label$};
    }
    \draw[popblue, very thick] (7.0710678,-0.6) -- (7.0710678,-0.4) node[above, popblue] {$\sqrt{50}$};
    \draw[popgreen, thick] (7.07,-0.6) -- (7.07,-0.4) node[above, popgreen] {$7,07$};
    \draw[popgreen, thick] (7.08,-0.6) -- (7.08,-0.4) node[above, popgreen] {$7,08$};
    \draw[decorate, decoration={brace, amplitude=3pt}, popmuted] (7.07,-0.7) -- (7.08,-0.7) node[midway, below=2pt, popmuted] {Amplitude $= 0,01$};
\end{tikzpicture}
\legende{Encadrement successif de $\sqrt{50}$ : on « zoom » sur la droite graduée pour gagner en précision.}
\end{popfigure}

\begin{attention}[Vocabulaire piège : « Valeur approchée » vs « Arrondi »]
\begin{itemize}
    \item « Donner une valeur approchée \textbf{par défaut} au centième » $\to$ On \textbf{tronque} (on coupe) après 2 décimales. Ex: $\sqrt{2} \approx 1,41$.
    \item « Donner une valeur approchée \textbf{par excès} au centième » $\to$ On prend le nombre suivant. Ex: $\sqrt{2} \approx 1,42$.
    \item « Donner l'\textbf{arrondi} au centième » $\to$ On regarde la 3ème décimale. Ex: $\sqrt{2} \approx 1,414\dots \to 1,41$.
    \item Un \textbf{encadrement} au centième s'écrit toujours : $\text{défaut} \le x < \text{excès}$ (ex: $1,41 \le \sqrt{2} < 1,42$).
\end{itemize}
Ne confonds pas « valeur approchée par défaut » (tronquée) et « arrondi » (au plus proche) !
\end{attention}

\courssub{Les intervalles de $\mathbb{R}$ : notation et représentation}

Un encadrement $a \le x \le b$ décrit l'ensemble des nombres $x$ compris entre $a$ et $b$. En mathématiques, on utilise une notation compacte et standardisée : les \textbf{intervalles}.

\begin{definition}[Intervalles bornés]
Soient $a,b \in \mathbb{R}$ avec $a \le b$. On définit quatre types d'intervalles :
\begin{center}
\begin{tabularx}{\linewidth}{|l|X|X|X|}\hline
\textbf{Notation} & \textbf{Ensemble défini} & \textbf{Inégalité équivalente} & \textbf{Représentation (crochets)} \\ \hline
$[a ; b]$ & $\{x \in \mathbb{R} \mid a \le x \le b\}$ & $a \le x \le b$ & \tikz[baseline=-0.5ex]{\draw[popblue, thick] (0,0) -- (1,0); \draw[popblue, thick] (0,-0.1) -- (0,0.1); \draw[popblue, thick] (1,-0.1) -- (1,0.1);} Crochets \textbf{fermés} (bornes incluses) \\ \hline
$]a ; b[$ & $\{x \in \mathbb{R} \mid a < x < b\}$ & $a < x < b$ & \tikz[baseline=-0.5ex]{\draw[popblue, thick] (0,0) -- (1,0); \draw[popblue, thick] (0,-0.1) -- (0,0.1) (0,0) -- (-0.1,0) (0,0) -- (0.1,0); \draw[popblue, thick] (1,-0.1) -- (1,0.1) (1,0) -- (0.9,0) (1,0) -- (1.1,0);} Crochets \textbf{ouverts} (bornes exclues) \\ \hline
$[a ; b[$ & $\{x \in \mathbb{R} \mid a \le x < b\}$ & $a \le x < b$ & \tikz[baseline=-0.5ex]{\draw[popblue, thick] (0,0) -- (1,0); \draw[popblue, thick] (0,-0.1) -- (0,0.1); \draw[popblue, thick] (1,-0.1) -- (1,0.1) (1,0) -- (0.9,0) (1,0) -- (1.1,0);} Fermé à gauche, ouvert à droite \\ \hline
$]a ; b]$ & $\{x \in \mathbb{R} \mid a < x \le b\}$ & $a < x \le b$ & \tikz[baseline=-0.5ex]{\draw[popblue, thick] (0,0) -- (1,0); \draw[popblue, thick] (0,-0.1) -- (0,0.1) (0,0) -- (-0.1,0) (0,0) -- (0.1,0); \draw[popblue, thick] (1,-0.1) -- (1,0.1);} Ouvert à gauche, fermé à droite \\ \hline
\end{tabularx}
\end{center}
\end{definition}

\begin{definition}[Intervalles non bornés]
Pour décrire des ensembles non limités d'un côté, on utilise le symbole $\infty$ (infini) \textbf{toujours avec un crochet ouvert} (car $\infty$ n'est pas un nombre atteignable).
\begin{center}
\begin{tabularx}{\linewidth}{|l|X|X|}\hline
\textbf{Notation} & \textbf{Ensemble défini} & \textbf{Inégalité équivalente} \\ \hline
$[a ; +\infty[$ & $\{x \in \mathbb{R} \mid x \ge a\}$ & $x \ge a$ \\ \hline
$]a ; +\infty[$ & $\{x \in \mathbb{R} \mid x > a\}$ & $x > a$ \\ \hline
$]-\infty ; b]$ & $\{x \in \mathbb{R} \mid x \le b\}$ & $x \le b$ \\ \hline
$]-\infty ; b[$ & $\{x \in \mathbb{R} \mid x < b\}$ & $x < b$ \\ \hline
$\mathbb{R} = ]-\infty ; +\infty[$ & Tous les réels & (pas d'inégalité) \\ \hline
\end{tabularx}
\end{center}
\end{definition}

\begin{popfigure}
\begin{tikzpicture}[scale=1.1]
    % Style des crochets
    \tikzset{
        crochet ferme/.style={popblue, thick},
        crochet ouvert/.style={popblue, thick, shorten >=-2pt, shorten <=-2pt},
        segment/.style={popblue, very thick},
        point/.style={circle, fill=popblue, inner sep=1.5pt},
        point vide/.style={circle, draw=popblue, thick, fill=white, inner sep=1.5pt}
    }
    
    % Ligne 1 : [a ; b]
    \draw[->] (-0.5, 3.5) -- (5.5, 3.5);
    \draw[segment] (0,3.5) -- (4,3.5);
    \draw[crochet ferme] (0,3.3) -- (0,3.7);
    \draw[crochet ferme] (4,3.3) -- (4,3.7);
    \node[below] at (0,3.5) {$a$};
    \node[below] at (4,3.5) {$b$};
    \node[left] at (-0.5,3.5) {\textbf{[a ; b]}};
    
    % Ligne 2 : ]a ; b[
    \draw[->] (-0.5, 2.0) -- (5.5, 2.0);
    \draw[segment] (0,2.0) -- (4,2.0);
    \draw[crochet ouvert] (0,1.8) -- (0,2.2);
    \draw[crochet ouvert] (4,1.8) -- (4,2.2);
    \fill[white] (0,2.0) circle (2.5pt); \draw[popblue, thick] (0,2.0) circle (2.5pt);
    \fill[white] (4,2.0) circle (2.5pt); \draw[popblue, thick] (4,2.0) circle (2.5pt);
    \node[below] at (0,2.0) {$a$};
    \node[below] at (4,2.0) {$b$};
    \node[left] at (-0.5,2.0) {\textbf{]a ; b[}};
    
    % Ligne 3 : [a ; b[
    \draw[->] (-0.5, 0.5) -- (5.5, 0.5);
    \draw[segment] (0,0.5) -- (4,0.5);
    \draw[crochet ferme] (0,0.3) -- (0,0.7);
    \draw[crochet ouvert] (4,0.3) -- (4,0.7);
    \fill[white] (4,0.5) circle (2.5pt); \draw[popblue, thick] (4,0.5) circle (2.5pt);
    \node[below] at (0,0.5) {$a$};
    \node[below] at (4,0.5) {$b$};
    \node[left] at (-0.5,0.5) {\textbf{[a ; b[}};
    
    % Ligne 4 : ]a ; b]
    \draw[->] (-0.5, -1.0) -- (5.5, -1.0);
    \draw[segment] (0,-1.0) -- (4,-1.0);
    \draw[crochet ouvert] (0,-1.2) -- (0,-0.8);
    \draw[crochet ferme] (4,-1.2) -- (4,-0.8);
    \fill[white] (0,-1.0) circle (2.5pt); \draw[popblue, thick] (0,-1.0) circle (2.5pt);
    \node[below] at (0,-1.0) {$a$};
    \node[below] at (4,-1.0) {$b$};
    \node[left] at (-0.5,-1.0) {\textbf{]a ; b]}};
\end{tikzpicture}
\legende{Représentation graphique des quatre intervalles bornés. \textbf{Point plein = borne incluse (crochet fermé)}. \textbf{Point vide = borne exclue (crochet ouvert)}.}
\end{popfigure}

\begin{methode}[Traduire une inégalité en intervalle (et réciproquement)]
\begin{enumerate}
    \item Identifier les bornes $a$ et $b$ (les nombres qui encadrent $x$).
    \item Regarder le sens des inégalités ($\le$ ou $<$ ; $\ge$ ou $>$).
    \item Choisir le crochet :
    \begin{itemize}
        \item $\le$ ou $\ge \to$ crochet \textbf{fermé} $[$ ou $]$.
        \item $<$ ou $> \to$ crochet \textbf{ouvert} $]$ ou $[$.
    \end{itemize}
    \item Écrire l'intervalle en mettant la plus petite borne à gauche.
\end{enumerate}
\end{methode}

\begin{exemplebox}[Traductions usuelles]
\begin{align*}
    2 \le x < 5 &\iff x \in [2 ; 5[ \\
    -3 < x \le 0 &\iff x \in ]-3 ; 0] \\
    x > 4 &\iff x \in ]4 ; +\infty[ \\
    x \le -1 &\iff x \in ]-\infty ; -1] \\
    x \in \mathbb{R} &\iff x \in ]-\infty ; +\infty[
\end{align*}
\end{exemplebox}

\courssub{Intersection et réunion d'intervalles (Initiation)}

En Seconde, vous manipulerez sans cesse des ensembles de solutions d'inéquations. Deux opérations fondamentales permettent de combiner des intervalles.

\begin{definition}[Intersection et Réunion]
Soient $I$ et $J$ deux intervalles (ou ensembles de réels).
\begin{itemize}
    \item L'\textbf{intersection} $I \cap J$ est l'ensemble des nombres qui appartiennent \textbf{à la fois} à $I$ \textbf{ET} à $J$ (le « commun »).
    \item La \textbf{réunion} $I \cup J$ est l'ensemble des nombres qui appartiennent à $I$ \textbf{OU} à $J$ (ou aux deux).
\end{itemize}
\end{definition}

\begin{propriete}[Intersection de deux intervalles]
L'intersection de deux intervalles est \textbf{toujours un intervalle} (éventuellement vide $\varnothing$).
Pour la trouver : on garde la \textbf{plus grande} des bornes inférieures et la \textbf{plus petite} des bornes supérieures. Attention au type de crochet (ouvert/fermé) : il faut prendre le plus « restrictif » (fermé l'emporte sur ouvert si la borne est commune).
\end{propriete}

\begin{exoresolu}[{Intersection de $I = [1 ; 5]$ et $J = ]3 ; 8]$}]
\begin{enumerate}
    \item \textbf{Borne inférieure :} $\max(1, 3) = 3$.
        $I$ a $[$ (inclut 1), $J$ a $]$ (exclut 3). La borne commune est $3$. Comme $J$ l'exclut, l'intersection l'exclut $\to$ crochet \textbf{ouvert} $]$.
    \item \textbf{Borne supérieure :} $\min(5, 8) = 5$.
        $I$ a $]$ (inclut 5), $J$ a $]$ (inclut 5). Les deux incluent $\to$ crochet \textbf{fermé} $]$.
    \item \textbf{Résultat :} $I \cap J = \boxed{]3 ; 5]}$.
\end{enumerate}
\end{exoresolu}

\begin{popfigure}
\begin{tikzpicture}[scale=1.2]
    % Intervalle I = [1;5]
    \draw[popblue, very thick] (1,2) -- (5,2);
    \draw[popblue, thick] (1,1.9) -- (1,2.1);
    \draw[popblue, thick] (5,1.9) -- (5,2.1);
    \node[popblue, left] at (0.5,2) {$I = [1;5]$};
    \node[popblue, below] at (1,2) {$1$};
    \node[popblue, below] at (5,2) {$5$};
    
    % Intervalle J = ]3;8]
    \draw[poporange, very thick] (3,1) -- (8,1);
    \draw[poporange, thick] (3,0.9) -- (3,1.1); \fill[white] (3,1) circle (2pt); \draw[poporange, thick] (3,1) circle (2pt);
    \draw[poporange, thick] (8,0.9) -- (8,1.1);
    \node[poporange, left] at (0.5,1) {$J = ]3;8]$};
    \node[poporange, below] at (3,1) {$3$};
    \node[poporange, below] at (8,1) {$8$};
    
    % Intersection I ∩ J = ]3;5]
    \draw[popgreen!70!black, very thick, line width=3pt, opacity=0.6] (3,0) -- (5,0);
    \draw[popgreen!70!black, thick] (3,-0.1) -- (3,0.1); \fill[white] (3,0) circle (2pt); \draw[popgreen!70!black, thick] (3,0) circle (2pt);
    \draw[popgreen!70!black, thick] (5,-0.1) -- (5,0.1);
    \node[popgreen!70!black, left] at (0.5,0) {$I \cap J = ]3;5]$};
    \node[popgreen!70!black, below] at (3,0) {$3$};
    \node[popgreen!70!black, below] at (5,0) {$5$};
    
    % Graduations communes
    \foreach \x in {1,2,3,4,5,6,7,8} {
        \draw (\x, -0.3) -- (\x, 2.3) node[above] {$\x$};
    }
\end{tikzpicture}
\legende{Visualisation de l'intersection $I \cap J$. La zone verte (superposition) correspond aux nombres appartenant aux deux intervalles.}
\end{popfigure}

\begin{attention}[Réunion : pas toujours un intervalle !]
Contrairement à l'intersection, la réunion de deux intervalles \textbf{n'est pas toujours un intervalle}.
Exemple : $I = [0;1]$ et $J = [2;3]$. $I \cup J$ a un « trou » entre 1 et 2. Ce n'est pas un intervalle unique.
Si les intervalles se chevauchent ou se touchent (ex: $[0;2] \cup [2;4] = [0;4]$), la réunion est un intervalle.
\end{attention}

\begin{savaistu}[L'infini en pratique : les tolérances industrielles]
Au Cameroun, comme partout, l'industrie (usines de transformation du cacao à Douala, cimenteries de Figuil, brasseries) utilise des \textbf{tolérances}. Une pièce doit mesurer $10,00 \text{ mm} \pm 0,02 \text{ mm}$.
Cela s'écrit mathématiquement : $x \in [9,98 ; 10,02]$.
Si la machine produit une pièce à $10,021 \text{ mm}$, elle est \textbf{hors tolérance} (rebut). La rigueur des intervalles (crochet fermé = accepté, crochet ouvert = refusé) a un impact économique direct : coût matière, sécurité, exportation. Les ingénieurs qualité passent leur vie à calculer des intersections d'intervalles de tolérance !
\end{savaistu}

\begin{exoresolu}[Synthèse : Encadrement, comparaison et intervalles]
Soit $x = \sqrt{30}$.
\begin{enumerate}
    \item Encadrer $x$ entre deux entiers consécutifs.
    \item Encadrer $x$ au dixième près.
    \item Écrire l'ensemble des valeurs de $x$ sous forme d'intervalle au dixième près.
    \item Comparer $x$ et $5,4$. Justifier.
\end{enumerate}

\textbf{Solution détaillée :}
\begin{enumerate}
    \item $5^2 = 25 \le 30 < 36 = 6^2$. Donc $\boxed{5 \le \sqrt{30} < 6}$.
    \item On teste les dixièmes à partir de 5 :
        $5,4^2 = 29,16 \le 30$.
        $5,5^2 = 30,25 > 30$.
        Donc $\boxed{5,4 \le \sqrt{30} < 5,5}$.
    \item L'encadrement au dixième s'écrit : $x \in \boxed{[5,4 ; 5,5[}$.
    \item Comparer $\sqrt{30}$ et $5,4$.
        Les deux sont positifs. On compare les carrés :
        $(\sqrt{30})^2 = 30$.
        $5,4^2 = 29,16$.
        $30 > 29,16 \implies \sqrt{30} > 5,4$.
        \textbf{Conclusion :} $\boxed{\sqrt{30} > 5,4}$.
        \emph{Remarque :} Cela se voit aussi sur l'encadrement : $\sqrt{30} \in [5,4 ; 5,5[$, donc il est $\ge 5,4$. Comme $30 \neq 29,16$, c'est strictement $>$.
\end{enumerate}
\end{exoresolu}

\begin{exercice}[Entraînement : Comparer, Encadrer, Intervalle]
\begin{enumerate}
    \item Comparer les nombres suivants (justifier par le carré) :
    \begin{itemize}
        \item[a)] $\sqrt{17}$ et $4$
        \item[b)] $\sqrt{99}$ et $10$
        \item[c)] $\sqrt{5}+2$ et $4$
    \end{itemize}
    \item Encadrer $\sqrt{75}$ au dixième près, puis au centième près. Donner l'arrondi au centième.
    \item Traduire en inégalité double (encadrement) puis en intervalle :
    \begin{itemize}
        \item[a)] $x \in ]-2 ; 5]$
        \item[b)] $x \in [3 ; +\infty[$
        \item[c)] $-1 < x \le 4$
    \end{itemize}
    \item Soit $A = [0 ; 4]$ et $B = ]2 ; 6]$. Déterminer $A \cap B$ et $A \cup B$. Représenter sur une droite graduée.
    \item \textbf{Problème concret :} Un menuisier à Yaoundé doit couper une poutre de longueur $L = \sqrt{120}$ cm. Sa scie a une précision au millimètre (dixième de cm).
    \begin{itemize}
        \item[a)] Entre quelles deux valeurs (en cm, au mm près) se trouve la longueur exacte ?
        \item[b)] S'il coupe à la valeur approchée par défaut au mm, quelle longueur de bois reste-t-il (perte) ? Et s'il coupe à la valeur par excès ?
    \end{itemize}
\end{enumerate}
\end{exercice}


\begin{corrige}[Exercice n°1]
\begin{enumerate}
    \item[a)] $4^2=16 < 17=(\sqrt{17})^2 \implies \sqrt{17} > 4$.
    \item[b)] $10^2=100 > 99=(\sqrt{99})^2 \implies \sqrt{99} < 10$.
    \item[c)] $(\sqrt{5}+2)^2 = 5 + 4\sqrt{5} + 4 = 9 + 4\sqrt{5}$. $4^2=16$. Comparer $9+4\sqrt{5}$ et $16 \iff 4\sqrt{5}$ et $7 \iff \sqrt{5}$ et $1,75$. $5 > 3,0625 \implies \sqrt{5} > 1,75 \implies \sqrt{5}+2 > 4$.
    \item $\sqrt{75}$ : $8^2=64, 9^2=81 \to 8 \le \sqrt{75} < 9$.
    Dixièmes : $8,6^2=73,96 \le 75 < 75,69=8,7^2 \to 8,6 \le \sqrt{75} < 8,7$.
    Centièmes : $8,66^2=74,9956 \le 75 < 75,1689=8,67^2 \to 8,66 \le \sqrt{75} < 8,67$.
    Arrondi au centième : $8,66$ (plus proche de $75$ que $8,67^2$).
    \item[a)] $-2 < x \le 5$. [b)] $x \ge 3$. [c)] $x \in ]-1 ; 4]$.
    \item $A \cap B = ]2 ; 4]$ (borne inf max(0,2)=2 ouverte car B ouverte ; borne sup min(4,6)=4 fermée car A fermée).
    $A \cup B = [0 ; 6]$ (ils se chevauchent sur ]2;4], réunion = intervalle continu de 0 à 6).
    \item[a)] $L=\sqrt{120}$. $10^2=100, 11^2=121 \to 10 \le L < 11$.
    Dixièmes : $10,9^2=118,81 \le 120 < 121=11,0^2$. Donc $L \in [10,9 ; 11,0[$.
    En mm : entre $10,9$ cm et $11,0$ cm.
    [b)] Défaut : $10,9$ cm. Perte = $L - 10,9 \approx 10,954 - 10,9 = 0,054$ cm ($0,54$ mm).
    Excès : $11,0$ cm. Impossible physiquement (bois trop court), ou alors il faut une pièce de $11,0$ cm et il gaspille $11,0 - L \approx 0,046$ cm.
\end{enumerate}
\end{corrige}

\begin{aretenir}[Bilan : Comparaison, Encadrement, Intervalles]
\begin{itemize}
    \item \textbf{Comparer} $A$ et $B$ ($\ge 0$) : comparer $A^2$ et $B^2$.
    \item \textbf{Encadrer} $\sqrt{N}$ : chercher entiers, puis dixièmes, puis centièmes... en testant les carrés.
    \item \textbf{Amplitude} = borne sup $-$ borne inf. Précision $= 10^{-n}$ pour $n$ décimales.
    \item \textbf{Intervalle} : Notation $[a;b], ]a;b[, [a;b[, ]a;b]$. Crochet \textbf{fermé} $[$ ou $]$ = borne \textbf{incluse} ($\le, \ge$). Crochet \textbf{ouvert} $]$ ou $[$ = borne \textbf{exclue} ($<, >$). $\infty$ toujours avec crochet ouvert.
    \item \textbf{Intersection} $I \cap J$ : commun (toujours un intervalle). \textbf{Réunion} $I \cup J$ : ensemble des deux (pas toujours un intervalle).
\end{itemize}
\end{aretenir}

\newpage

\coursec{Activité d'intégration}

\courssub{Objectifs de l'activité}
Cette activité de synthèse vise à mobiliser l'ensemble des savoirs et savoir-faire de l'unité « Nombres réels » dans une situation concrète d'aménagement urbain. Elle permet à l'élève de :
\begin{itemize}
    \item Manipuler les racines carrées et les opérations sur les radicaux (somme, produit, quotient, simplification) pour calculer des grandeurs géométriques exactes.
    \item Comparer des nombres réels comportant des radicaux en utilisant la comparaison par les carrés (croissance de la fonction racine carrée sur $\mathbb{R}_+$).
    \item Effectuer des encadrements décimaux (au dixième, à l'entier) pour prendre des décisions budgétaires réalistes.
    \item Utiliser la notation d'intervalles pour modéliser une contrainte technique (zone d'éclairement d'un projecteur).
    \item Rédiger une argumentation mathématique complète, structurée et justifiée, adaptée à un destinataire non mathématicien (le maire).
\end{itemize}

\courssub{Situation}
La mairie de Bafoussam, dans la région de l'Ouest, lance un projet d'aménagement d'un terrain de football de quartier sur une parcelle rectangulaire d'une superficie de \textbf{5000 m²}. Le cahier des charges impose que la \textbf{longueur} du terrain soit exactement le \textbf{double de sa largeur}.

Pour la clôture, l'entreprise locale « \emph{Grillage Express} » propose un grillage rigide au prix de \textbf{2 500 FCFA le mètre linéaire}.

Pour l'éclairage, un projecteur de forte puissance doit être installé au centre du terrain. La notice technique du fabricant (importé via le port de Douala) précise que l'éclairement au sol est conforme aux normes de la FECAFOOT (Fédération Camerounaise de Football) si et seulement si la distance entre le projecteur et le point éclairé appartient à l'intervalle \textbf{[40 ; 60] mètres}.

Un technicien propose de fixer le projecteur sur un poteau déjà existant, situé précisément à l'un des \textbf{coins du terrain}.

\emph{En tant qu'expert mathématicien consultant pour la mairie, votre mission est de valider la faisabilité technique et budgétaire de ce projet.}

\begin{popfigure}
\begin{center}
\begin{tikzpicture}[scale=0.055]
    \draw[fill=popgreenL,draw=popdark,thick] (0,0) rectangle (100,50);
    \draw[popblue,dashed] (0,0) -- (100,50);
    \draw[popblue,dashed] (0,50) -- (100,0);
    \fill[popink] (50,25) circle (1.2) node[below right=1pt] {$O$};
    \fill[poporange] (100,50) circle (1.2) node[above left=1pt] {$A$};
    \draw[poporange,very thick] (50,25) -- (100,50) node[midway,above,sloped] {\small $OA$};
    \node[below] at (50,0) {Longueur $= 2x$};
    \node[left,rotate=90] at (0,25) {Largeur $= x$};
\end{tikzpicture}
\end{center}
\legende{Schéma du terrain : rectangle de largeur $x$ et de longueur $2x$, de centre $O$ (milieu des diagonales) ; le poteau est au sommet $A$.}
\end{popfigure}

\courssub{Consignes}
\emph{Toutes les réponses doivent être justifiées par un raisonnement écrit. Les calculs exacts (sous forme de radicaux simplifiés) sont exigés avant tout encadrement décimal.}

\begin{enumerate}
    \item \textbf{Dimensions du terrain.} Soit $x$ la largeur du terrain en mètres. Exprimer la longueur en fonction de $x$, puis écrire l'équation traduisant la superficie. Résoudre cette équation pour trouver la largeur $x$ et la longueur. Vérifier que les valeurs trouvées sont compatibles avec le contexte.
    \item \textbf{Longueur des diagonales.} Calculer la longueur exacte d'une diagonale du terrain (sous la forme $a\sqrt{b}$ avec $a,b$ entiers, $b$ le plus petit possible). Donner ensuite un encadrement de cette longueur au dixième de mètre près.
    \item \textbf{Coût du grillage.} Calculer le périmètre exact du terrain. En déduire le coût total exact du grillage (en FCFA). Exprimer ce coût en milliers de FCFA pour établir un budget prévisionnel.
    \item \textbf{Validation de l'éclairage.} Le projecteur est au centre $O$ du terrain. Le poteau existant est au sommet $A$ du rectangle. Calculer la distance exacte $OA$. En utilisant la comparaison par les carrés, déterminer si la distance $OA$ appartient à l'intervalle $[40 ; 60]$. Conclure sur la conformité de l'installation proposée.
    \item \textbf{Synthèse et recommandation.} Rédiger une note courte (10 à 15 lignes) adressée au Maire de Bafoussam, résumant les résultats techniques (dimensions, diagonale, coût, éclairage) et formulant une recommandation claire (valider / modifier / rejeter le projet tel que proposé) avec justification mathématique.
\end{enumerate}

\courssub{Ressources à mobiliser}
\begin{itemize}
    \item \textbf{Racine carrée :} Pour tout $a \ge 0$, $\sqrt{a}$ est l'unique réel positif dont le carré est $a$. $\sqrt{a^2}=a$ si $a\ge 0$.
    \item \textbf{Simplification de radicaux :} $\sqrt{ab}=\sqrt{a}\times\sqrt{b}$ ($a,b\ge 0$) ; on cherche à extraire les facteurs carrés parfaits (4, 9, 16, 25, 36, 49, 64, 81, 100...).
    \item \textbf{Opérations :} $k\sqrt{a} \pm k'\sqrt{a} = (k\pm k')\sqrt{a}$ ; $\sqrt{a}\times\sqrt{b}=\sqrt{ab}$ ; $\dfrac{\sqrt{a}}{\sqrt{b}}=\sqrt{\dfrac{a}{b}}$ ($b>0$). Rationalisation : $\dfrac{k}{\sqrt{a}} = \dfrac{k\sqrt{a}}{a}$.
    \item \textbf{Comparaison :} Pour $a,b \ge 0$, $a < b \iff \sqrt{a} < \sqrt{b} \iff a^2 < b^2$. On compare souvent un radical à un entier en élevant au carré.
    \item \textbf{Encadrement :} Si $n^2 < A < (n+1)^2$ alors $n < \sqrt{A} < n+1$. On affine ensuite en testant les dixièmes ($2,2^2$, $2,3^2$, ...).
    \item \textbf{Intervalles :} $[a;b] = \{x \in \mathbb{R} \mid a \le x \le b\}$.
    \item \textbf{Géométrie :} Aire du rectangle $= L \times l$ ; Périmètre $= 2(L+l)$ ; Diagonale $d = \sqrt{L^2+l^2}$ (théorème de Pythagore) ; le centre du rectangle est le milieu des diagonales, donc la distance centre-sommet vaut $\dfrac{d}{2}$.
\end{itemize}

\courssub{Démarche et corrigé commenté}
\begin{exoresolu}{Corrigé détaillé de l'activité d'intégration}
\textbf{1. Dimensions du terrain.}\par
Soit $x$ la largeur en mètres. La longueur est $2x$.
La superficie $S$ est donnée par $S = \text{longueur} \times \text{largeur} = 2x \times x = 2x^2$.
On a donc l'équation : $2x^2 = 5000$.
\emph{Résolution :} $x^2 = \dfrac{5000}{2} = 2500$.
Puisque $x$ est une longueur, $x > 0$, donc $x = \sqrt{2500}$.
Or $50^2 = 2500$, donc $\sqrt{2500} = 50$.
\emph{Résultat :} Largeur $= 50$ m ; Longueur $= 2 \times 50 = 100$ m.
\emph{Vérification :} $50 \times 100 = 5000$ m². Les valeurs sont positives et cohérentes avec le contexte.

\textbf{2. Longueur des diagonales.}\par
Dans le rectangle, la diagonale $d$ vérifie, d'après le théorème de Pythagore :
\[ d^2 = 100^2 + 50^2 = 10000 + 2500 = 12500. \]
Donc $d = \sqrt{12500}$.
\emph{Simplification du radical :} on cherche le plus grand carré parfait divisant $12500$ :
\[ 12500 = 2500 \times 5 \quad\text{donc}\quad \sqrt{12500} = \sqrt{2500 \times 5} = \sqrt{2500} \times \sqrt{5} = 50\sqrt{5}. \]
La forme exacte est \textbf{$50\sqrt{5}$ mètres}.
\emph{Encadrement au dixième :} on encadre d'abord $\sqrt{5}$.
\begin{align*}
2^2 = 4 < 5 < 9 = 3^2 &\Rightarrow 2 < \sqrt{5} < 3,\\
2,2^2 = 4,84 < 5 < 5,29 = 2,3^2 &\Rightarrow 2,2 < \sqrt{5} < 2,3,\\
2,23^2 = 4,9729 < 5 < 5,0176 = 2,24^2 &\Rightarrow 2,23 < \sqrt{5} < 2,24.
\end{align*}
En multipliant par 50 : $50 \times 2,23 < 50\sqrt{5} < 50 \times 2,24$, c'est-à-dire $111,5 < d < 112$. Pour trancher au dixième, on affine : $2,236^2 = 4,999696 < 5 < 5,000049 = 2,237^2$, donc $2,236 < \sqrt{5} < 2,237$, puis :
\[ 111,8 \text{ m} < d < 111,85 \text{ m}, \quad\text{donc}\quad 111,8 \text{ m} < d < 111,9 \text{ m}. \]
La diagonale mesure environ \textbf{111,8 m}.

\textbf{3. Coût du grillage.}\par
Périmètre : $P = 2 \times (\text{Longueur} + \text{Largeur}) = 2 \times (100 + 50) = 2 \times 150 = 300$ m.
Coût total : $C = P \times 2500 = 300 \times 2500 = 750\,000$ FCFA.
\emph{Budget prévisionnel :} $750\,000$ FCFA $= 750 \times 1000$ FCFA. Le coût est exactement \textbf{750 milliers de FCFA}, valeur entière : aucun encadrement n'est nécessaire ici, le budget est connu exactement. (Si le périmètre avait contenu un radical, on aurait encadré le coût entre deux entiers consécutifs de milliers de FCFA.)

\textbf{4. Validation de l'éclairage.}\par
Le centre $O$ du rectangle est le milieu des diagonales. La distance du centre au sommet $A$ est donc la demi-diagonale :
\[ OA = \frac{d}{2} = \frac{50\sqrt{5}}{2} = 25\sqrt{5} \text{ mètres}. \]
Il faut tester si $OA \in [40 ; 60]$, c'est-à-dire si $40 \le 25\sqrt{5} \le 60$.
\emph{Méthode : comparaison par les carrés (tous les nombres en jeu sont positifs, ce qui autorise cette méthode).}
On calcule d'abord : $(25\sqrt{5})^2 = 25^2 \times (\sqrt{5})^2 = 625 \times 5 = 3125$.
\begin{itemize}
    \item \textbf{Borne inférieure :} $40^2 = 1600$. Comme $3125 > 1600$, on a $25\sqrt{5} > 40$. \emph{Condition respectée.}
    \item \textbf{Borne supérieure :} $60^2 = 3600$. Comme $3125 < 3600$, on a $25\sqrt{5} < 60$. \emph{Condition respectée.}
\end{itemize}
Donc $40 < 25\sqrt{5} < 60$, ce qui signifie que $25\sqrt{5} \in [40 ; 60]$.
\emph{Conclusion technique :} la distance du centre au coin ($25\sqrt{5} \approx 55,9$ m) appartient bien à l'intervalle de conformité. \textbf{L'installation du projecteur sur le poteau au coin est validée.}

\textbf{5. Synthèse et recommandation (Note au Maire).}\par
\textbf{Objet :} Avis technique sur l'aménagement du terrain de football de quartier — Bafoussam.

Monsieur le Maire,

Suite à l'étude du dossier d'aménagement de la parcelle de 5000 m², je vous confirme la faisabilité technique du projet tel que proposé.

Le terrain respectera les proportions imposées (longueur = 2 × largeur) avec des dimensions exactes de \textbf{100 m × 50 m}. La diagonale mesure exactement $50\sqrt{5}$ m (environ 111,8 m).

Le budget prévisionnel pour la clôture (grillage à 2 500 FCFA/m) s'élève exactement à \textbf{750 000 FCFA} (750 milliers de FCFA), sans marge d'incertitude.

Concernant l'éclairage, l'installation du projecteur sur le poteau existant au coin du terrain est \textbf{conforme aux normes FECAFOOT} : la distance centre-poteau ($25\sqrt{5}$ m $\approx$ 55,9 m) appartient bien à l'intervalle requis $[40 ; 60]$ m. Aucun déplacement de poteau n'est nécessaire, ce qui génère une économie substantielle.

\emph{Recommandation :} \textbf{valider le projet en l'état.} Les travaux peuvent être lancés immédiatement selon ce cahier des charges.

Respectueusement,
L'Expert Mathématicien Consultant.
\end{exoresolu}

\begin{attention}[Pièges fréquents dans ce type d'activité]
\begin{itemize}
    \item \textbf{Oublier la condition de positivité :} l'équation $x^2 = 2500$ admet deux solutions ($50$ et $-50$), mais une longueur est positive : on ne retient que $x = 50$.
    \item \textbf{Comparer sans élever au carré :} écrire « $25\sqrt{5} < 60$ car $\sqrt{5} < 2,4$ » sans justification est insuffisant au BEPC ; la comparaison par les carrés ($3125 < 3600$) est la méthode rigoureuse attendue.
    \item \textbf{Mal simplifier le radical :} $\sqrt{12500} \neq 125\sqrt{...}$ ; il faut extraire le \emph{plus grand} carré parfait : $12500 = 2500 \times 5$.
    \item \textbf{Confondre diagonale et demi-diagonale :} la distance centre-sommet est $\dfrac{d}{2}$, pas $d$.
    \item \textbf{Mélanger exact et approché :} on raisonne avec la valeur exacte $25\sqrt{5}$ et l'on n'arrondit qu'à la toute fin, pour la communication.
\end{itemize}
\end{attention}

\courssub{Bilan de l'activité}
Cette activité a permis de mettre en évidence la puissance de l'outil « nombres réels et radicaux » pour modéliser et résoudre un problème réel d'ingénierie locale. Les points clés acquis ou consolidés sont :
\begin{itemize}
    \item La \textbf{modélisation algébrique} (équation $2x^2=5000$) résolue par extraction de racine carrée, en retenant le signe imposé par le contexte géométrique ($x > 0$).
    \item La \textbf{simplification systématique des radicaux} ($\sqrt{12500}=50\sqrt{5}$) comme forme exacte indispensable avant tout calcul approché.
    \item L'articulation \textbf{exact / approché} : la forme exacte ($50\sqrt{5}$) sert au raisonnement rigoureux (comparaison par les carrés), l'encadrement décimal (111,8 m) sert à la représentation concrète et au chiffrage.
    \item L'utilisation de la \textbf{comparaison par les carrés} pour trancher une appartenance à un intervalle sans calculatrice, méthode rigoureuse et attendue au BEPC.
    \item La \textbf{traduction en langage naturel} (note au maire) d'un raisonnement mathématique, compétence transversale essentielle (savoir-être / communication).
    \item La gestion des \textbf{unités et grandeurs} (m, m², FCFA, milliers de FCFA) dans un contexte camerounais authentique (Bafoussam, FECAFOOT, coûts locaux).
\end{itemize}
L'élève a ainsi réinvesti l'ensemble du programme de l'unité (racines carrées, opérations sur les radicaux, puissances, comparaison, encadrement, intervalles) dans une démarche unique, cohérente et professionnalisante.

\newpage

\coursec{Méthodes et exercices résolus}

\courssub{Simplifier une expression contenant des racines carrées}

\begin{methode}[Simplifier et calculer avec des racines carrées]
\textbf{Objectif :} Écrire sous la forme $a\sqrt{b}$ avec $a \in \mathbb{Q}$ et $b \in \mathbb{N}$ le plus petit possible (radical simplifié).
\begin{enumerate}
    \item \textbf{Identifier les carrés parfaits} : repère les diviseurs de l'entier sous le radical qui sont des carrés ($4, 9, 16, 25, 36, 49, 64, 81, 100\dots$).
    \item \textbf{Décomposer} : écris le radicande en produit $k^2 \times b$ où $k^2$ est le plus grand carré parfait diviseur.
    \item \textbf{Appliquer la propriété} : $\sqrt{k^2 \times b} = \sqrt{k^2} \times \sqrt{b} = k\sqrt{b}$ (car $k \ge 0$).
    \item \textbf{Simplifier les coefficients} : si l'expression contient plusieurs radicaux, réduis-les tous avant d'additionner ou soustraire les termes \emph{semblables} (même radicande).
    \item \textbf{Vérifier} : le nombre sous le radical n'a plus de diviseur carré parfait (autre que 1).
\end{enumerate}
\cle{Règle d'or} : $\sqrt{a^2} = |a|$. Pour $a \ge 0$, $\sqrt{a^2}=a$. Au BEPC, on précise souvent « $a$ réel positif » pour lever l'ambiguïté.
\end{methode}

\begin{exoresolu}[Simplification et calcul exact]
\textbf{Énoncé :} Soit l'expression $E = 3\sqrt{75} - 2\sqrt{12} + \sqrt{27}$.
\begin{enumerate}
    \item Simplifier chaque radical.
    \item Écrire $E$ sous la forme $a\sqrt{b}$ avec $a,b \in \mathbb{N}$.
    \item Donner une valeur approchée de $E$ au millième près.
\end{enumerate}
\tcblower
\textbf{Solution détaillée :}
\begin{enumerate}
    \item \textbf{Simplification des radicaux :}
    \begin{itemize}
        \item $\sqrt{75} = \sqrt{25 \times 3} = \sqrt{25}\times\sqrt{3} = 5\sqrt{3}$.
        \item $\sqrt{12} = \sqrt{4 \times 3} = \sqrt{4}\times\sqrt{3} = 2\sqrt{3}$.
        \item $\sqrt{27} = \sqrt{9 \times 3} = \sqrt{9}\times\sqrt{3} = 3\sqrt{3}$.
    \end{itemize}
    \item \textbf{Calcul de $E$ :} On remplace dans l'expression :
    \begin{align*}
    E &= 3(5\sqrt{3}) - 2(2\sqrt{3}) + 3\sqrt{3} \\
      &= 15\sqrt{3} - 4\sqrt{3} + 3\sqrt{3} \quad \text{(termes semblables : même radicande $\sqrt{3}$)} \\
      &= (15 - 4 + 3)\sqrt{3} \\
      &= 14\sqrt{3}.
    \end{align*}
    Donc $a=14$ et $b=3$.
    \item \textbf{Valeur approchée :} $\sqrt{3} \approx 1,73205\dots$
    $E \approx 14 \times 1,73205 = 24,2487\dots$
    Arrondi au millième : $\boxed{E \approx 24,249}$.
\end{enumerate}
\textbf{Conclusion :} $E = 14\sqrt{3} \approx 24,249$.
\end{exoresolu}

\courssub{Identifier et classifier les nombres réels}

\begin{methode}[Déterminer la nature d'un nombre réel (rationnel / irrationnel)]
\textbf{Objectif :} Dire si un nombre donné appartient à $\mathbb{Q}$ ou à $\mathbb{R}\setminus\mathbb{Q}$ (irrationnel).
\begin{enumerate}
    \item \textbf{Rappel des définitions} :
    \begin{itemize}
        \item $x \in \mathbb{Q} \iff$ il existe $p \in \mathbb{Z}, q \in \mathbb{N}^*$ tel que $x = \frac{p}{q}$ (écriture fractionnaire). En écriture décimale : développement \textbf{fini} ou \textbf{périodique}.
        \item $x \notin \mathbb{Q}$ (irrationnel) $\iff$ développement décimal \textbf{infini non périodique}. Les racines carrées d'entiers non carrés parfaits sont irrationnels ($\sqrt{2}, \sqrt{3}, \sqrt{5}, \sqrt{7}, \sqrt{10}\dots$).
    \end{itemize}
    \item \textbf{Analyser l'expression} :
    \begin{itemize}
        \item Entier relatif, décimal, fraction $\to \mathbb{Q}$.
        \item $\sqrt{n}$ avec $n$ non carré parfait $\to$ Irrationnel.
        \item Somme/Produit/Quotient : 
        \begin{itemize}
            \item Rationnel $\pm$ Rationnel = Rationnel.
            \item Rationnel $\times$ Rationnel = Rationnel.
            \item Rationnel $\neq 0$ $\times$ Irrationnel = Irrationnel.
            \item Irrationnel $\pm$ Irrationnel = \emph{Indéterminé} (peut être rationnel : $\sqrt{2} + (-\sqrt{2}) = 0$).
        \end{itemize}
    \end{itemize}
    \item \textbf{Conclure} par une phrase : « $x$ est rationnel car... » ou « $x$ est irrationnel car... ».
\end{enumerate}
\end{methode}

\begin{exoresolu}[Classification de nombres réels]
\textbf{Énoncé :} Parmi les nombres suivants, lesquels sont rationnels ? Justifier chaque réponse.
\[ A = \frac{22}{7} \quad ; \quad B = \sqrt{0,04} \quad ; \quad C = \sqrt{2} + \sqrt{8} \quad ; \quad D = \frac{\pi}{2} \quad ; \quad E = 0,\overline{142857} \]
\tcblower
\textbf{Solution détaillée :}
On examine chaque nombre :
\begin{itemize}
    \item \textbf{$A = \frac{22}{7}$} : C'est un quotient de deux entiers ($22 \in \mathbb{Z}, 7 \in \mathbb{N}^*$). Donc $\boxed{A \in \mathbb{Q}}$.
    \item \textbf{$B = \sqrt{0,04}$} : $0,04 = \frac{4}{100} = \frac{1}{25} = \left(\frac{1}{5}\right)^2$. Donc $\sqrt{0,04} = \frac{1}{5} = 0,2$. C'est un nombre décimal, donc rationnel. $\boxed{B \in \mathbb{Q}}$.
    \item \textbf{$C = \sqrt{2} + \sqrt{8}$} : Simplifions $\sqrt{8} = \sqrt{4\times 2} = 2\sqrt{2}$.
    $C = \sqrt{2} + 2\sqrt{2} = 3\sqrt{2}$.
    $\sqrt{2}$ est irrationnel (carré non parfait). $3 \in \mathbb{Q}^*$. Produit d'un rationnel non nul par un irrationnel = irrationnel. $\boxed{C \notin \mathbb{Q}}$.
    \item \textbf{$D = \frac{\pi}{2}$} : $\pi$ est irrationnel (admis en 3ème). $\frac{1}{2} \in \mathbb{Q}^*$. Produit d'un rationnel non nul par un irrationnel = irrationnel. $\boxed{D \notin \mathbb{Q}}$.
    \item \textbf{$E = 0,\overline{142857}$} : La barre au-dessus des chiffres indique une écriture décimale \textbf{périodique} (période 142857). Tout nombre à développement décimal périodique est rationnel. $\boxed{E \in \mathbb{Q}}$.
    \emph{(Note : $E = \frac{1}{7}$)}.
\end{itemize}
\textbf{Bilan :} Les nombres rationnels sont $A, B, E$. Les nombres irrationnels sont $C, D$.
\end{exoresolu}

\courssub{Calculer une somme ou un produit d'expressions radicales}

\begin{methode}[Développer, factoriser et simplifier des expressions avec radicaux]
\textbf{Objectif :} Maîtriser le calcul littéral sur les expressions de la forme $a\sqrt{b} \pm c\sqrt{d}$.
\begin{enumerate}
    \item \textbf{Simplifier chaque radical} au préalable (Méthode 1) pour faire apparaître les termes semblables.
    \item \textbf{Addition / Soustraction} : On ne combine que les termes ayant \textbf{exactement le même radicande}.
    \[ k\sqrt{n} + k'\sqrt{n} = (k+k')\sqrt{n} \]
    \item \textbf{Multiplication} : On utilise la distributivité et la règle $\sqrt{a}\times\sqrt{b} = \sqrt{ab}$ (pour $a,b \ge 0$).
    \begin{itemize}
        \item $(a\sqrt{b})(c\sqrt{d}) = ac\sqrt{bd}$.
        \item Identités remarquables (valables car $\sqrt{b} \ge 0$) :
        \begin{align*}
        (x+y)^2 &= x^2 + 2xy + y^2 \\
        (x-y)^2 &= x^2 - 2xy + y^2 \\
        (x-y)(x+y) &= x^2 - y^2 \quad \text{(Très utile pour rationaliser !)}
        \end{align*}
    \end{itemize}
    \item \textbf{Simplifier le résultat final} (radical simplifié, coefficients réduits).
\end{enumerate}
\cle{Attention} : $\sqrt{a+b} \neq \sqrt{a} + \sqrt{b}$ et $\sqrt{a-b} \neq \sqrt{a} - \sqrt{b}$.
\end{methode}

\begin{exoresolu}[Développement et factorisation avec radicaux]
\textbf{Énoncé :} Soient $x = \sqrt{5} + \sqrt{2}$ et $y = \sqrt{5} - \sqrt{2}$.
\begin{enumerate}
    \item Calculer $x+y$ et $x-y$.
    \item Calculer $xy$.
    \item Calculer $x^2 + y^2$.
    \item En déduire la valeur de $x^4 + y^4$.
\end{enumerate}
\tcblower
\textbf{Solution détaillée :}
\begin{enumerate}
    \item \textbf{Somme et différence :}
    \begin{align*}
    x + y &= (\sqrt{5} + \sqrt{2}) + (\sqrt{5} - \sqrt{2}) = 2\sqrt{5}. \\
    x - y &= (\sqrt{5} + \sqrt{2}) - (\sqrt{5} - \sqrt{2}) = 2\sqrt{2}.
    \end{align*}
    \item \textbf{Produit $xy$ :} On utilise l'identité $(a+b)(a-b)=a^2-b^2$ avec $a=\sqrt{5}, b=\sqrt{2}$.
    \[ xy = (\sqrt{5})^2 - (\sqrt{2})^2 = 5 - 2 = 3. \]
    \item \textbf{Calcul de $x^2 + y^2$ :} On utilise $(x+y)^2 = x^2 + y^2 + 2xy \iff x^2+y^2 = (x+y)^2 - 2xy$.
    \begin{align*}
    x^2 + y^2 &= (2\sqrt{5})^2 - 2\times 3 \\
              &= 4 \times 5 - 6 \\
              &= 20 - 6 = 14.
    \end{align*}
    \item \textbf{Calcul de $x^4 + y^4$ :} On utilise $x^4+y^4 = (x^2+y^2)^2 - 2x^2y^2 = (x^2+y^2)^2 - 2(xy)^2$.
    \begin{align*}
    x^4 + y^4 &= 14^2 - 2 \times 3^2 \\
              &= 196 - 2 \times 9 \\
              &= 196 - 18 = 178.
    \end{align*}
\end{enumerate}
\textbf{Conclusion :} $x+y=2\sqrt{5},\; x-y=2\sqrt{2},\; xy=3,\; x^2+y^2=14,\; x^4+y^4=178$.
\end{exoresolu}

\courssub{Rationaliser le dénominateur d'une fraction}

\begin{methode}[Mettre un quotient de réels sous forme standard (dénominateur rationnel)]
\textbf{Objectif :} Transformer $\frac{A}{B}$ avec $B$ contenant des radicaux en $\frac{A'}{B'}$ avec $B' \in \mathbb{Q}$.
\begin{enumerate}
    \item \textbf{Identifier la forme du dénominateur} :
    \begin{itemize}
        \item \textbf{Cas 1 : Monôme} $\sqrt{b}$ ou $k\sqrt{b}$. $\to$ Multiplier numérateur et dénominateur par $\sqrt{b}$.
        \item \textbf{Cas 2 : Binôme} $a\sqrt{b} \pm c\sqrt{d}$ (somme/différence de deux radicaux). $\to$ Multiplier numérateur et dénominateur par le \textbf{conjugué} : $a\sqrt{b} \mp c\sqrt{d}$.
        \item \textbf{Cas 3 : Trinôme ou plus} : Rare en 3ème, on procède par étapes successives (regrouper deux termes, traiter comme binôme).
    \end{itemize}
    \item \textbf{Appliquer} : $\frac{N}{D} = \frac{N \times C}{D \times C}$ où $C$ est le facteur choisi.
    \item \textbf{Calculer le nouveau dénominateur} : Il doit devenir un entier (grâce à $(x-y)(x+y)=x^2-y^2$ ou $(\sqrt{b})^2=b$).
    \item \textbf{Développer le numérateur} et simplifier par un facteur commun éventuel avec le dénominateur.
    \item \textbf{Écrire la réponse finale} sous la forme $\frac{a\sqrt{b} + c\sqrt{d}}{k}$ ou $a\sqrt{b} + c\sqrt{d}$.
\end{enumerate}
\cle{Règle} : On ne laisse \textbf{jamais} de radical au dénominateur dans une réponse finale au BEPC.
\end{methode}

\begin{exoresolu}[Rationalisation et simplification]
\textbf{Énoncé :} Simplifier les expressions suivantes de façon à ce que le dénominateur soit un entier.
\begin{enumerate}
    \item $A = \dfrac{6}{\sqrt{3}}$
    \item $B = \dfrac{\sqrt{5}+1}{\sqrt{5}-1}$
    \item $C = \dfrac{3\sqrt{2}}{2\sqrt{3}-\sqrt{6}}$
\end{enumerate}
\tcblower
\textbf{Solution détaillée :}
\begin{enumerate}
    \item \textbf{$A = \frac{6}{\sqrt{3}}$} (Cas monôme). On multiplie par $\sqrt{3}$ :
    \[ A = \frac{6 \times \sqrt{3}}{\sqrt{3} \times \sqrt{3}} = \frac{6\sqrt{3}}{3} = \boxed{2\sqrt{3}}. \]
    \item \textbf{$B = \frac{\sqrt{5}+1}{\sqrt{5}-1}$} (Cas binôme). Le conjugué du dénominateur est $\sqrt{5}+1$.
    \begin{align*}
    B &= \frac{(\sqrt{5}+1)(\sqrt{5}+1)}{(\sqrt{5}-1)(\sqrt{5}+1)} \\
      &= \frac{(\sqrt{5}+1)^2}{(\sqrt{5})^2 - 1^2} \quad \text{(Identité remarquable au dénominateur)} \\
      &= \frac{5 + 2\sqrt{5} + 1}{5 - 1} \\
      &= \frac{6 + 2\sqrt{5}}{4} = \frac{2(3+\sqrt{5})}{4} = \boxed{\frac{3+\sqrt{5}}{2}}.
    \end{align*}
    \item \textbf{$C = \frac{3\sqrt{2}}{2\sqrt{3}-\sqrt{6}}$} (Cas binôme). Conjugué : $2\sqrt{3}+\sqrt{6}$.
    \begin{align*}
    C &= \frac{3\sqrt{2}(2\sqrt{3}+\sqrt{6})}{(2\sqrt{3}-\sqrt{6})(2\sqrt{3}+\sqrt{6})} \\
      \text{Dénominateur : } & (2\sqrt{3})^2 - (\sqrt{6})^2 = 4\times 3 - 6 = 12 - 6 = 6. \\
      \text{Numérateur : } & 3\sqrt{2} \times 2\sqrt{3} + 3\sqrt{2} \times \sqrt{6} = 6\sqrt{6} + 3\sqrt{12}. \\
      & \sqrt{12} = 2\sqrt{3} \implies \text{Num} = 6\sqrt{6} + 6\sqrt{3} = 6(\sqrt{6}+\sqrt{3}). \\
      C &= \frac{6(\sqrt{6}+\sqrt{3})}{6} = \boxed{\sqrt{6}+\sqrt{3}}.
    \end{align*}
\end{enumerate}
\end{exoresolu}

\courssub{Calculer des puissances entières (positives, nulles, négatives)}

\begin{methode}[Maîtriser les lois des puissances entières sur $\mathbb{R}$]
\textbf{Objectif :} Calculer et simplifier des expressions comme $a^n$ avec $n \in \mathbb{Z}$.
\begin{enumerate}
    \item \textbf{Rappel des définitions} (pour $a \in \mathbb{R}^*$) :
    \begin{itemize}
        \item $a^0 = 1$.
        \item $a^{-n} = \frac{1}{a^n}$ pour $n \in \mathbb{N}^*$.
        \item $a^n = a \times a \times \dots \times a$ ($n$ fois) pour $n \in \mathbb{N}^*$.
    \end{itemize}
    \item \textbf{Lois des puissances} (valides pour $a,b \in \mathbb{R}^*$, $m,n \in \mathbb{Z}$) :
    \begin{align*}
    a^m \times a^n &= a^{m+n} & \frac{a^m}{a^n} &= a^{m-n} \\
    (a^m)^n &= a^{mn} & (ab)^n &= a^n b^n \\
    \left(\frac{a}{b}\right)^n &= \frac{a^n}{b^n}
    \end{align*}
    \item \textbf{Stratégie de simplification} :
    \begin{enumerate}
        \item Écrire tous les termes en \textbf{puissances de base simple} (ex: $4=2^2, 8=2^3, \sqrt{2}=2^{1/2}$ — attention, exposants fractionnaires hors programme 3ème, on reste en notation radicale ou puissance entière).
        \item Appliquer les lois pour regrouper les puissances de même base.
        \item Gérer les signes : $(-a)^n = a^n$ si $n$ pair, $=-a^n$ si $n$ impair.
        \item Résultat final : puissance unique par base, exposant entier (positif ou négatif), ou valeur numérique.
    \end{enumerate}
\end{enumerate}
\cle{Piège} : $-2^2 = -(2^2) = -4$ alors que $(-2)^2 = 4$. Les parenthèses changent tout !
\end{methode}

\begin{exoresolu}[Simplification d'expressions avec puissances entières]
\textbf{Énoncé :} Simplifier au maximum les expressions suivantes (écrire sous forme de puissance de base 2 ou 3, ou valeur exacte).
\begin{enumerate}
    \item $A = 2^{-3} \times 2^5 \times 2^{-1}$
    \item $B = \frac{3^4 \times 9^{-2}}{27^{-1}}$
    \item $C = \left(\frac{2^{-2}}{4^{-1}}\right)^3$
    \item $D = (\sqrt{2})^6$ \emph{(indice : $(\sqrt{a})^2 = a$)}
\end{enumerate}
\tcblower
\textbf{Solution détaillée :}
\begin{enumerate}
    \item \textbf{$A$} : Même base 2. On additionne les exposants :
    $A = 2^{-3+5-1} = 2^{1} = \boxed{2}$.
    \item \textbf{$B$} : Passage en base 3. $9 = 3^2 \Rightarrow 9^{-2} = (3^2)^{-2} = 3^{-4}$. $27 = 3^3 \Rightarrow 27^{-1} = 3^{-3}$.
    \begin{align*}
    B &= \frac{3^4 \times 3^{-4}}{3^{-3}} = \frac{3^{4-4}}{3^{-3}} = \frac{3^0}{3^{-3}} = \frac{1}{3^{-3}} = 3^3 = \boxed{27}.
    \end{align*}
    \item \textbf{$C$} : $4 = 2^2 \Rightarrow 4^{-1} = 2^{-2}$.
    \begin{align*}
    C &= \left(\frac{2^{-2}}{2^{-2}}\right)^3 = (1)^3 = \boxed{1}.
    \end{align*}
    \item \textbf{$D$} : $(\sqrt{2})^6 = ((\sqrt{2})^2)^3 = 2^3 = \boxed{8}$.
    \emph{(Ou : $\sqrt{2} = 2^{1/2}$, $(2^{1/2})^6 = 2^3 = 8$, mais notation radicale préférée en 3ème).}
\end{enumerate}
\end{exoresolu}

\courssub{Comparer et encadrer des nombres réels}

\begin{methode}[Comparer deux réels et encadrer un réel à $10^{-n}$ près]
\textbf{Objectif :} Ordonner des nombres, trouver un encadrement décimal.
\begin{enumerate}
    \item \textbf{Comparer $a$ et $b$ (avec radicaux)} :
    \begin{itemize}
        \item Si $a \ge 0$ et $b \ge 0$ : $a < b \iff a^2 < b^2$ (fonction carrée croissante sur $\mathbb{R}^+$).
        \item Stratégie : Mettre au carré les deux nombres (après simplification des radicaux) et comparer les entiers obtenus.
        \item Si signes différents : le négatif est plus petit.
    \end{itemize}
    \item \textbf{Encadrer un nombre $x$ à $10^{-n}$ près} (ex: au centième $10^{-2}$) :
    \begin{enumerate}
        \item Calculer une valeur approchée de $x$ (calculatrice ou méthode manuelle) avec au moins $n+1$ décimales.
        \item Écrire l'encadrement : $k \times 10^{-n} \le x < (k+1) \times 10^{-n}$.
        \item \textbf{Vérifier} par le calcul : $(k \times 10^{-n})^2 \le x^2 < ((k+1) \times 10^{-n})^2$ si $x=\sqrt{\dots}$.
    \end{enumerate}
    \item \textbf{Notation intervalle} : $[a;b] = \{x \in \mathbb{R} \mid a \le x \le b\}$, $]a;b[$, $[a;b[$, $]a;b]$. $+\infty$ et $-\infty$ toujours avec crochet ouvert $[$ ou $]$.
\end{enumerate}
\end{methode}

\begin{exoresolu}[Comparaison, encadrement et intervalles]
\textbf{Énoncé :}
\begin{enumerate}
    \item Comparer $\sqrt{7}$ et $2,6$.
    \item Encadrer $\sqrt{10}$ au centième ($10^{-2}$) près.
    \item Écrire sous forme d'intervalle l'ensemble $E = \{x \in \mathbb{R} \mid -2 < x \le 5\}$.
    \item Résoudre l'inéquation $3x - 7 \le 2(x+1)$ et donner l'ensemble des solutions sous forme d'intervalle.
\end{enumerate}
\tcblower
\textbf{Solution détaillée :}
\begin{enumerate}
    \item \textbf{Comparaison :} Les deux nombres sont positifs. On compare les carrés.
    $(\sqrt{7})^2 = 7$.
    $(2,6)^2 = 6,76$.
    Comme $7 > 6,76$, on a $\sqrt{7} > 2,6$.
    \item \textbf{Encadrement de $\sqrt{10}$ au centième :}
    On cherche $k$ tel que $\frac{k}{100} \le \sqrt{10} < \frac{k+1}{100}$.
    $\sqrt{10} \approx 3,1622\dots$ (calculatrice).
    Au centième : $3,16 \le \sqrt{10} < 3,17$.
    \textbf{Vérification rigoureuse (sans calculatrice) :}
    $3,16^2 = 9,9856 < 10$ et $3,17^2 = 10,0489 > 10$.
    L'encadrement est correct : $\boxed{3,16 \le \sqrt{10} < 3,17}$.
    \item \textbf{Intervalle :} $x$ strictement plus grand que $-2$ (parenthèse) et inférieur ou égal à $5$ (crochet).
    $\boxed{E = ]-2 ; 5]}$.
    \item \textbf{Inéquation :}
    \begin{align*}
    3x - 7 &\le 2(x+1) \\
    3x - 7 &\le 2x + 2 \\
    3x - 2x &\le 2 + 7 \\
    x &\le 9
    \end{align*}
    Ensemble des solutions : $\boxed{S = ]-\infty ; 9]}$.
\end{enumerate}

\begin{popfigure}
\begin{tikzpicture}[scale=0.7, >=Stealth]
% Ligne pour ]-2; 5]
\draw[thick, ->] (-3,0) -- (6,0) node[right] {$x$};
\foreach \x in {-2, 5} \draw (\x,2pt) -- (\x,-2pt) node[below] {$\x$};
\draw[popblue, very thick] (-2,0.4) -- (5,0.4);
\fill[white, draw=popblue, thick] (-2,0.4) circle (3pt); % ouvert
\fill[popblue] (5,0.4) circle (3pt); % fermé
\node[popblue, above] at (1.5,0.7) {$E = ]-2; 5]$};

% Ligne pour ]-infty; 9]
\draw[thick, ->] (-3,-1.5) -- (10,-1.5) node[right] {$x$};
\draw (9,2pt-1.5) -- (9,-2pt-1.5) node[below] {$9$};
\draw[poporange, very thick, <-] (-3,-1.1) -- (9,-1.1);
\fill[poporange] (9,-1.1) circle (3pt);
\node[poporange, above] at (3,-0.8) {$S = ]-\infty; 9]$};
\end{tikzpicture}
\legende{Représentation des ensembles solutions $E$ et $S$ sur la droite graduée.}
\end{popfigure}
\end{exoresolu}

\courssub{Résoudre des problèmes concrets mettant en jeu les nombres réels}

\begin{methode}[Modéliser une situation concrète avec les nombres réels]
\textbf{Objectif :} Traduire un problème (géométrie, vie courante) en calculs sur $\mathbb{R}$ et répondre par une phrase.
\begin{enumerate}
    \item \textbf{Lire et comprendre} : Identifier les grandeurs, les données numériques, la question.
    \item \textbf{Modéliser} :
    \begin{itemize}
        \item Choisir une variable $x$ pour l'inconnue principale.
        \item Traduire les relations (Pythagore, Thalès, Aire, Volume, Vitesse, Pourcentage) en égalités ou inéquations.
        \item Attention aux unités (conversion m $\leftrightarrow$ cm, km/h $\leftrightarrow$ m/s, FCFA).
    \end{itemize}
    \item \textbf{Résoudre} : Appliquer les techniques du chapitre (simplification radicaux, résolution équation/inéquation 1er degré, puissances).
    \item \textbf{Interpréter} : Vérifier la cohérence (longueur positive, temps positif, nombre entier si comptage). Écrire la réponse finale dans une phrase complète avec unité.
\end{enumerate}
\end{methode}

\begin{exoresolu}[Problème concret : Électrification rurale]
\textbf{Énoncé :} Pour électrifier un village, on doit tendre un câble entre deux poteaux distants de $20$ m. Le câble doit pendre de $2$ m au milieu (flèche). On modélise le câble par une arc de cercle (approximation parabolique trop complexe, on utilise Pythagore sur le triangle rectangle demi-portée/flèche/demi-câble).
On note $R$ le rayon du cercle, $L$ la longueur du câble (approximée par la corde ? Non, ici on suppose que le câble tendu forme l'hypoténuse).
\textit{Correction du modèle pour 3ème :} On considère le triangle rectangle formé par la moitié de la distance entre poteaux ($10$ m), la flèche ($2$ m) et la moitié du câble tendu $d$.
\begin{enumerate}
    \item Calculer la longueur exacte $d$ de la demi-corde (demi-câble) en utilisant le théorème de Pythagore. Simplifier le radical.
    \item En déduire la longueur totale $L$ du câble nécessaire sous la forme $a\sqrt{b}$.
    \item Le câble est vendu par bobines de $50$ m au prix de $45\,000$ FCFA. Combien de bobines faut-il acheter ? Quel est le coût total ?
\end{enumerate}
\tcblower
\textbf{Solution détaillée :}
\begin{enumerate}
    \item \textbf{Modélisation :} Triangle rectangle : hypoténuse $d$, côtés $10$ m et $2$ m.
    D'après Pythagore : $d^2 = 10^2 + 2^2 = 100 + 4 = 104$.
    $d = \sqrt{104} = \sqrt{4 \times 26} = 2\sqrt{26}$ m. (Longueur exacte).
    \item \textbf{Longueur totale :} $L = 2d = 2 \times 2\sqrt{26} = \boxed{4\sqrt{26} \text{ m}}$.
    \item \textbf{Achat :} Valeur approchée : $\sqrt{26} \approx 5,099$. $L \approx 4 \times 5,099 = 20,396$ m.
    Il faut au moins $20,4$ m. Une bobine fait $50$ m. \textbf{1 bobine suffit}.
    Coût total = $\boxed{45\,000 \text{ FCFA}}$.
\end{enumerate}
\textbf{Conclusion :} Il faut acheter 1 bobine pour un coût de 45 000 FCFA. Le câble restant servira pour les raccordements.

\begin{popfigure}
\begin{tikzpicture}[scale=0.5, >=Stealth]
% Coordinates
\coordinate (A) at (0,0);
\coordinate (B) at (10,0);
\coordinate (C) at (10,2);
% Triangle
\draw[popblue, thick] (A) -- (B) -- (C) -- cycle;
% Right angle mark
\draw[popblue] (9.5,0) -- (9.5,0.5) -- (10,0.5);
% Labels
\node[below] at (5,0) {10 m (demi-portée)};
\node[right] at (10,1) {2 m (flèche)};
\node[above left] at (5,1) {$d$ (demi-câble)};
% Points
\fill[popdark] (A) circle (3pt) node[below left] {Poteau 1};
\fill[popdark] (B) circle (3pt) node[below right] {Milieu};
\fill[popdark] (C) circle (3pt) node[above right] {Poteau 2};
\end{tikzpicture}
\legende{Modélisation par triangle rectangle pour le câble (exercice électrification).}
\end{popfigure}
\end{exoresolu}

\begin{attention}[Les 5 erreurs qui coûtent des points au BEPC]
\begin{enumerate}
    \item \textbf{Oublier la valeur absolue ou la condition de positivité} : Écrire $\sqrt{x^2} = x$ au lieu de $|x|$, ou $\sqrt{a}\sqrt{b} = \sqrt{ab}$ sans préciser $a,b \ge 0$. \textbf{Réflexe} : Vérifier le signe des nombres sous le radical.
    \item \textbf{Additionner les radicandes} : Croire que $\sqrt{a} + \sqrt{b} = \sqrt{a+b}$. \textbf{FAUX}. Ex: $\sqrt{9}+\sqrt{16}=3+4=7 \neq \sqrt{25}=5$. On n'additionne que les \textbf{coefficients} des radicaux \textbf{semblables}.
    \item \textbf{Laisser un radical au dénominateur} : Réponse $\frac{3}{\sqrt{2}}$ au lieu de $\frac{3\sqrt{2}}{2}$. \textbf{Obligatoire} : Rationaliser le dénominateur à la dernière étape.
    \item \textbf{Erreur de signe aux puissances négatives} : Confondre $a^{-n} = \frac{1}{a^n}$ et $a^{-n} = -a^n$. Ex: $2^{-3} = \frac{1}{8}$, pas $-8$. Aussi : $-3^2 = -9$ mais $(-3)^2 = 9$. \textbf{Parenthèses !}
    \item \textbf{Encadrement mal écrit} : Écrire $3,16 < \sqrt{10} < 3,17$ au lieu de $3,16 \le \sqrt{10} < 3,17$ (le nombre \textbf{est} dans l'intervalle, la borne inférieure est incluse car $3,16^2 < 10$). Vérifier toujours avec les carrés. Pour les intervalles : $]-\infty; 5]$ (crochet ouvert à l'infini, fermé à 5).
\end{enumerate}
\end{attention}

\newpage

\coursec{Exercices}

\courssub{Je vérifie mes connaissances}

\begin{exercice}[QCM : Racines carrées et ensemble $\mathbb{R}$]
Parmi les affirmations suivantes, une seule est vraie. Laquelle ? Justifier brièvement.
\begin{enumerate}
    \item $\sqrt{(-5)^2} = -5$
    \item $\sqrt{18} = 3\sqrt{2}$
    \item $\sqrt{9+16} = \sqrt{9} + \sqrt{16}$
    \item $\dfrac{\sqrt{50}}{\sqrt{2}} = 25$
\end{enumerate}
\end{exercice}

\begin{exercice}[Vrai ou Faux ? Justifier]
Dire si les affirmations suivantes sont vraies ou fausses. En cas de fausseté, donner un contre-exemple ou corriger l'égalité.
\begin{enumerate}
    \item Pour tout réel $a \ge 0$, $\sqrt{a^2} = a$.
    \item Pour tous réels $a \ge 0$ et $b \ge 0$, $\sqrt{a+b} = \sqrt{a} + \sqrt{b}$.
    \item $\sqrt{12} - \sqrt{3} = \sqrt{9}$.
    \item $(\sqrt{7})^2 = 7$.
    \item $\dfrac{1}{\sqrt{2}} = \dfrac{\sqrt{2}}{2}$.
\end{enumerate}
\end{exercice}

\begin{exercice}[Vocabulaire et définitions]
Compléter les phrases suivantes en utilisant le vocabulaire mathématique précis.
\begin{enumerate}
    \item Si $x^2 = 20$, alors $x = \dots$ ou $x = \dots$ (on donnera la forme exacte simplifiée).
    \item Un nombre qui ne peut pas s'écrire sous la forme $\frac{p}{q}$ avec $p \in \mathbb{Z}$ et $q \in \mathbb{Z}^*$ est dit \dots
    \item La forme simplifiée de $\sqrt{72}$ s'écrit $a\sqrt{b}$ avec $a$ et $b$ entiers, $b$ le plus petit possible. Ici $a = \dots$ et $b = \dots$.
    \item Rationaliser le dénominateur de $\frac{5}{\sqrt{3}}$ consiste à multiplier le numérateur et le dénominateur par \dots
\end{enumerate}
\end{exercice}

\begin{exercice}[Calculs directs et simplification]
Effectuer les calculs suivants et donner les résultats sous la forme $a\sqrt{b}$ avec $a,b$ entiers et $b$ le plus petit possible, ou sous forme décimale exacte si possible.
\begin{enumerate}
    \item $\sqrt{48}$
    \item $\sqrt{75}$
    \item $\sqrt{200}$
    \item $\sqrt{12} \times \sqrt{27}$
    \item $(\sqrt{5})^2 \times \sqrt{5}$
    \item $2\sqrt{3} \times 3\sqrt{2}$
\end{enumerate}
\end{exercice}

\courssub{Je m'entraîne}

\begin{exercice}[Simplification et calcul approché]
Réduire l'expression $E = 5\sqrt{12} - 3\sqrt{27} + 2\sqrt{75}$ sous la forme $a\sqrt{b}$ ($a,b$ entiers, $b$ minimal).
Ensuite, donner une valeur approchée de $E$ au centième près.
\end{exercice}

\begin{exercice}[Rationalisation des dénominateurs]
Rendre rationnel le dénominateur des fractions suivantes et simplifier le résultat final.
\begin{enumerate}
    \item $A = \dfrac{7}{\sqrt{3}}$
    \item $B = \dfrac{6}{\sqrt{7} - \sqrt{2}}$
    \item $C = \dfrac{\sqrt{5} + 1}{\sqrt{5} - 1}$
    \item $D = \dfrac{4\sqrt{2}}{\sqrt{6} + \sqrt{2}}$
\end{enumerate}
\end{exercice}

\begin{exercice}[Puissances entières et radicaux]
Calculer les expressions suivantes et donner le résultat sous forme simplifiée (entier, fraction irréductible ou forme $a\sqrt{b}$).
\begin{enumerate}
    \item $(\sqrt{2})^4 \times (\sqrt{2})^{-3}$
    \item $(2\sqrt{5})^3$
    \item $(3^{-2} \times 3^5)^2$
    \item $\left(\dfrac{\sqrt{3}}{3}\right)^{-2}$
    \item $(\sqrt{6} \times \sqrt{2})^2$
\end{enumerate}
\end{exercice}

\begin{exercice}[Comparaison, encadrement et intervalles]
\begin{enumerate}
    \item Comparer les nombres suivants sans calculatrice (justifier par le calcul des carrés ou factorisation) :
        \begin{itemize}
            \item $3\sqrt{5}$ et $7$
            \item $\sqrt{11} + \sqrt{6}$ et $\sqrt{17}$
            \item $2\sqrt{3}$ et $3\sqrt{2}$
        \end{itemize}
    \item Donner un encadrement de $\sqrt{83}$ à $10^{-2}$ près, puis un arrondi au dixième en justifiant.
    \item Traduire en intervalles de $\mathbb{R}$ puis représenter sur une droite graduée :
        \begin{itemize}
            \item $I_1 = \{x \in \mathbb{R} \mid -1 \le x < 4\}$
            \item $I_2 = \{x \in \mathbb{R} \mid x > 2\}$
            \item $I_3 = \{x \in \mathbb{R} \mid x \le 0\}$
        \end{itemize}
    \item Déterminer l'intervalle $I_1 \cap I_2$ et le représenter.
\end{enumerate}
\end{exercice}

\courssub{Je me prépare au Probatoire}

\begin{exercice}[Problème de synthèse : Clôture de champs (Type BEPC)]
Un agriculteur de la région de l'Ouest dispose de deux terrains de même aire $300~\text{m}^2$.
\begin{itemize}
    \item Le premier terrain est un \textbf{carré}.
    \item Le second est un \textbf{rectangle} dont la longueur est supérieure de $5~\text{m}$ à la largeur.
\end{itemize}
\begin{enumerate}
    \item Soit $c$ la longueur du côté du carré. Exprimer $c$ en fonction de l'aire. Calculer la valeur exacte de $c$ sous la forme $a\sqrt{b}$ ($a,b$ entiers). Donner un arrondi au centimètre près.
    \item Calculer le périmètre $P_c$ du carré (valeur exacte puis arrondi au cm).
    \item Soit $x$ la largeur du rectangle (en mètres). Écrire une équation du second degré vérifiée par $x$. Résoudre cette équation et donner les dimensions exactes du rectangle (forme $a\sqrt{b} \pm c$). Donner les arrondis au centimètre.
    \item Calculer le périmètre $P_r$ du rectangle (valeur exacte simplifiée et arrondi au cm).
    \item Comparer $P_c$ et $P_r$. Quel terrain nécessite le moins de grillage pour être clôturé ?
    \item Si le grillage coûte $2500~\text{FCFA}$ le mètre linéaire, calculer la différence de coût entre les deux clôtures (arrondi au franc CFA).
\end{enumerate}
\textit{Rédaction complète, justifiée et structurée attendue.}
\end{exercice}

\begin{exercice}[Démonstration guidée : Irrationalité de $\sqrt{2}$]
On veut démontrer que $\sqrt{2}$ est irrationnel (ne s'écrit pas sous forme de fraction irréductible $\frac{p}{q}$).
Suivre le raisonnement par l'absurde ci-dessous.
\begin{enumerate}
    \item Supposons que $\sqrt{2}$ est rationnel. On peut alors écrire $\sqrt{2} = \frac{p}{q}$ avec $p,q$ entiers naturels non nuls, $q \neq 0$ et la fraction \textbf{irréductible} (PGCD$(p,q)=1$).
    \item Élever au carré les deux membres. En déduire une relation entre $p^2$ et $q^2$.
    \item Montrer que $p^2$ est pair. En déduire que $p$ est pair (rappeler la propriété : si un carré est pair, le nombre est pair).
    \item Puisque $p$ est pair, on peut écrire $p = 2k$ ($k \in \mathbb{N}^*$). Remplacer dans la relation trouvée à la question 2 et simplifier par 2.
    \item Montrer que $q^2$ est pair, donc $q$ est pair.
    \item Conclure : $p$ et $q$ sont tous les deux pairs. Cela contredit l'hypothèse de départ (fraction irréductible). Quelle est la conclusion finale sur $\sqrt{2}$ ?
\end{enumerate}
\end{exercice}

\begin{exercice}[Problème concret : Antenne relais et zone de couverture (Contexte Cameroun)]
Une entreprise de télécommunications (type MTN ou Orange) installe une antenne relais au point $O$. La zone de couverture au sol est modélisée par un disque de centre $O$ et de rayon $R = \sqrt{5000}$ mètres.
Un village se situe au point $V$ tel que $OV = 70~\text{m}$.
Une route nationale passe par les points $A$ et $B$ tels que $OA = 50~\text{m}$ et $OB = 80~\text{m}$. On suppose $O, A, B$ alignés dans cet ordre.
\begin{enumerate}
    \item Simplifier l'expression du rayon $R = \sqrt{5000}$ sous la forme $a\sqrt{b}$ ($a,b$ entiers, $b$ minimal).
    \item Calculer la valeur approchée de $R$ au mètre près.
    \item Le village $V$ est-il couvert par l'antenne ? Justifier par un calcul ou une comparaison.
    \item Déterminer si les points $A$ et $B$ sont dans la zone de couverture.
    \item La route traverse-t-elle la zone de couverture ? Justifier.
    \item On veut installer un second village $W$ sur la droite $(OB)$ de manière qu'il soit \textbf{juste à la limite} de la zone de couverture. Calculer la distance $OW$ exacte et approchée au mètre.
\end{enumerate}
\end{exercice}

\newpage

\coursec{Corrigés des exercices}

\begin{corrige}[Exercice 1 — QCM : Racines carrées et ensemble $\mathbb{R}$]
\textbf{Analyse de chaque proposition :}

\begin{enumerate}
    \item \textbf{Fausse.} $\sqrt{(-5)^2} = \sqrt{25} = 5 \neq -5$. Rappel : pour tout réel $a$, $\sqrt{a^2} = |a|$. Ici $|-5| = 5$.
    \item \textbf{Vraie.} $\sqrt{18} = \sqrt{9 \times 2} = \sqrt{9} \times \sqrt{2} = 3\sqrt{2}$. C'est la \cle{forme simplifiée} car $2$ n'a pas de diviseur \cle{carré parfait} (autre que $1$).
    \item \textbf{Fausse.} $\sqrt{9+16} = \sqrt{25} = 5$, alors que $\sqrt{9} + \sqrt{16} = 3 + 4 = 7$. La racine carrée n'est pas distributive par rapport à l'addition : $\sqrt{a+b} \neq \sqrt{a} + \sqrt{b}$ en général.
    \item \textbf{Fausse.} $\dfrac{\sqrt{50}}{\sqrt{2}} = \sqrt{\dfrac{50}{2}} = \sqrt{25} = 5 \neq 25$.
\end{enumerate}
\textbf{Conclusion :} La seule affirmation vraie est la \textbf{proposition 2}.
\end{corrige}

\begin{corrige}[Exercice 2 — Vrai ou Faux ? Justifier]
\begin{enumerate}
    \item \textbf{Vrai.} Par définition, pour tout réel $a \ge 0$, $\sqrt{a^2} = a$ (car $|a| = a$ si $a \ge 0$).
    \item \textbf{Faux.} Contre-exemple : $a=9, b=16$. $\sqrt{9+16} = \sqrt{25} = 5$ mais $\sqrt{9}+\sqrt{16}=3+4=7$. L'égalité $\sqrt{a+b}=\sqrt{a}+\sqrt{b}$ est fausse en général.
    \item \textbf{Faux.} $\sqrt{12} - \sqrt{3} = \sqrt{4\times 3} - \sqrt{3} = 2\sqrt{3} - \sqrt{3} = \sqrt{3}$. Or $\sqrt{9} = 3$. Comme $\sqrt{3} \approx 1,73 \neq 3$, l'égalité est fausse. Correction : $\sqrt{12} - \sqrt{3} = \sqrt{3}$.
    \item \textbf{Vrai.} Par définition de la racine carrée, pour tout $a \ge 0$, $(\sqrt{a})^2 = a$. Ici $a=7 \ge 0$, donc $(\sqrt{7})^2 = 7$.
    \item \textbf{Vrai.} On \cle{rationalise le dénominateur} : $\dfrac{1}{\sqrt{2}} = \dfrac{1 \times \sqrt{2}}{\sqrt{2} \times \sqrt{2}} = \dfrac{\sqrt{2}}{2}$.
\end{enumerate}
\end{corrige}

\begin{corrige}[Exercice 3 — Vocabulaire et défininitions]
\begin{enumerate}
    \item Si $x^2 = 20$, alors $x = \sqrt{20}$ ou $x = -\sqrt{20}$. \cle{Forme exacte simplifiée} : $x = 2\sqrt{5}$ ou $x = -2\sqrt{5}$.
    \item Un nombre qui ne peut pas s'écrire sous la forme $\frac{p}{q}$ avec $p \in \mathbb{Z}$ et $q \in \mathbb{Z}^*$ est dit \cle{irrationnel}.
    \item $\sqrt{72} = \sqrt{36 \times 2} = 6\sqrt{2}$. Donc $a = 6$ et $b = 2$.
    \item \cle{Rationaliser le dénominateur} de $\frac{5}{\sqrt{3}}$ consiste à multiplier le numérateur et le dénominateur par \textbf{$\sqrt{3}$}.
\end{enumerate}
\end{corrige}

\begin{corrige}[Exercice 4 — Calculs directs et simplification]
Rappel : On cherche la forme $a\sqrt{b}$ avec $a,b \in \mathbb{Z}$, $b \ge 0$ le plus petit possible (sans facteur carré).
\begin{enumerate}
    \item $\sqrt{48} = \sqrt{16 \times 3} = \sqrt{16}\times\sqrt{3} = 4\sqrt{3}$.
    \item $\sqrt{75} = \sqrt{25 \times 3} = \sqrt{25}\times\sqrt{3} = 5\sqrt{3}$.
    \item $\sqrt{200} = \sqrt{100 \times 2} = \sqrt{100}\times\sqrt{2} = 10\sqrt{2}$.
    \item $\sqrt{12} \times \sqrt{27} = \sqrt{12 \times 27} = \sqrt{324} = 18$ (car $18^2=324$). \textit{Autre méthode :} $2\sqrt{3} \times 3\sqrt{3} = 6 \times 3 = 18$.
    \item $(\sqrt{5})^2 \times \sqrt{5} = 5 \times \sqrt{5} = 5\sqrt{5}$.
    \item $2\sqrt{3} \times 3\sqrt{2} = (2\times 3) \times (\sqrt{3}\times\sqrt{2}) = 6\sqrt{6}$.
\end{enumerate}
\end{corrige}

\begin{corrige}[Exercice 5 — Simplification et calcul approché]
\textbf{Simplification de $E = 5\sqrt{12} - 3\sqrt{27} + 2\sqrt{75}$ :}
On simplifie chaque radical :
$\sqrt{12} = \sqrt{4\times 3} = 2\sqrt{3}$
$\sqrt{27} = \sqrt{9\times 3} = 3\sqrt{3}$
$\sqrt{75} = \sqrt{25\times 3} = 5\sqrt{3}$

On substitue :
$E = 5(2\sqrt{3}) - 3(3\sqrt{3}) + 2(5\sqrt{3})$
$E = 10\sqrt{3} - 9\sqrt{3} + 10\sqrt{3}$
$E = (10 - 9 + 10)\sqrt{3}$
$E = 11\sqrt{3}$

\textbf{Valeur approchée au centième près :}
$\sqrt{3} \approx 1,73205...$
$E \approx 11 \times 1,73205 = 19,05255...$
Arrondi au centième : $\boxed{E \approx 19,05}$.
\end{corrige}

\begin{corrige}[Exercice 6 — Rationalisation des dénominateurs]
\begin{enumerate}
    \item $A = \dfrac{7}{\sqrt{3}} = \dfrac{7 \times \sqrt{3}}{\sqrt{3} \times \sqrt{3}} = \dfrac{7\sqrt{3}}{3}$.
    \item $B = \dfrac{6}{\sqrt{7} - \sqrt{2}}$. On multiplie par le \cle{conjugué} $\sqrt{7} + \sqrt{2}$ :
    $B = \dfrac{6(\sqrt{7} + \sqrt{2})}{(\sqrt{7} - \sqrt{2})(\sqrt{7} + \sqrt{2})} = \dfrac{6(\sqrt{7} + \sqrt{2})}{7 - 2} = \dfrac{6(\sqrt{7} + \sqrt{2})}{5}$.
    \item $C = \dfrac{\sqrt{5} + 1}{\sqrt{5} - 1}$. On multiplie par le conjugué $\sqrt{5} + 1$ :
    $C = \dfrac{(\sqrt{5} + 1)^2}{(\sqrt{5})^2 - 1^2} = \dfrac{5 + 2\sqrt{5} + 1}{5 - 1} = \dfrac{6 + 2\sqrt{5}}{4} = \dfrac{3 + \sqrt{5}}{2}$.
    \item $D = \dfrac{4\sqrt{2}}{\sqrt{6} + \sqrt{2}}$. On multiplie par le conjugué $\sqrt{6} - \sqrt{2}$ :
    $D = \dfrac{4\sqrt{2}(\sqrt{6} - \sqrt{2})}{(\sqrt{6})^2 - (\sqrt{2})^2} = \dfrac{4\sqrt{12} - 4\sqrt{4}}{6 - 2} = \dfrac{4(2\sqrt{3}) - 4(2)}{4} = \dfrac{8\sqrt{3} - 8}{4} = 2\sqrt{3} - 2$.
\end{enumerate}
\end{corrige}

\begin{corrige}[Exercice 7 — Puissances entières et radicaux]
Rappels : $a^n \times a^m = a^{n+m}$, $(a^m)^n = a^{mn}$, $a^{-n} = \frac{1}{a^n}$, $(\sqrt{a})^2 = a$.
\begin{enumerate}
    \item $(\sqrt{2})^4 \times (\sqrt{2})^{-3} = (\sqrt{2})^{4-3} = (\sqrt{2})^1 = \sqrt{2}$.
    \item $(2\sqrt{5})^3 = 2^3 \times (\sqrt{5})^3 = 8 \times (\sqrt{5})^2 \times \sqrt{5} = 8 \times 5 \times \sqrt{5} = 40\sqrt{5}$.
    \item $(3^{-2} \times 3^5)^2 = (3^{-2+5})^2 = (3^3)^2 = 3^{6} = 729$.
    \item $\left(\dfrac{\sqrt{3}}{3}\right)^{-2} = \left(\dfrac{3}{\sqrt{3}}\right)^{2} = \left(\sqrt{3}\right)^2 = 3$. \textit{Autre méthode :} $\left(\dfrac{\sqrt{3}}{3}\right)^{-2} = \dfrac{(\sqrt{3})^{-2}}{3^{-2}} = \dfrac{\frac{1}{3}}{\frac{1}{9}} = \frac{1}{3} \times 9 = 3$.
    \item $(\sqrt{6} \times \sqrt{2})^2 = (\sqrt{12})^2 = 12$. \textit{Autre méthode :} $(\sqrt{6})^2 \times (\sqrt{2})^2 = 6 \times 2 = 12$.
\end{enumerate}
\end{corrige}

\begin{corrige}[Exercice 8 — Comparaison, encadrement et intervalles]
\begin{enumerate}
    \item \textbf{Comparaisons (méthode : mise au carré car nombres positifs) :}
    \begin{itemize}
        \item $3\sqrt{5}$ et $7$. $(3\sqrt{5})^2 = 9 \times 5 = 45$. $7^2 = 49$. $45 < 49 \Rightarrow 3\sqrt{5} < 7$.
        \item $\sqrt{11} + \sqrt{6}$ et $\sqrt{17}$. $(\sqrt{11} + \sqrt{6})^2 = 11 + 6 + 2\sqrt{66} = 17 + 2\sqrt{66}$. $(\sqrt{17})^2 = 17$. Comme $2\sqrt{66} > 0$, le premier carré est plus grand. Donc $\sqrt{11} + \sqrt{6} > \sqrt{17}$.
        \item $2\sqrt{3}$ et $3\sqrt{2}$. $(2\sqrt{3})^2 = 4 \times 3 = 12$. $(3\sqrt{2})^2 = 9 \times 2 = 18$. $12 < 18 \Rightarrow 2\sqrt{3} < 3\sqrt{2}$.
    \end{itemize}
    \item \textbf{Encadrement de $\sqrt{83}$ à $10^{-2}$ près :}
    On cherche $a, b$ entiers tels que $a^2 < 83 < b^2$. $9^2=81$, $10^2=100$. Donc $9 < \sqrt{83} < 10$.
    On affine au dixième : $9,1^2 = 82,81$ ; $9,2^2 = 84,64$. Donc $9,1 < \sqrt{83} < 9,2$.
    On affine au centième : $9,11^2 = 82,9921$ ; $9,12^2 = 83,1744$.
    Donc $9,11 < \sqrt{83} < 9,12$. C'est l'\cle{encadrement} à $10^{-2}$ près.
    \textbf{Arrondi au dixième :} Comme $\sqrt{83} \approx 9,11...$, le chiffre des centièmes est $1 < 5$. Arrondi au dixième : $\boxed{9,1}$.
    \item \textbf{Traduction en intervalles et représentation :}
    \begin{itemize}
        \item $I_1 = \{x \in \mathbb{R} \mid -1 \le x < 4\} = \boxed{[-1 ; 4[}$.
        \item $I_2 = \{x \in \mathbb{R} \mid x > 2\} = \boxed{]2 ; +\infty[}$.
        \item $I_3 = \{x \in \mathbb{R} \mid x \le 0\} = \boxed{]-\infty ; 0]}$.
    \end{itemize}
    \textit{Représentation sur droite graduée :}
    \begin{popfigure}
    \begin{tikzpicture}[scale=0.6]
        \draw[->] (-3,0) -- (6,0) node[right] {$x$};
        \foreach \x in {-2,-1,0,1,2,3,4,5} \draw (\x,0.1) -- (\x,-0.1) node[below] {\x};
        % I1 = [-1; 4[
        \draw[popblue, very thick] (-1,0.3) -- (4,0.3);
        \fill[popblue] (-1,0.3) circle (3pt);
        \draw[popblue] (4,0.3) circle (3pt);
        \node[above, text=popblue] at (1.5,0.5) {$I_1 = [-1; 4[$};
        % I2 = ]2; +inf[
        \draw[poporange, very thick] (2,-0.3) -- (5.5,-0.3);
        \draw[poporange] (2,-0.3) circle (3pt);
        \draw[poporange, ->] (5.5,-0.3) -- (5.8,-0.3);
        \node[below, text=poporange] at (3.5,-0.7) {$I_2 = ]2; +\infty[$};
        % I3 = ]-inf; 0]
        \draw[popgreen, very thick] (-2.5,-0.9) -- (0,-0.9);
        \draw[popgreen, <-] (-2.5,-0.9) -- (-2.8,-0.9);
        \fill[popgreen] (0,-0.9) circle (3pt);
        \node[below, text=popgreen] at (-1.5,-1.3) {$I_3 = ]-\infty; 0]$};
    \end{tikzpicture}
    \legende{Représentation graphique des intervalles $I_1$, $I_2$ et $I_3$.}
    \end{popfigure}
    \item \textbf{Intersection $I_1 \cap I_2$ :}
    $I_1 = [-1 ; 4[$, $I_2 = ]2 ; +\infty[$.
    Les nombres appartenant aux deux intervalles sont ceux $> 2$ et $< 4$ (et $\ge -1$ qui est implicite).
    $I_1 \cap I_2 = \boxed{]2 ; 4[}$.
    \textit{Représentation :} Sur la droite, portion commune entre $2$ (exclu) et $4$ (exclu).
\end{enumerate}
\end{corrige}

\begin{corrige}[Exercice 9 — Problème de synthèse : Clôture de champs (Type BEPC)]
\textbf{Barème indicatif (sur 20 pts) : Q1 (4 pts), Q2 (2 pts), Q3 (5 pts), Q4 (3 pts), Q5 (3 pts), Q6 (3 pts).}
\textbf{Points de vigilance :} Unités (\SI{}{m}, \SI{}{m^2}), forme exacte $a\sqrt{b}$, arrondis demandés, rédaction structurée, conclusion finale.

\begin{enumerate}
    \item \textbf{Carré :} Aire $A = c^2 = 300$.
    $c = \sqrt{300} = \sqrt{100 \times 3} = 10\sqrt{3}$ m. (Forme exacte : $a=10, b=3$).
    Valeur approchée : $c \approx 10 \times 1,73205 = 17,3205$ m.
    Arrondi au centimètre près (2 décimales) : $\boxed{c \approx \SI{17,32}{m}}$.
    
    \item \textbf{Périmètre du carré :} $P_c = 4c = 4 \times 10\sqrt{3} = 40\sqrt{3}$ m.
    Valeur approchée : $P_c \approx 40 \times 1,73205 = 69,282$ m.
    Arrondi au cm : $\boxed{P_c \approx \SI{69,28}{m}}$.
    
    \item \textbf{Rectangle :} Largeur $x$, Longueur $x+5$. Aire $x(x+5) = 300$.
    Équation : $x^2 + 5x - 300 = 0$.
    Discriminant $\Delta = 5^2 - 4(1)(-300) = 25 + 1200 = 1225 = 35^2$.
    Racines : $x_1 = \frac{-5 - 35}{2} = -20$ (rejetée, longueur négative).
    $x_2 = \frac{-5 + 35}{2} = 15$.
    Largeur $x = 15$ m. Longueur $L = 20$ m.
    \textit{Note :} Les dimensions sont des entiers ici, la forme $a\sqrt{b} \pm c$ se simplifie en entiers. Dimensions exactes : $\boxed{\SI{15}{m} \text{ et } \SI{20}{m}}$. Arrondis au cm : $\SI{15,00}{m}$ et $\SI{20,00}{m}$.
    
    \item \textbf{Périmètre du rectangle :} $P_r = 2(x + L) = 2(15 + 20) = 70$ m.
    Valeur exacte : $\boxed{\SI{70}{m}}$. Arrondi : $\SI{70,00}{m}$.
    
    \item \textbf{Comparaison :} $P_c = 40\sqrt{3} \approx \SI{69,28}{m}$. $P_r = \SI{70}{m}$.
    $40\sqrt{3} < 70$ car $(40\sqrt{3})^2 = 4800 < 4900 = 70^2$.
    Le carré a un périmètre plus petit. \textbf{Le terrain carré nécessite le moins de grillage.}
    
    \item \textbf{Coût :} Prix = \SI{2500}{FCFA/m}.
    Coût carré $C_c = P_c \times 2500 = 40\sqrt{3} \times 2500 = 100\,000\sqrt{3}$ FCFA.
    Coût rectangle $C_r = 70 \times 2500 = 175\,000$ FCFA.
    Différence $\Delta C = C_r - C_c = 175\,000 - 100\,000\sqrt{3}$.
    Valeur approchée : $100\,000 \times 1,73205 = 173\,205$ FCFA.
    $\Delta C \approx 175\,000 - 173\,205 = 1\,795$ FCFA.
    Arrondi au franc CFA : $\boxed{\SI{1795}{FCFA}}$.
    \textit{Le rectangle coûte 1 795 FCFA de plus à clôturer.}
\end{enumerate}
\end{corrige}

\begin{corrige}[Exercice 10 — Démonstration guidée : Irrationalité de $\sqrt{2}$]
Voici la rédaction complète de la démonstration par l'absurde.

\begin{enumerate}
    \item \textbf{Hypothèse :} Supposons que $\sqrt{2}$ est rationnel.
    Alors il existe deux entiers naturels non nuls $p$ et $q$ ($q \neq 0$) tels que $\sqrt{2} = \frac{p}{q}$, avec la fraction \textbf{irréductible} (i.e. $\text{PGCD}(p,q)=1$).
    
    \item \textbf{Élévation au carré :} $(\sqrt{2})^2 = \left(\frac{p}{q}\right)^2 \iff 2 = \frac{p^2}{q^2}$.
    En multipliant par $q^2$ : $\boxed{p^2 = 2q^2}$. \hfill (Relation 1)
    
    \item \textbf{Parité de $p$ :} D'après (1), $p^2 = 2q^2$, donc $p^2$ est un multiple de $2$, c'est-à-dire que $p^2$ est \textbf{pair}.
    Propriété : Si le carré d'un entier est pair, alors cet entier est pair.
    (Preuve rapide : contraposée : si $p$ est impair, $p=2k+1 \Rightarrow p^2=4k^2+4k+1=2(2k^2+2k)+1$ impair).
    Donc $p$ est \textbf{pair}.
    
    \item \textbf{Substitution :} Puisque $p$ est pair, on peut écrire $p = 2k$ avec $k \in \mathbb{N}^*$.
    On remplace dans (1) : $(2k)^2 = 2q^2 \iff 4k^2 = 2q^2$.
    On simplifie par $2$ : $\boxed{2k^2 = q^2}$. \hfill (Relation 2)
    
    \item \textbf{Parité de $q$ :} D'après (2), $q^2 = 2k^2$, donc $q^2$ est un multiple de $2$, c'est-à-dire que $q^2$ est \textbf{pair}.
    Par la même propriété, $q$ est \textbf{pair}.
    
    \item \textbf{Conclusion :} Nous avons montré que $p$ est pair et $q$ est pair.
    Cela signifie que $p$ et $q$ ont un diviseur commun : $2$.
    Or, au départ, nous avions supposé que la fraction $\frac{p}{q}$ était \textbf{irréductible} (PGCD$(p,q)=1$).
    C'est une \textbf{contradiction}.
    L'hypothèse de départ est donc fausse.
    \textbf{$\sqrt{2}$ n'est pas rationnel, c'est-à-dire que $\sqrt{2}$ est irrationnel.} $\blacksquare$
\end{enumerate}
\end{corrige}

\begin{corrige}[Exercice 11 — Problème concret : Antenne relais et zone de couverture (Contexte Cameroun)]
\begin{enumerate}
    \item \textbf{Simplification du rayon $R$ :}
    $R = \sqrt{5000} = \sqrt{100 \times 50} = 10\sqrt{50} = 10\sqrt{25 \times 2} = 10 \times 5 \sqrt{2} = \boxed{50\sqrt{2} \text{ m}}$.
    ($a=50, b=2$).
    
    \item \textbf{Valeur approchée au mètre près :}
    $\sqrt{2} \approx 1,4142$.
    $R \approx 50 \times 1,4142 = 70,71$ m.
    Arrondi au mètre près : $\boxed{R \approx \SI{71}{m}}$.
    
    \item \textbf{Couverture du village $V$ :}
    $OV = \SI{70}{m}$. Rayon $R \approx \SI{70,71}{m}$ (exact : $50\sqrt{2}$).
    On compare $OV$ et $R$ : $70 < 50\sqrt{2}$ car $70^2 = 4900$ et $(50\sqrt{2})^2 = 5000$.
    Comme $OV < R$, le village $V$ est \textbf{à l'intérieur du disque}, donc \textbf{couvert par l'antenne}.
    
    \item \textbf{Points $A$ et $B$ :}
    $OA = \SI{50}{m}$. $50 < 50\sqrt{2}$ (car $\sqrt{2}>1$). Donc $A$ est \textbf{couvert}.
    $OB = \SI{80}{m}$. $80 > 50\sqrt{2} \approx 70,71$. Donc $B$ n'est \textbf{pas couvert} (il est hors zone).
    
    \item \textbf{Traversée de la route :}
    Les points $O, A, B$ sont alignés dans cet ordre. $OA = \SI{50}{m}$ (dans la zone), $OB = \SI{80}{m}$ (hors zone).
    La zone de couverture est le disque de centre $O$ et rayon $R \approx \SI{70,71}{m}$.
    Sur la demi-droite $[OB)$, les points de distance à $O$ inférieure à $R$ sont couverts.
    Comme $OA < R < OB$, par continuité (théorème des valeurs intermédiaires appliqué à la fonction distance), il existe un point $M$ sur le segment $[AB]$ tel que $OM = R$.
    Donc la route \textbf{traverse la zone de couverture} (elle y entre en $A$ et en sort entre $A$ et $B$).
    
    \item \textbf{Second village $W$ à la limite :}
    $W$ est sur la droite $(OB)$, du même côté que $B$ (puisque $B$ est hors zone), tel que $OW = R$.
    Distance exacte : $\boxed{OW = 50\sqrt{2} \text{ m}}$.
    Distance approchée au mètre : $\boxed{OW \approx \SI{71}{m}}$.
\end{enumerate}
\end{corrige}

\newpage

\coursec{Fiche bilan}

\begin{popfigure}
\centering
\begin{center}\tcbox[colback=popblue,colframe=popblue,coltext=white,arc=9pt,boxsep=4pt,left=6pt,right=6pt]{\bfseries\small Nombres Réels ($\mathbb{R}$)}\end{center}
\vspace{-2pt}
\begin{tcbraster}[raster columns=3,raster equal height=rows,raster column skip=6pt,raster row skip=6pt,enhanced,blankest]
\begin{tcolorbox}[enhanced,colback=white,colframe=poporange,boxrule=0.9pt,arc=3pt,title={Racines carrées},fonttitle=\bfseries\scriptsize,coltitle=white,colbacktitle=poporange,left=4pt,right=4pt,top=2pt,bottom=2pt,boxsep=1pt,toptitle=1.5pt,bottomtitle=1.5pt,halign=left]
\begin{itemize}[nosep,leftmargin=*,itemsep=1pt,font=\scriptsize]
\item Définition : $\sqrt{a}$
\item Propriétés : $\sqrt{ab}$, $\frac{\sqrt{a}}{\sqrt{b}}$
\end{itemize}
\end{tcolorbox}
\begin{tcolorbox}[enhanced,colback=white,colframe=popgreen,boxrule=0.9pt,arc=3pt,title={Ensemble $\mathbb{R}$},fonttitle=\bfseries\scriptsize,coltitle=white,colbacktitle=popgreen,left=4pt,right=4pt,top=2pt,bottom=2pt,boxsep=1pt,toptitle=1.5pt,bottomtitle=1.5pt,halign=left]
\begin{itemize}[nosep,leftmargin=*,itemsep=1pt,font=\scriptsize]
\item Rationnels $\mathbb{Q}$
\item Irrationnels
\item Droite graduée
\end{itemize}
\end{tcolorbox}
\begin{tcolorbox}[enhanced,colback=white,colframe=poppurple,boxrule=0.9pt,arc=3pt,title={Opérations sur radicaux},fonttitle=\bfseries\scriptsize,coltitle=white,colbacktitle=poppurple,left=4pt,right=4pt,top=2pt,bottom=2pt,boxsep=1pt,toptitle=1.5pt,bottomtitle=1.5pt,halign=left]
\begin{itemize}[nosep,leftmargin=*,itemsep=1pt,font=\scriptsize]
\item Somme / Produit
\item Quotient / Rationalisation
\item Écriture $a\sqrt{b}$
\end{itemize}
\end{tcolorbox}
\begin{tcolorbox}[enhanced,colback=white,colframe=poppink,boxrule=0.9pt,arc=3pt,title={Puissances entières},fonttitle=\bfseries\scriptsize,coltitle=white,colbacktitle=poppink,left=4pt,right=4pt,top=2pt,bottom=2pt,boxsep=1pt,toptitle=1.5pt,bottomtitle=1.5pt,halign=left]
\begin{itemize}[nosep,leftmargin=*,itemsep=1pt,font=\scriptsize]
\item Règles : $a^n a^m$, $(a^n)^m$
\item Exposants négatifs
\item Notation scientifique
\end{itemize}
\end{tcolorbox}
\begin{tcolorbox}[enhanced,colback=white,colframe=popgold,boxrule=0.9pt,arc=3pt,title={Comparaison \& Encadrement},fonttitle=\bfseries\scriptsize,coltitle=white,colbacktitle=popgold,left=4pt,right=4pt,top=2pt,bottom=2pt,boxsep=1pt,toptitle=1.5pt,bottomtitle=1.5pt,halign=left]
\begin{itemize}[nosep,leftmargin=*,itemsep=1pt,font=\scriptsize]
\item Ordre sur $\mathbb{R}$
\item Encadrement décimal
\item Valeur approchée
\end{itemize}
\end{tcolorbox}
\begin{tcolorbox}[enhanced,colback=white,colframe=popblue!50!popgreen,boxrule=0.9pt,arc=3pt,title={Intervalles},fonttitle=\bfseries\scriptsize,coltitle=white,colbacktitle=popblue!50!popgreen,left=4pt,right=4pt,top=2pt,bottom=2pt,boxsep=1pt,toptitle=1.5pt,bottomtitle=1.5pt,halign=left]
\begin{itemize}[nosep,leftmargin=*,itemsep=1pt,font=\scriptsize]
\item Ouvert / Fermé
\item Borné / Non borné
\item Appartenance
\end{itemize}
\end{tcolorbox}
\end{tcbraster}

\legende{Carte mentale de synthèse : les nombres réels en classe de 3ème}
\end{popfigure}

\coursec{Bilan de la leçon : Nombres réels}

\courssub{Définitions et notations essentielles}
\begin{itemize}
    \item \cle{Racine carrée} : pour tout réel $a \ge 0$, $\sqrt{a}$ est l'\emph{unique} réel \emph{positif ou nul} dont le carré est égal à $a$ : $(\sqrt{a})^2 = a$ et $\sqrt{a} \ge 0$. Un nombre strictement négatif n'a pas de racine carrée dans $\mathbb{R}$.
    \item \cle{Nombre rationnel} : nombre qui peut s'écrire sous la forme $\frac{p}{q}$ avec $p \in \mathbb{Z}$ et $q \in \mathbb{N}^*$. Son écriture décimale est finie ou infinie \emph{périodique}. L'ensemble des rationnels est noté $\mathbb{Q}$.
    \item \cle{Nombre irrationnel} : réel qui n'est pas rationnel ; son écriture décimale est infinie et non périodique. Exemples : $\sqrt{2}$, $\sqrt{3}$, $\pi$.
    \item \cle{Nombre réel} : tout nombre rationnel ou irrationnel. L'ensemble des réels est noté $\mathbb{R}$ ; on a les inclusions $\mathbb{N} \subset \mathbb{Z} \subset \mathbb{D} \subset \mathbb{Q} \subset \mathbb{R}$. À chaque point de la droite graduée correspond un unique réel, et réciproquement.
    \item \cle{Radical simplifié} : écriture $a\sqrt{b}$ avec $a$ entier (ou rationnel) et $b \in \mathbb{N}^*$ \emph{sans facteur carré} (c'est-à-dire que $b$ n'est divisible par aucun carré parfait autre que $1$).
    \item \cle{Notation scientifique} : tout nombre décimal non nul s'écrit de manière unique $x = a \times 10^n$ avec $1 \le a < 10$ et $n \in \mathbb{Z}$.
    \item \cle{Intervalles} : $[a;b]$ (fermé), $]a;b[$ (ouvert), $[a;b[$ et $]a;b]$ (semi-ouverts), $]-\infty; a]$, $]a; +\infty[$, etc. On a l'équivalence fondamentale : $x \in [a;b] \iff a \le x \le b$.
\end{itemize}

\courssub{Formules et propriétés à connaître par cœur}
\begin{center}
\begin{tabularx}{\linewidth}{|l|X|}\hline
\rowcolor{popblueL} \textbf{Domaine} & \textbf{Règles (conditions : $a \ge 0$ et $b \ge 0$ pour les radicaux)} \\ \hline
Radicaux : produit / quotient & $\sqrt{a \times b} = \sqrt{a}\times\sqrt{b}$ \quad ; \quad $\sqrt{\dfrac{a}{b}} = \dfrac{\sqrt{a}}{\sqrt{b}}$ \ (avec $b > 0$) \\ \hline
Radicaux : carré & $(\sqrt{a})^2 = a$ \ (pour $a \ge 0$) \quad ; \quad $\sqrt{a^2} = |a|$ \ (pour tout réel $a$) \\ \hline
Simplification & $\sqrt{k^2 \times a} = k\sqrt{a}$ \ (avec $k \ge 0$) \\ \hline
Rationalisation & $\dfrac{1}{\sqrt{a}} = \dfrac{\sqrt{a}}{a}$ \ ($a > 0$) \quad ; \quad $\dfrac{1}{\sqrt{a}+\sqrt{b}} = \dfrac{\sqrt{a}-\sqrt{b}}{a-b}$ \ ($a \neq b$) \\ \hline
Puissances entières ($a \neq 0$, $b \neq 0$) & $a^n \times a^m = a^{n+m}$ \ ; \ $\dfrac{a^n}{a^m} = a^{n-m}$ \ ; \ $(a^n)^m = a^{n \times m}$ \\
& $(a \times b)^n = a^n \times b^n$ \ ; \ $\left(\dfrac{a}{b}\right)^n = \dfrac{a^n}{b^n}$ \ ; \ $a^{-n} = \dfrac{1}{a^n}$ \ ; \ $a^0 = 1$ \\ \hline
Ordre et comparaison & $a < b \iff b - a > 0$ \quad ; \quad si $0 \le a < b$ alors $\sqrt{a} < \sqrt{b}$ et $a^2 < b^2$ \\ \hline
Encadrement & Encadrer $\sqrt{a}$ entre deux entiers : trouver $n \in \mathbb{N}$ tel que $n^2 \le a < (n+1)^2$, alors $n \le \sqrt{a} < n+1$ \\ \hline
Intervalles & $x \in [a;b] \iff a \le x \le b$ \quad ; \quad $x \in ]a;b[ \iff a < x < b$ \\
& $x \in [a;+\infty[ \iff x \ge a$ \quad ; \quad $x \in ]-\infty; a[ \iff x < a$ \\ \hline
\end{tabularx}
\end{center}

\begin{attention}[Conditions d'application à ne jamais oublier]
\begin{itemize}
    \item Les formules $\sqrt{a \times b} = \sqrt{a}\times\sqrt{b}$ et $\sqrt{\dfrac{a}{b}} = \dfrac{\sqrt{a}}{\sqrt{b}}$ ne sont valables que pour des nombres \textbf{positifs ou nuls} (et $b > 0$ pour le quotient). Écrire $\sqrt{-4} \times \sqrt{-9} = \sqrt{36} = 6$ est \textbf{faux} : $\sqrt{-4}$ n'existe pas dans $\mathbb{R}$.
    \item $\sqrt{a^2} = |a|$ et non pas $a$ : par exemple $\sqrt{(-3)^2} = \sqrt{9} = 3 = |-3|$.
    \item Il n'existe \textbf{aucune formule} pour $\sqrt{a+b}$ : en général $\sqrt{a+b} \neq \sqrt{a} + \sqrt{b}$. Contre-exemple : $\sqrt{9+16} = \sqrt{25} = 5$ alors que $\sqrt{9} + \sqrt{16} = 3 + 4 = 7$.
    \item Pour rationaliser $\dfrac{1}{\sqrt{a}-\sqrt{b}}$, on multiplie par le conjugué $\sqrt{a}+\sqrt{b}$, et le dénominateur devient $a - b$ (identité remarquable $(x-y)(x+y) = x^2 - y^2$).
\end{itemize}
\end{attention}

\courssub{Méthodes express}
\begin{methode}[Les six gestes du chapitre]
\begin{enumerate}
    \item \textbf{Simplifier $\sqrt{n}$ ($n \in \mathbb{N}^*$) :} décomposer $n$ en produit de facteurs premiers, regrouper les facteurs par paires, sortir chaque paire de la racine. Exemple : $\sqrt{180} = \sqrt{2^2 \times 3^2 \times 5} = 2 \times 3 \times \sqrt{5} = 6\sqrt{5}$.
    \item \textbf{Réduire une somme de radicaux :} simplifier chaque terme sous la forme $a\sqrt{b}$, puis factoriser par le radical commun. Exemple : $3\sqrt{2} + 2\sqrt{8} = 3\sqrt{2} + 2 \times 2\sqrt{2} = 3\sqrt{2} + 4\sqrt{2} = 7\sqrt{2}$.
    \item \textbf{Rationaliser un dénominateur :}
    \begin{itemize}
        \item dénominateur monôme $\sqrt{a}$ : multiplier numérateur et dénominateur par $\sqrt{a}$ ;
        \item dénominateur binôme $\sqrt{a} \pm \sqrt{b}$ : multiplier par l'expression conjuguée $\sqrt{a} \mp \sqrt{b}$.
    \end{itemize}
    \item \textbf{Comparer deux réels $A$ et $B$ :} étudier le signe de la différence $A - B$ (mise au même dénominateur, factorisation, expression conjuguée pour les radicaux). On peut aussi comparer leurs carrés lorsque $A$ et $B$ sont positifs.
    \item \textbf{Encadrer $\sqrt{a}$ :} trouver d'abord l'entier $n$ tel que $n^2 \le a < (n+1)^2$, puis affiner au dixième, au centième, à l'aide de la calculatrice ou par essais successifs.
    \item \textbf{Écrire en notation scientifique :} placer la virgule juste après le premier chiffre non nul, compter le nombre de rangs déplacés : l'exposant est positif pour un grand nombre, négatif pour un nombre inférieur à $1$. Exemple : $\num{0,00045} = 4,5 \times 10^{-4}$.
\end{enumerate}
\end{methode}

\begin{exemplebox}[Deux exemples corrigés de rédaction]
\textbf{Exemple 1.} Écrire $A = \dfrac{3}{\sqrt{5} - \sqrt{2}}$ sans radical au dénominateur.

On multiplie numérateur et dénominateur par l'expression conjuguée $\sqrt{5} + \sqrt{2}$ :
\[ A = \dfrac{3(\sqrt{5} + \sqrt{2})}{(\sqrt{5} - \sqrt{2})(\sqrt{5} + \sqrt{2})} = \dfrac{3(\sqrt{5} + \sqrt{2})}{5 - 2} = \dfrac{3(\sqrt{5} + \sqrt{2})}{3} = \sqrt{5} + \sqrt{2}. \]

\textbf{Exemple 2.} Comparer $2\sqrt{3}$ et $3\sqrt{2}$.

Les deux nombres sont positifs ; on compare leurs carrés :
$(2\sqrt{3})^2 = 4 \times 3 = 12$ et $(3\sqrt{2})^2 = 9 \times 2 = 18$.
Comme $12 < 18$ et que les deux nombres sont positifs, on conclut : $2\sqrt{3} < 3\sqrt{2}$.
\end{exemplebox}

\begin{savaistu}[Pourquoi $\sqrt{2}$ n'est-il pas rationnel ?]
Les mathématiciens grecs de l'école de Pythagore ont découvert, en étudiant la diagonale d'un carré de côté $1$, que $\sqrt{2}$ ne peut pas s'écrire comme un quotient de deux entiers : on dit que $\sqrt{2}$ est \emph{irrationnel}. Cette découverte a bouleversé l'histoire des mathématiques ! La démonstration classique raisonne par l'absurde : si l'on suppose $\sqrt{2} = \frac{p}{q}$ avec $p$ et $q$ entiers sans facteur commun, on montre que $p$ et $q$ seraient tous les deux pairs, ce qui est contradictoire. Aujourd'hui, les nombres irrationnels sont partout : en géométrie ($\pi$ pour le cercle), en physique et dans les calculs d'ingénieurs, y compris ceux qui dimensionnent les réseaux de téléphonie mobile au Cameroun.
\end{savaistu}

\begin{aretenir}[Checklist avant le BEPC]
\begin{itemize}
    \item $\square$ Je sais donner la définition de $\sqrt{a}$ et expliquer pourquoi $\sqrt{a}$ existe seulement pour $a \ge 0$ et pourquoi $\sqrt{a} \ge 0$.
    \item $\square$ Je distingue $\sqrt{a^2} = |a|$ (valable pour tout réel $a$) et $(\sqrt{a})^2 = a$ (valable seulement pour $a \ge 0$).
    \item $\square$ J'applique correctement $\sqrt{a \times b} = \sqrt{a}\times\sqrt{b}$ et $\sqrt{a/b} = \sqrt{a}/\sqrt{b}$ \textbf{uniquement si $a \ge 0$ et $b > 0$}.
    \item $\square$ Je sais simplifier un radical ($\sqrt{75} = 5\sqrt{3}$) et réduire une somme de radicaux ($3\sqrt{2} + 2\sqrt{8} = 7\sqrt{2}$).
    \item $\square$ Je sais rationaliser un dénominateur (monôme et binôme) et écrire le résultat sous la forme $a + b\sqrt{c}$.
    \item $\square$ Je maîtrise les règles de calcul sur les puissances entières, y compris les exposants négatifs ($a^{-n} = \frac{1}{a^n}$) et $a^0 = 1$ pour $a \neq 0$.
    \item $\square$ Je sais passer de l'écriture décimale à la notation scientifique et inversement.
    \item $\square$ Je compare deux réels en étudiant le signe de leur différence ou en comparant leurs carrés (s'ils sont positifs) ; je ne me contente pas de valeurs approchées sauf demande explicite.
    \item $\square$ J'encadre un nombre réel (entier, décimal, racine carrée) à l'unité, au dixième, au centième près.
    \item $\square$ Je lis, j'écris et j'interprète les intervalles ($[a;b[$, $]-\infty; a]$, etc.) et je teste l'appartenance d'un nombre à un intervalle.
    \item $\square$ Je résous des problèmes concrets (géométrie, vie courante) en modélisant par des calculs avec radicaux ou puissances.
    \item $\square$ Je rédige mes démonstrations et justifications avec des phrases complètes, en citant les propriétés utilisées.
\end{itemize}
\end{aretenir}

\findecours

\end{document}
